% Coniques — Mathématiques Tle E — Cours signé OnBuch+
% Programme officiel MINESEC. Compiler avec : tectonic M15-coniques.tex
\newcommand{\DOCMATIERE}{Mathématiques}
\newcommand{\DOCNIVEAU}{Tle E}
\newcommand{\DOCMODULE}{Configurations et transformations élémentaires du plan}
\newcommand{\DOCLECON}{15}
\newcommand{\DOCTITRE}{Coniques}
% OnBuch+ — préambule « cours signés » ENRICHI (compilé avec tectonic / XeLaTeX)
% Système visuel « Pop » d'OnBuch+ : fond crème (jamais blanc), bordures encre
% épaisses, ombres « dures » décalées, titres Archivo Black en capitales.
% Version MATHÉMATIQUES : graphes (pgfplots/tikz), boîtes pédagogiques
% (définition, propriété, méthode, attention, le savais-tu, exercice résolu,
% exercice, TP, bilan), page de garde et signature OnBuch+ ancrée en bas.
%
% Macros attendues dans main.tex AVANT \input{preamble} :
%   \DOCMATIERE  \DOCNIVEAU  \DOCMODULE  \DOCLECON (numéro)  \DOCTITRE
\documentclass[11pt]{article}

\usepackage{fontspec}
\usepackage[french]{babel}
\usepackage[a4paper,margin=2cm,top=3.4cm,bottom=2.8cm,headheight=1.4cm,footskip=1.3cm]{geometry}
\usepackage{amsmath,amssymb,amsfonts}
\usepackage{mathtools} % \xmapsto, \coloneqq... souvent utilisés par les modèles
\usepackage{cancel} % simplifications barrées
% \abs{z} (module, valeur absolue) et \norm{u} (norme), utilisés par les modèles
\providecommand{\abs}{}\let\abs\relax\DeclarePairedDelimiter{\abs}{\lvert}{\rvert}
\providecommand{\norm}{}\let\norm\relax\DeclarePairedDelimiter{\norm}{\lVert}{\rVert}
\usepackage[table]{xcolor} % [table] : fournit aussi \rowcolors (alternance de lignes)
\usepackage{pagecolor}
\usepackage{tikz}
\usetikzlibrary{arrows.meta,positioning,calc,shapes.geometric,shapes.misc,shapes.arrows,decorations.pathmorphing,decorations.markings,patterns,shadows,mindmap,trees,fit,backgrounds,matrix,intersections}
\usepackage{pgfplots}
\usepackage{pgf-pie} % diagrammes circulaires \pie (statistiques)
\pgfplotsset{compat=1.18}
\usepgfplotslibrary{fillbetween,groupplots} % aires entre courbes (calcul intégral)
\usepackage{enumitem}
\usepackage{fancyhdr}
% [most] charge toutes les bibliothèques tcolorbox (listings, minted, raster,
% documentation...) et gonfle inutilement la mémoire TeX sur un document long
% (~300 boîtes) : on ne charge que ce qui est réellement utilisé.
\usepackage[skins,breakable]{tcolorbox}
\usepackage{graphicx}
\usepackage{colortbl}
\usepackage{booktabs}
\usepackage{tabularx}
\usepackage{diagbox} % case d'en-tête diagonale des tableaux à double entrée
% Colonnes extensibles centrée / gauche / droite (C, L, R), souvent utilisées par les modèles
\newcolumntype{C}{>{\centering\arraybackslash}X}
\newcolumntype{L}{>{\raggedright\arraybackslash}X}
\newcolumntype{R}{>{\raggedleft\arraybackslash}X}
\usepackage{array}
\usepackage{multicol}
\usepackage{siunitx}
% parse-numbers=false : accepte aussi \SI{2,5 \times 10^{-3}}{...}
\sisetup{locale=FR,output-decimal-marker={,},group-minimum-digits=4,parse-numbers=false}
\DeclareSIUnit\ohm{\ensuremath{\symup{\Omega}}} % U+2126 absent de la police
\usepackage{qrcode}
% \degree : commande fréquemment utilisée par les modèles pour un angle ou une
% température en dehors de \SI{}{} (non fournie par les paquets chargés).
\providecommand{\degree}{^{\circ}}
% \qed (fin de démonstration), utilisé par les modèles sans amsthm
\providecommand{\qed}{\hfill$\square$}
% \barre : double barre des valeurs interdites dans un tableau de variations (tkz-tab)
\providecommand{\barre}{\ensuremath{\|}}
% environnement proof (amsthm non chargé) utilisé par les modèles
\ifdefined\proof\else\newenvironment{proof}[1][Démonstration]{\par\noindent\textit{#1.}\ }{\qed\par}\fi
\usepackage{lastpage}
\usepackage{hyperref}
\hypersetup{hidelinks}

% ============================================================
% Palette « Pop » OnBuch+ — mêmes hex que la classe P (plus_ui.dart)
% ============================================================
\definecolor{poporange}{HTML}{F59321}
\definecolor{popdark}{HTML}{E07A0C}
\definecolor{popcream}{HTML}{FFF6E8}
\definecolor{popink}{HTML}{1C1714}
\definecolor{poppurple}{HTML}{7A5AE0}
\definecolor{popgreen}{HTML}{2E9E6A}
\definecolor{poppink}{HTML}{E0457B}
\definecolor{popblue}{HTML}{2F7DE1}
\definecolor{popgold}{HTML}{C98A12}
\definecolor{popmuted}{HTML}{8A7F76}
\definecolor{popline}{HTML}{EADFD2}
\definecolor{popgray}{HTML}{8A7F76}
\colorlet{popgrey}{popgray}
% Alias tolérants (le modèle nomme parfois les couleurs différemment)
\colorlet{popwhite}{white}
\colorlet{popblack}{popink}
\colorlet{popred}{poppink}
\colorlet{popyellow}{popgold}
% Teintes claires pour fonds de tableaux / zones
\colorlet{popblueL}{popblue!10!white}
\colorlet{poporangeL}{poporange!14!white}
\colorlet{popgreenL}{popgreen!12!white}
\colorlet{poppurpleL}{poppurple!12!white}
\colorlet{poppinkL}{poppink!12!white}
\colorlet{popgoldL}{popgold!14!white}

\pagecolor{popcream}

% ============================================================
% Typographie
% ============================================================
\defaultfontfeatures{Ligatures=TeX}
\setmainfont{PlusJakartaSans-Medium.ttf}[Path=../../../fonts/,BoldFont=PlusJakartaSans-Bold.ttf]
% Maths en sans-serif (Fira Math) avec chiffres et lettres droites de Plus
% Jakarta, pour que les formules se fondent dans le texte.
\usepackage[math-style=ISO,bold-style=ISO]{unicode-math}
\setmathfont{FiraMath-Regular.otf}[Path=../../../fonts/]
\setmathfont{PlusJakartaSans-Medium.ttf}[Path=../../../fonts/,range={up/num,up/{Latin,latin},\mathup}]
\usepackage{icomma} % virgule décimale sans espace en mode math
% Caractères mathématiques Unicode absents des polices de texte (Plus Jakarta,
% Archivo Black) : rendus par leur équivalent mathématique, en texte comme en
% formule — sinon ils s'affichent en carré vide (ex. « Arithmétique dans ℤ »).
\usepackage{newunicodechar}
\newunicodechar{ℝ}{\ensuremath{\mathbb{R}}}
\newunicodechar{ℤ}{\ensuremath{\mathbb{Z}}}
\newunicodechar{ℂ}{\ensuremath{\mathbb{C}}}
\newunicodechar{ℕ}{\ensuremath{\mathbb{N}}}
\newunicodechar{ℚ}{\ensuremath{\mathbb{Q}}}
\newunicodechar{𝔻}{\ensuremath{\mathbb{D}}}
\newunicodechar{α}{\ensuremath{\alpha}}
\newunicodechar{β}{\ensuremath{\beta}}
\newunicodechar{γ}{\ensuremath{\gamma}}
\newunicodechar{δ}{\ensuremath{\delta}}
\newunicodechar{ε}{\ensuremath{\varepsilon}}
\newunicodechar{θ}{\ensuremath{\theta}}
\newunicodechar{λ}{\ensuremath{\lambda}}
\newunicodechar{μ}{\ensuremath{\mu}}
\newunicodechar{σ}{\ensuremath{\sigma}}
\newunicodechar{τ}{\ensuremath{\tau}}
\newunicodechar{φ}{\ensuremath{\varphi}}
\newunicodechar{ψ}{\ensuremath{\psi}}
\newunicodechar{ω}{\ensuremath{\omega}}
\newunicodechar{Σ}{\ensuremath{\Sigma}}
% majuscules produites par \MakeUppercase dans les titres (α-aminés -> Α)
\newunicodechar{Α}{\ensuremath{\alpha}}
\newunicodechar{Β}{\ensuremath{\beta}}
\newunicodechar{Γ}{\ensuremath{\Gamma}}
\newunicodechar{Δ}{\ensuremath{\Delta}}
\newunicodechar{Ε}{\ensuremath{\varepsilon}}
\newunicodechar{Θ}{\ensuremath{\Theta}}
\newunicodechar{Λ}{\ensuremath{\Lambda}}
\newunicodechar{Μ}{\ensuremath{\mu}}
\newunicodechar{Τ}{\ensuremath{\tau}}
\newunicodechar{Φ}{\ensuremath{\Phi}}
\newunicodechar{Ψ}{\ensuremath{\Psi}}
\newunicodechar{Ω}{\ensuremath{\Omega}}
\newunicodechar{↦}{\ensuremath{\mapsto}}
\newunicodechar{∈}{\ensuremath{\in}}
\newunicodechar{∉}{\ensuremath{\notin}}
\newunicodechar{∧}{\ensuremath{\wedge}}
\newunicodechar{⇔}{\ensuremath{\Leftrightarrow}}
\newunicodechar{⟺}{\ensuremath{\Longleftrightarrow}}
\newunicodechar{⇒}{\ensuremath{\Rightarrow}}
\newunicodechar{⟹}{\ensuremath{\Longrightarrow}}
\newunicodechar{∘}{\ensuremath{\circ}}
\newunicodechar{∪}{\ensuremath{\cup}}
\newunicodechar{∩}{\ensuremath{\cap}}
\newunicodechar{≡}{\ensuremath{\equiv}}
\newunicodechar{⊂}{\ensuremath{\subset}}
\newunicodechar{⊥}{\ensuremath{\perp}}
\newunicodechar{∃}{\ensuremath{\exists}}
\newunicodechar{∀}{\ensuremath{\forall}}
\newunicodechar{∛}{\ensuremath{\sqrt[3]{\,}}}
\newunicodechar{∜}{\ensuremath{\sqrt[4]{\,}}}
\newunicodechar{⃗}{\ensuremath{{}^{\to}}}
\newunicodechar{ⁿ}{\ifmmode^{n}\else\textsuperscript{\ensuremath{n}}\fi}
\newunicodechar{ˣ}{\ifmmode^{x}\else\textsuperscript{\ensuremath{x}}\fi}
\newunicodechar{ᵇ}{\ifmmode^{b}\else\textsuperscript{\ensuremath{b}}\fi}
\newunicodechar{ʳ}{\ifmmode^{r}\else\textsuperscript{\ensuremath{r}}\fi}
\newunicodechar{ᵘ}{\ifmmode^{u}\else\textsuperscript{\ensuremath{u}}\fi}
\newunicodechar{ʸ}{\ifmmode^{y}\else\textsuperscript{\ensuremath{y}}\fi}
\newunicodechar{ᵃ}{\ifmmode^{a}\else\textsuperscript{\ensuremath{a}}\fi}
\newunicodechar{ᵉ}{\ifmmode^{e}\else\textsuperscript{\ensuremath{e}}\fi}
\newunicodechar{ᵏ}{\ifmmode^{k}\else\textsuperscript{\ensuremath{k}}\fi}
\newunicodechar{ᵖ}{\ifmmode^{p}\else\textsuperscript{\ensuremath{p}}\fi}
\newunicodechar{ᴺ}{\ifmmode^{N}\else\textsuperscript{\ensuremath{N}}\fi}
\newunicodechar{ᶜ}{\ifmmode^{c}\else\textsuperscript{\ensuremath{c}}\fi}
\newunicodechar{ʰ}{\ifmmode^{h}\else\textsuperscript{\ensuremath{h}}\fi}
\newunicodechar{ᵠ}{\ifmmode^{q}\else\textsuperscript{\ensuremath{q}}\fi}
\newunicodechar{ₙ}{\ifmmode_{n}\else\textsubscript{\ensuremath{n}}\fi}
\newunicodechar{ₐ}{\ifmmode_{a}\else\textsubscript{\ensuremath{a}}\fi}
\newunicodechar{ᵢ}{\ifmmode_{i}\else\textsubscript{\ensuremath{i}}\fi}
\newunicodechar{ᵣ}{\ifmmode_{r}\else\textsubscript{\ensuremath{r}}\fi}
\newunicodechar{ₖ}{\ifmmode_{k}\else\textsubscript{\ensuremath{k}}\fi}
\newunicodechar{ᵦ}{\ifmmode_{\beta}\else\textsubscript{\ensuremath{\beta}}\fi}
\newunicodechar{ₓ}{\ifmmode_{x}\else\textsubscript{\ensuremath{x}}\fi}
\newfontfamily\popdisplay{ArchivoBlack-Regular.ttf}[Path=../../../fonts/]
\newfontfamily\poplogo{SpaceGrotesk-Bold.ttf}[Path=../../../fonts/]
\newfontfamily\popsemi{PlusJakartaSans-SemiBold.ttf}[Path=../../../fonts/]
\newfontfamily\popxbold{PlusJakartaSans-ExtraBold.ttf}[Path=../../../fonts/]
\newfontfamily\popxboldls{PlusJakartaSans-ExtraBold.ttf}[Path=../../../fonts/,LetterSpace=3.0]

\setlist[enumerate]{leftmargin=1.7em,itemsep=5pt,topsep=4pt}
\setlist[itemize]{leftmargin=1.7em,itemsep=4pt,topsep=4pt,label={\color{poporange}\textbullet}}
\setlength{\parindent}{0pt}
\setlength{\parskip}{5pt}
\renewcommand{\arraystretch}{1.25}

% pgfplots : style Pop par défaut
\pgfplotsset{
  every axis/.append style={axis line style={popink,line width=1pt},tick style={popink},
    grid style={popline,line width=0.5pt},label style={font=\small},tick label style={font=\footnotesize}},
  popcurve/.style={poporange,line width=1.8pt,no markers,smooth},
}
% Même style disponible dans un \draw TikZ simple (hors pgfplots)
\tikzset{popcurve/.style={poporange,line width=1.8pt}}

% ============================================================
% Pill / squiggle
% ============================================================
\newcommand{\poppill}[3]{%
  \tikz[baseline=-0.55ex]\node[draw=popink,line width=1.2pt,rounded corners=2.6pt,%
    fill=#2,inner xsep=6.5pt,inner ysep=3pt]{%
    \popxboldls\fontsize{7}{7}\selectfont\color{#3}\MakeUppercase{#1}};%
}
\newcommand{\popsquiggle}[1]{%
  \tikz[baseline=-0.4ex]{
    \draw[#1,line width=2.6pt,line cap=round]
      plot [smooth] coordinates {(0,0.09) (0.34,0.01) (0.68,0.21) (1.02,0.01) (1.36,0.10)};
  }%
}
% Mot-clé surligné (marqueur orange sous le texte)
\newcommand{\cle}[1]{{\popxbold\color{popdark}#1}}

% ============================================================
% En-tête / pied
% ============================================================
\pagestyle{fancy}
\fancyhf{}
\renewcommand{\headrulewidth}{0pt}
\renewcommand{\footrulewidth}{0pt}
\lhead{\raisebox{-0.15ex}{\poplogo\fontsize{13}{13}\selectfont\color{popink}OnBuch\color{poporange}+}}
\rhead{\poppill{\DOCMATIERE\ \textbullet\ \DOCNIVEAU\ \textbullet\ Leçon \DOCLECON}{white}{popink}}
\lfoot{\popxbold\fontsize{7.6}{7.6}\selectfont\color{popmuted}\MakeUppercase{Cours signé OnBuch+}\ \textbullet\ {\popsemi\DOCTITRE}}
\rfoot{\tikz[baseline=-0.6ex]\node[rounded corners=8.5pt,fill=popink,minimum height=17pt,minimum width=17pt,inner xsep=5pt,inner sep=0]{\popxbold\fontsize{8}{8}\selectfont\color{popcream}\thepage\,/\,\pageref*{LastPage}};}

% ============================================================
% Titres
% ============================================================
\newcommand{\chaptitre}[2]{%
  \begin{tcolorbox}[enhanced,colback=poporange,colframe=popink,boxrule=2.6pt,arc=11pt,
      drop shadow=popink,left=15pt,right=15pt,top=12pt,bottom=11pt,before skip=3pt,after skip=18pt]
    \poppill{#2}{popink}{popcream}\par\vspace{9pt}
    {\raggedright\hyphenpenalty=10000\exhyphenpenalty=10000\popdisplay\fontsize{22}{24}\selectfont\color{popcream}\MakeUppercase{#1}\par}
  \end{tcolorbox}
}
% Section numérotée : pastille orange avec le numéro + titre Archivo
\newcounter{popsec}
\newcommand{\coursec}[1]{%
  \stepcounter{popsec}\setcounter{popsub}{0}%
  \par\vspace{16pt}\noindent
  \begin{minipage}{\linewidth}
  \tikz[baseline=-0.7ex]\node[draw=popink,line width=1.6pt,fill=poporange,rounded corners=4pt,minimum size=22pt,inner sep=0]{\popdisplay\fontsize{12}{12}\selectfont\color{popcream}\thepopsec};%
  \hspace{8pt}{\popdisplay\fontsize{15}{17}\selectfont\color{popink}\MakeUppercase{#1}}\par
  \vspace{4pt}\noindent{\color{poporange}\rule{36pt}{3.6pt}}
  \end{minipage}\par\nopagebreak\vspace{8pt}
}
\newcounter{popsub}
\newcommand{\courssub}[1]{%
  \stepcounter{popsub}%
  \par\vspace{9pt}\noindent{\popxbold\fontsize{11.5}{13}\selectfont\color{popink}{\color{poporange}\thepopsec.\thepopsub}\hspace{6pt}#1}\par\nopagebreak\vspace{4pt}
}

% ============================================================
% Boîtes pédagogiques — même recette : fond blanc, bordure épaisse colorée,
% ombre dure encre, badge plein. Titre optionnel [#1].
% ============================================================
\tcbset{popbase/.style={breakable,enhanced,colback=white,boxrule=2.3pt,arc=9pt,
  shadow={3pt}{-3pt}{0pt}{popink},left=14pt,right=14pt,top=11pt,bottom=11pt,before skip=11pt,after skip=15pt,
  fontupper=\popsemi\fontsize{10.6}{14.5}\selectfont\color{popink},
  fontlower=\popsemi\fontsize{10.6}{14.5}\selectfont\color{popink}}}
% Coche dessinée (le glyphe ✓ n'existe pas dans les polices du document)
\newcommand{\popcheck}[1]{\tikz[baseline=-0.5ex]\draw[#1,line width=1.6pt,line cap=round,line join=round](-0.13,0.0)--(-0.04,-0.09)--(0.13,0.1);}
\newcommand{\popboxhead}[3]{\poppill{#1}{#2}{white}\ifx\relax#3\relax\else\hspace{6pt}{\popxbold\fontsize{10.6}{12}\selectfont\color{popink}#3}\fi\par\vspace{7pt}}

\newtcolorbox{aretenir}[1][]{popbase,colframe=popgreen,before upper={\popboxhead{À retenir}{popgreen}{#1}}}
\newtcolorbox{exemplebox}[1][]{popbase,colframe=poporange,before upper={\popboxhead{Exemple}{poporange}{#1}}}
\newtcolorbox{definition}[1][]{popbase,colframe=popblue,before upper={\popboxhead{Définition}{popblue}{#1}}}
\newtcolorbox{propriete}[1][]{popbase,colframe=poppurple,before upper={\popboxhead{Propriété}{poppurple}{#1}}}
\newtcolorbox{methode}[1][]{popbase,colframe=popgold,before upper={\popboxhead{Méthode}{popgold}{#1}}}
\newtcolorbox{attention}[1][]{popbase,colframe=poppink,before upper={\popboxhead{Attention}{poppink}{#1}}}
\newtcolorbox{experience}[1][]{popbase,colframe=popblue,colback=popblueL,before upper={\popboxhead{Expérience}{popblue}{#1}}}
% « Le savais-tu ? » : carte violette pleine (contexte, culture, Cameroun)
\newtcolorbox{savaistu}[1][]{popbase,colback=poppurple,colframe=popink,
  fontupper=\popsemi\fontsize{10.4}{14.2}\selectfont\color{white},
  before upper={\poppill{Le savais-tu\,?}{white}{poppurple}\ifx\relax#1\relax\else\hspace{6pt}{\popxbold\color{white}#1}\fi\par\vspace{7pt}}}
% Exercice résolu : énoncé en haut, \tcblower puis la solution sur fond crème
\newcounter{popexo}\newcounter{popexr}
\newtcolorbox{exoresolu}[1][]{popbase,colframe=poporange,colbacklower=popcream,
  segmentation style={popink,line width=1.2pt,dashed},
  before upper={\stepcounter{popexr}\popboxhead{Exercice résolu \thepopexr}{poporange}{#1}},
  before lower={\poppill{Solution}{popink}{popcream}\par\vspace{6pt}}}
% Exercice (non résolu en place) — la correction est dans la partie Corrigés
\newtcolorbox{exercice}[1][]{popbase,colframe=popink,
  before upper={\stepcounter{popexo}\popboxhead{Exercice \thepopexo}{popink}{#1}}}
% Corrigé
\newtcolorbox{corrige}[1][]{popbase,colframe=popgreen,colback=popgreenL,
  before upper={\popboxhead{Corrigé}{popgreen}{#1}}}
% Objectifs / prérequis / bilan
\newtcolorbox{objectifs}{popbase,colframe=popink,colback=poporangeL,
  before upper={\popboxhead{Objectifs de la leçon}{popink}{}}}
\newtcolorbox{prerequis}{popbase,colframe=popink,colback=popblueL,
  before upper={\popboxhead{Prérequis}{popblue}{}}}
\newtcolorbox{bilan}{popbase,colframe=popink,colback=white,boxrule=2.8pt,
  before upper={\popboxhead{Fiche bilan}{poporange}{L'essentiel à connaître pour l'examen}}}
% Figure encadrée avec légende
\newtcolorbox{popfigure}[1][]{enhanced,breakable=false,colback=white,colframe=popline,boxrule=1.4pt,arc=8pt,
  left=10pt,right=10pt,top=10pt,bottom=8pt,before skip=10pt,after skip=12pt,halign=center,
  lower separated=false,halign lower=center,
  fontlower=\popsemi\fontsize{9}{11.5}\selectfont\color{popmuted},
  #1}
\newcommand{\legende}[1]{\par\vspace{6pt}{\popsemi\fontsize{9}{11.5}\selectfont\color{popmuted}#1}}

% ============================================================
% Signature OnBuch+
% ============================================================
\newcommand{\poplogoinline}[1]{{\poplogo\fontsize{#1}{#1}\selectfont\color{popink}OnBuch\color{poporange}+}}
\newcommand{\signatureonbuch}{%
  \begin{tcolorbox}[enhanced,colback=popink,colframe=popink,boxrule=2.2pt,arc=10pt,
      drop shadow=poporange,left=14pt,right=14pt,top=10pt,bottom=10pt,before skip=0pt,after skip=0pt]
    \tikz[baseline=-0.7ex]\node[circle,fill=popgreen,draw=popcream,line width=1.3pt,minimum size=22pt,inner sep=0]{\popcheck{white}};%
    \hspace{9pt}\begin{minipage}[c]{0.62\linewidth}
      {\popxbold\fontsize{12}{13}\selectfont\color{popcream}Signé\ \textbullet\ {\poplogo OnBuch\color{poporange}+}}\par\vspace{2pt}
      {\raggedright\popsemi\fontsize{8}{10.5}\selectfont\color{popline}\MakeUppercase{Équipe pédagogique OnBuch+}\\\MakeUppercase{Conforme au programme officiel MINESEC}\par}
    \end{minipage}\hfill
    {\poplogo\fontsize{22}{22}\selectfont\color{popcream}OnBuch\color{poporange}+}
  \end{tcolorbox}%
}

% ============================================================
% Page de garde
% #1 titre · #2 matière · #3 niveau · #4 module · #5 n° leçon · #6 durée ·
% #7 description / objectifs (texte ou itemize)
% ============================================================
\newcommand{\pagegarde}[7]{%
  \thispagestyle{empty}
  \newgeometry{a4paper,margin=1.7cm,top=1.6cm,bottom=1.4cm}%
  \begin{tikzpicture}[remember picture,overlay]
    \fill[poporange!18] ([xshift=-3.2cm,yshift=-3.6cm]current page.north east) circle (4.6cm);
    \fill[poppurple!14] ([xshift=2.4cm,yshift=7.6cm]current page.south west) circle (3.8cm);
  \end{tikzpicture}%
  {\poplogo\fontsize{30}{30}\selectfont\color{popink}OnBuch\color{poporange}+}\hfill
  \raisebox{0.6ex}{\poppill{Cours signé \textbullet\ Premium}{popink}{popcream}}\par
  \vspace{1pt}\popsquiggle{poppurple}\par
  \vspace{8pt}
  \begin{tcolorbox}[enhanced,colback=poporange,colframe=popink,boxrule=2.8pt,arc=13pt,
      drop shadow=popink,left=18pt,right=18pt,top=12pt,bottom=13pt,after skip=10pt]
    \poppill{#2 \textbullet\ #3}{popink}{popcream}\hspace{5pt}\poppill{#4}{white}{popink}\par\vspace{8pt}
    {\popdisplay\fontsize{14}{14}\selectfont\color{popink}LEÇON #5}\par\vspace{4pt}
    {\raggedright\hyphenpenalty=10000\exhyphenpenalty=10000\popdisplay\fontsize{25}{27}\selectfont\color{popcream}\MakeUppercase{#1}\par}
    \vspace{8pt}
    {\popxbold\fontsize{9.5}{9.5}\selectfont\color{popink}\MakeUppercase{Durée conseillée : #6}}
  \end{tcolorbox}
  \begin{tcolorbox}[enhanced,colback=white,colframe=popink,boxrule=2.2pt,arc=10pt,
      drop shadow=popink,left=16pt,right=16pt,top=10pt,bottom=10pt,after skip=8pt]
    \poppill{Dans cette leçon}{poppurple}{white}\par\vspace{5pt}
    \popsemi\fontsize{9.3}{12.6}\selectfont\color{popink}#7
  \end{tcolorbox}
  \noindent\begin{minipage}[t]{0.487\linewidth}
    \begin{tcolorbox}[enhanced,colback=popgreen,colframe=popink,boxrule=2.2pt,arc=10pt,
        drop shadow=popink,left=11pt,right=11pt,top=10pt,bottom=10pt,height=3.1cm,valign=center]
      \tcbox[colback=white,colframe=popink,boxrule=1.3pt,arc=5pt,boxsep=0pt,nobeforeafter,
        left=3pt,right=3pt,top=3pt,bottom=3pt,valign=center]{\qrcode[height=1.9cm]{https://play.google.com/store/apps/details?id=com.luvvix.onbuch&hl=fr}}%
      \hspace{8pt}\begin{minipage}[c]{0.5\linewidth}
        {\popdisplay\fontsize{11}{12.5}\selectfont\color{white}\MakeUppercase{Télécharge OnBuch}}\par\vspace{4pt}
        {\raggedright\popsemi\fontsize{8.4}{11}\selectfont\color{white}Gratuit \textbullet\ Google Play\\Tuteur IA, annales, concours, résultats\par}
      \end{minipage}
    \end{tcolorbox}
  \end{minipage}\hfill
  \begin{minipage}[t]{0.487\linewidth}
    \begin{tcolorbox}[enhanced,colback=poppurple,colframe=popink,boxrule=2.2pt,arc=10pt,
        drop shadow=popink,left=11pt,right=11pt,top=10pt,bottom=10pt,height=3.1cm,valign=center]
      \tikz[baseline=-0.7ex]\node[circle,fill=white,draw=popink,line width=1.3pt,minimum size=1.6cm,inner sep=0]{\popdisplay\fontsize{24}{24}\selectfont\color{poppurple}+};%
      \hspace{8pt}\begin{minipage}[c]{0.6\linewidth}
        {\popdisplay\fontsize{11}{12.5}\selectfont\color{white}\MakeUppercase{Passe à OnBuch+}}\par\vspace{4pt}
        {\raggedright\popsemi\fontsize{8.4}{11}\selectfont\color{white}Tous les cours signés, Tuteur IA illimité, fiches PDF — onglet OnBuch+ de l'app\par}
      \end{minipage}
    \end{tcolorbox}
  \end{minipage}
  \vspace{8pt}
  \popcontenu
  \vfill
  \signatureonbuch
  \restoregeometry
  \newpage
}

% Rangée « Ce cours contient » (page de garde)
\newcommand{\poptile}[3]{% #1 couleur #2 gros texte #3 légende
  \begin{tcolorbox}[enhanced,colback=#1,colframe=popink,boxrule=2pt,arc=9pt,shadow={2.5pt}{-2.5pt}{0pt}{popink},
      left=6pt,right=6pt,top=9pt,bottom=9pt,halign=center,valign=center,height=1.75cm,width=\linewidth,nobeforeafter]
    {\popdisplay\fontsize{12.5}{13}\selectfont\color{white}\MakeUppercase{#2}}\par\vspace{3pt}
    {\centering\popsemi\fontsize{7.8}{9.5}\selectfont\color{white}#3\par}
  \end{tcolorbox}}
\newcommand{\popcontenu}{%
  \par\noindent{\popxboldls\fontsize{7.5}{7.5}\selectfont\color{popmuted}CE COURS SIGNÉ CONTIENT}\par\vspace{7pt}
  \noindent\begin{minipage}[t]{0.235\linewidth}\poptile{popblue}{Cours}{complet, illustré, pas à pas}\end{minipage}\hfill
  \begin{minipage}[t]{0.235\linewidth}\poptile{poporange}{Méthodes}{et exercices résolus}\end{minipage}\hfill
  \begin{minipage}[t]{0.235\linewidth}\poptile{poppink}{Exercices}{type examen + corrigés}\end{minipage}\hfill
  \begin{minipage}[t]{0.235\linewidth}\poptile{popgreen}{Fiche bilan}{l'essentiel à retenir}\end{minipage}\par}

% Sommaire
\newtcolorbox{sommaire}{enhanced,colback=white,colframe=popink,boxrule=2.2pt,arc=10pt,
  drop shadow=popink,left=18pt,right=18pt,top=13pt,bottom=13pt,before skip=6pt,after skip=20pt,
  before upper={\poppill{Sommaire}{poppurple}{white}\par\vspace{9pt}\popsemi\fontsize{10.6}{14.5}\selectfont\color{popink}}}

% Signature de fin de cours — poussée tout en bas de la dernière page
\newcommand{\findecours}{%
  \par\vfill\nopagebreak
  \begin{center}\popsquiggle{poporange}\end{center}\vspace{4pt}
  \signatureonbuch
}
\AtBeginDocument{\def\llbracket{\mathopen{[\mkern-3.5mu[}}\def\rrbracket{\mathclose{]\mkern-3.5mu]}}} % intervalles d'entiers (glyphe ⟦ absent de la police)
% Glyphes absents de Fira Math : redéfinis à partir de glyphes disponibles
\AtBeginDocument{%
  \def\setminus{\mathbin{\backslash}}\let\smallsetminus\setminus
  \def\vdots{\mathord{\vbox{\baselineskip3.2pt\lineskiplimit0pt\kern4pt\hbox{.}\hbox{.}\hbox{.}}}}%
  \def\ddots{\mathinner{\mkern1mu\raise6.5pt\hbox{.}\mkern2mu\raise3.6pt\hbox{.}\mkern2mu\raise0.7pt\hbox{.}\mkern1mu}}%
  \def\triangle{\mathord{\scalebox{0.9}{$\symup{\Delta}$}}}%
  \def\nabla{\mathord{\raisebox{\depth}{\scalebox{1}[-1]{$\symup{\Delta}$}}}}%
  \def\checkmark{\popcheck{popdark}}%
  \def\star{\mathbin{\ast}}\def\bigstar{\mathord{\ast}}%
  \def\diamond{\mathbin{\scalebox{0.6}{$\Diamond$}}}%
}

%% FIGFIX-BEGIN v4 — correctifs globaux des figures (docs/onbuch-generation/tools/figfix)
\usepackage{adjustbox}\usepackage{etoolbox}
% toute figure TikZ/pgfplots plus large que la ligne est réduite à la largeur disponible
\BeforeBeginEnvironment{tikzpicture}{\begin{adjustbox}{max width=\linewidth}}
\AfterEndEnvironment{tikzpicture}{\end{adjustbox}}
% légendes pgfplots lisibles par-dessus les courbes
\pgfplotsset{every axis/.append style={legend style={fill=white,fill opacity=0.92,text opacity=1,draw=black!15,inner sep=3pt}}}
% étiquettes d'arêtes (syntaxe edge["..."]) avec fond blanc
\tikzset{every edge quotes/.append style={fill=white,inner sep=1.2pt}}
%% FIGFIX-END
\begin{document}

\pagegarde{Coniques}{Mathématiques}{Tle E}{Configurations et transformations élémentaires du plan}{15}{28 h}{Cette leçon présente les coniques (ellipse, parabole, hyperbole) comme sections planes d'un cône de révolution, puis par leur définition monofocale MF = e·d(M,D). Elle conduit à leurs équations réduites, à leurs éléments caractéristiques, à leurs tangentes et à la reconnaissance d'équations du second degré, avec des applications concrètes (antennes, orbites, ponts).
\vspace{4pt}
\begin{multicols}{2}\small
\begin{itemize}
\item À la fin de cette leçon, je sais décrire les coniques comme sections planes d'un cône de révolution
\item À la fin de cette leçon, je sais identifier le type d'une conique à partir de son excentricité e
\item À la fin de cette leçon, je sais construire point par point une conique à partir de sa définition monofocale
\item À la fin de cette leçon, je sais déterminer l'équation réduite d'une conique et ses éléments caractéristiques (foyers, directrices, sommets, axes)
\item À la fin de cette leçon, je sais déterminer l'équation de la tangente à une conique en un point donné
\item À la fin de cette leçon, je sais mettre sous forme réduite une équation du type ax² + by² + cx + dy + e = 0 et reconnaître la conique associée
\end{itemize}\end{multicols}}

\begin{sommaire}
\begin{multicols}{2}
\begin{enumerate}
\item Coniques comme sections d'un cône
\item Définition monofocale et excentricité
\item Équations réduites et éléments caractéristiques
\item Tangentes à une conique
\item Reconnaissance d'une équation du second degré
\item Coniques, similitudes et situations concrètes
\item Activité d'intégration
\item Méthodes et exercices résolus
\item Exercices
\item Corrigés des exercices
\item Fiche bilan
\end{enumerate}
\end{multicols}
\end{sommaire}

\begin{objectifs}
\begin{itemize}
\item À la fin de cette leçon, je sais décrire les coniques comme sections planes d'un cône de révolution
\item À la fin de cette leçon, je sais identifier le type d'une conique à partir de son excentricité e
\item À la fin de cette leçon, je sais construire point par point une conique à partir de sa définition monofocale
\item À la fin de cette leçon, je sais déterminer l'équation réduite d'une conique et ses éléments caractéristiques (foyers, directrices, sommets, axes)
\item À la fin de cette leçon, je sais déterminer l'équation de la tangente à une conique en un point donné
\item À la fin de cette leçon, je sais mettre sous forme réduite une équation du type ax² + by² + cx + dy + e = 0 et reconnaître la conique associée
\item À la fin de cette leçon, je sais déterminer l'image d'une conique par une similitude directe
\item À la fin de cette leçon, je sais reconnaître des coniques dans des situations concrètes (antennes paraboliques, orbites, ponts, ombres)
\end{itemize}
\end{objectifs}

\begin{prerequis}
\begin{itemize}
\item Distance d'un point à une droite et produit scalaire (Première)
\item Équations de droites, cercles et repérage dans le plan (Première)
\item Similitudes directes : définition, propriétés, effet sur les distances (leçon sur les transformations)
\item Formes canoniques des trinômes du second degré (Seconde/Première)
\item Trigonométrie et coordonnées paramétriques élémentaires (Première C/E)
\end{itemize}
\end{prerequis}

\newpage

\coursec{Coniques comme sections d'un cône}

À Yaoundé, une entreprise d'installation d'antennes paraboliques doit dimensionner un récepteur satellite : où placer exactement le récepteur pour capter le maximum de signal ? Au même moment, un ingénieur du port de Douala étudie la trajectoire d'un satellite au-dessus du Cameroun, et un architecte dessine l'arche d'un pont à Mbalmayo. Antenne, orbite, arche : trois objets très différents, trois courbes en apparence sans rapport. Et pourtant, chacune d'elles peut être obtenue en tranchant un même objet géométrique : un \cle{cône}. Comment définir, construire et reconnaître ces courbes appelées \cle{coniques} ? C'est tout l'objet de cette leçon, et nous commençons par leur origine géométrique : la section d'un cône par un plan.

\courssub{Le cône de révolution}

\textbf{Intuition.} Prends une droite $(\Delta)$ verticale et une autre droite $(d)$ qui la coupe en un point $S$ en faisant un angle constant avec elle. Fais maintenant tourner $(d)$ autour de $(\Delta)$, comme le manche d'une toupie ou le bras d'un tour de potier. La droite $(d)$ balaie une surface en forme de double entonnoir, infini vers le haut et vers le bas : c'est un cône de révolution à deux nappes.

\begin{definition}[Cône de révolution]
Soit $(\Delta)$ une droite, $S$ un point de $(\Delta)$ et $\alpha$ un réel tel que $0 < \alpha < \dfrac{\pi}{2}$. Le \cle{cône de révolution} de \cle{sommet} $S$, d'\cle{axe} $(\Delta)$ et de \cle{demi-angle au sommet} $\alpha$ est la surface engendrée par la rotation autour de $(\Delta)$ d'une droite $(d)$ passant par $S$ et faisant avec $(\Delta)$ un angle constant $\alpha$.
\begin{itemize}
\item Les droites obtenues par rotation de $(d)$ sont appelées les \cle{génératrices} du cône.
\item Le cône est formé de deux \cle{nappes} symétriques par rapport au sommet $S$.
\item Tout plan perpendiculaire à l'axe $(\Delta)$, ne passant pas par $S$, coupe le cône selon un \cle{cercle}.
\end{itemize}
\end{definition}

\begin{attention}[Le cône a deux nappes]
En géométrie, le cône n'est pas le « chapeau » limité de la vie courante : ses génératrices sont des \emph{droites} entières, donc le cône s'étend indéfiniment des deux côtés du sommet. Cette remarque est essentielle : c'est l'existence des deux nappes qui explique que l'hyperbole possède deux branches.
\end{attention}

\courssub{Sections planes d'un cône : les coniques}

\textbf{Intuition.} Éclaire un mur avec une torche : le faisceau lumineux est (approximativement) un cône, et le mur est un plan. La tache lumineuse dessinée sur le mur est précisément l'intersection du cône avec ce plan. Selon que tu diriges la torche perpendiculairement au mur, de biais, ou presque parallèlement au mur, la tache change de forme : cercle, ovale, courbe ouverte... Toutes les coniques naissent ainsi.

\begin{definition}[Conique]
On appelle \cle{conique} toute courbe obtenue comme intersection d'un cône de révolution et d'un plan. Lorsque le plan ne passe pas par le sommet du cône, la conique est dite \cle{non dégénérée} : c'est une ellipse, une parabole ou une hyperbole.
\end{definition}

La nature de la section dépend uniquement de la comparaison entre deux angles : le demi-angle au sommet $\alpha$ du cône, et l'angle $\beta$ que fait le plan de coupe avec un plan perpendiculaire à l'axe du cône (c'est le complémentaire de l'angle entre le plan et l'axe lui-même).

\begin{propriete}[Nature de la section selon l'inclinaison du plan (admise)]
Soit un cône de révolution de sommet $S$ et de demi-angle au sommet $\alpha$, et un plan $(P)$ \textbf{ne passant pas par} $S$. Notons $\beta$ l'angle que fait $(P)$ \textbf{avec un plan perpendiculaire à l'axe} (donc $\beta = 0$ signifie que $(P)$ est perpendiculaire à l'axe).
\begin{itemize}
\item Si $\beta = 0$ (plan perpendiculaire à l'axe) : la section est un \cle{cercle}.
\item Si $0 < \beta < \alpha$ (plan moins incliné que les génératrices) : la section est une \cle{ellipse} (courbe fermée, le cercle en est un cas particulier).
\item Si $\beta = \alpha$ (plan parallèle à une génératrice) : la section est une \cle{parabole} (courbe ouverte, une seule branche, une seule nappe coupée).
\item Si $\beta > \alpha$ (plan plus incliné que les génératrices) : la section est une \cle{hyperbole} (le plan coupe les deux nappes, la courbe a deux branches).
\end{itemize}
\end{propriete}

\begin{popfigure}
\begin{tikzpicture}[scale=0.85, line cap=round, >=Stealth]
% Styles
\tikzset{cone/.style={fill=popblueL, opacity=0.5, draw=popblue, thick},
         plan/.style={draw=poporange, thick},
         curve/.style={draw=popink, very thick, smooth},
         axis/.style={draw=popmuted, dashed, thin}}
% ---- Ellipse ----
\begin{scope}[xshift=0cm]
  % Nappes du cône (triangle isocèle)
  \fill[popblueL, opacity=0.55] (0,2.6) -- (-1.3,0) -- (1.3,0) -- cycle;
  \fill[popblueL, opacity=0.55] (0,-2.6) -- (-1.3,0) -- (1.3,0) -- cycle;
  \draw[popblue, thick] (0,2.6) -- (-1.3,0) (0,2.6) -- (1.3,0);
  \draw[popblue, thick] (0,-2.6) -- (-1.3,0) (0,-2.6) -- (1.3,0);
  \draw[axis] (0,-2.8) -- (0,2.8) node[above, font=\tiny, popmuted] {Axe};
  % Plan de coupe (incliné, beta < alpha)
  \draw[plan] (-1.5,0.5) -- (1.5,1.5);
  % Ellipse sur le plan (représentée par un arc d'ellipse incliné)
  \draw[curve, rotate around={15:(0,1)}] (-0.6,1) ellipse [x radius=0.6, y radius=0.15];
  \node[popdark, font=\small\bfseries] at (0,-3.1) {Ellipse ($\beta < \alpha$)};
\end{scope}
% ---- Parabole ----
\begin{scope}[xshift=5.2cm]
  \fill[popgreenL, opacity=0.55] (0,2.6) -- (-1.3,0) -- (1.3,0) -- cycle;
  \fill[popgreenL, opacity=0.55] (0,-2.6) -- (-1.3,0) -- (1.3,0) -- cycle;
  \draw[popgreen!70!black, thick] (0,2.6) -- (-1.3,0) (0,2.6) -- (1.3,0);
  \draw[popgreen!70!black, thick] (0,-2.6) -- (-1.3,0) (0,-2.6) -- (1.3,0);
  \draw[axis] (0,-2.8) -- (0,2.8);
  % Plan parallèle à la génératrice droite (pente -2 en coords (x,y) -> generatrice eq y = -2x+2.6? Non.
  % Generatrice droite : (0,2.6) to (1.3,0). Pente = -2.6/1.3 = -2. Eq: y = -2x + 2.6.
  % Plan parallèle : y = -2x + c. Prend c=0.2 pour couper la nappe sup.
  \draw[plan] (-1.15,-1.4) -- (0.35,1.4); % y = 2x + 1.4? Non, pente -2.
  % Points: x=-1.15 -> y=2*1.15+1.4=3.7 trop haut. x=0.35 -> y=-0.7+1.4=0.7.
  % Refaire : Plan y = -2x + 1.0. x=-1 -> y=3. x=0.5 -> y=0.
  \draw[plan] (-1.1, 3.2) -- (0.6, -0.2); % Coupe la nappe sup.
  % Parabole sur le plan : param eq dans le plan. Projection 2D : parabole standard.
  % On trace une courbe y = a(x-h)^2 + k approximative sur le segment du plan.
  \draw[curve, domain=-0.9:0.3, samples=60, variable=\t]
    plot ({\t}, {-2*\t + 1.0 - 1.5*(\t+0.3)*(\t+0.3)}); % Parabole concave vers le bas dans le repère du plan
  \node[popdark, font=\small\bfseries] at (0,-3.1) {Parabole ($\beta = \alpha$)};
\end{scope}
% ---- Hyperbole ----
\begin{scope}[xshift=10.4cm]
  \fill[poppurpleL, opacity=0.55] (0,2.6) -- (-1.3,0) -- (1.3,0) -- cycle;
  \fill[poppurpleL, opacity=0.55] (0,-2.6) -- (-1.3,0) -- (1.3,0) -- cycle;
  \draw[poppurple, thick] (0,2.6) -- (-1.3,0) (0,2.6) -- (1.3,0);
  \draw[poppurple, thick] (0,-2.6) -- (-1.3,0) (0,-2.6) -- (1.3,0);
  \draw[axis] (0,-2.8) -- (0,2.8);
  % Plan vertical (beta > alpha) : x = 0.45
  \draw[plan] (0.45,-2.2) -- (0.45,2.2);
  % Hyperbole : intersection x=0.45 avec cone x^2 = (0.5 y)^2 -> 0.45^2 = 0.25 y^2 - z^2 (3D)
  % En vue de côté (plan xOy), la courbe a eq x = 0.45 +/- sqrt(0.25 y^2 - 0.45^2)
  % Branche droite (x > 0.45)
  \draw[curve, domain=0.9:2.2, samples=50, variable=\t]
    plot ({0.45 + sqrt(max(0,0.25*\t*\t - 0.2025))}, {\t});
  \draw[curve, domain=-2.2:-0.9, samples=50, variable=\t]
    plot ({0.45 + sqrt(max(0,0.25*\t*\t - 0.2025))}, {\t});
  % Branche gauche (x < 0.45)
  \draw[curve, domain=0.9:2.2, samples=50, variable=\t]
    plot ({0.45 - sqrt(max(0,0.25*\t*\t - 0.2025))}, {\t});
  \draw[curve, domain=-2.2:-0.9, samples=50, variable=\t]
    plot ({0.45 - sqrt(max(0,0.25*\t*\t - 0.2025))}, {\t});
  \node[popdark, font=\small\bfseries] at (0,-3.1) {Hyperbole ($\beta > \alpha$)};
\end{scope}
\end{tikzpicture}
\legende{Un cône à deux nappes coupé par trois plans : selon l'inclinaison $\beta$ du plan par rapport au plan perpendiculaire à l'axe (demi-angle au sommet $\alpha$), on obtient une ellipse, une parabole ou une hyperbole.}
\end{popfigure}

\begin{experience}[Découvrir la classification avec une lampe torche]
Matériel : une lampe torche (ou le flash d'un téléphone) et un mur, dans une pièce sombre.
\begin{enumerate}
\item Dirige la torche \emph{perpendiculairement} au mur : la tache lumineuse est un \textbf{cercle}.
\item Incline légèrement la torche : la tache s'allonge en \textbf{ellipse}. Plus tu inclines, plus l'ellipse est allongée.
\item Incline davantage jusqu'à ce que le bord du faisceau soit \emph{parallèle au mur} : la tache devient une courbe ouverte qui « part à l'infini » d'un seul côté : c'est une \textbf{parabole}.
\item Incline encore : le faisceau éclaire à la fois le mur et, par exemple, le plafond ou le sol ; sur chaque surface apparaît une branche : c'est une \textbf{hyperbole}.
\end{enumerate}
Conclusion : le passage ellipse $\to$ parabole $\to$ hyperbole se fait de façon continue quand on fait pivoter le plan de coupe ; la parabole est le cas limite, le « moment exact » où le plan devient parallèle à une génératrice.
\end{experience}

\begin{attention}[Parallèle à une génératrice $\neq$ simplement incliné]
L'erreur classique consiste à dire : « le plan est incliné, donc c'est une parabole ». Faux ! Un plan simplement incliné (mais moins que les génératrices, c'est-à-dire $\beta < \alpha$) donne une \textbf{ellipse}. La parabole n'apparaît que pour l'inclinaison \emph{exacte} $\beta = \alpha$, c'est-à-dire lorsque le plan est \textbf{parallèle à une génératrice} du cône. C'est un cas limite, unique, qui sépare les ellipses des hyperboles. Retiens : incliné $\neq$ parallèle à une génératrice.
\end{attention}

\courssub{Les cas dégénérés : quand le plan passe par le sommet}

Que se passe-t-il si le plan de coupe passe par le sommet $S$ du cône ? La section n'est plus une belle courbe, mais une figure « effondrée » :

\begin{itemize}
\item Si le plan ne rencontre le cône qu'en $S$ (plan trop peu incliné, $\beta < \alpha$) : la section est réduite à un \cle{point}.
\item Si le plan est parallèle à une génératrice et passe par $S$ ($\beta = \alpha$) : il est \emph{tangent} au cône le long d'une génératrice ; la section est une \cle{droite} (une génératrice).
\item Si le plan est plus incliné ($\beta > \alpha$) : il traverse les deux nappes en passant par $S$ ; la section est formée de \cle{deux droites sécantes} en $S$ (deux génératrices).
\end{itemize}

Ces figures sont appelées \cle{coniques dégénérées}. On les mentionne pour être complet, mais dans cette leçon, sauf indication contraire, le mot « conique » désignera toujours une conique \emph{non dégénérée} : ellipse, parabole ou hyperbole.

\begin{aretenir}[Bilan de la classification]
\begin{center}
\begin{tabularx}{\linewidth}{|l|X|X|}
\hline
\rowcolor{popblueL} \textbf{Position du plan} & \textbf{Section non dégénérée} & \textbf{Section dégénérée (plan par $S$)} \\
\hline
Perpendiculaire à l'axe ($\beta=0$) & Cercle & Point $S$ \\
\hline
$0 < \beta < \alpha$ & Ellipse & Point $S$ \\
\hline
$\beta = \alpha$ (parallèle à une génératrice) & Parabole & Une droite \\
\hline
$\beta > \alpha$ & Hyperbole (deux branches) & Deux droites sécantes \\
\hline
\end{tabularx}
\end{center}
\end{aretenir}

\courssub{Les coniques autour de nous}

\begin{exemplebox}[L'ombre du lampadaire et la lumière de la torche]
Le soir, à Ngaoundéré, l'ampoule d'un lampadaire est assimilable à un point lumineux $S$. Le poteau du lampadaire bloque la lumière : l'ombre du poteau sur le sol est délimitée par l'intersection du cône de lumière (de sommet $S$, s'appuyant sur le sommet du poteau) avec le plan du sol. Selon l'inclinaison, le bord de l'ombre est une branche d'\textbf{hyperbole}. De même, une torche dirigée obliquement vers un mur dessine une ellipse, et un projecteur de chantier réglé de façon que le bord de son faisceau soit parallèle au mur dessine une parabole.
\end{exemplebox}

\begin{popfigure}
\begin{tikzpicture}[scale=1.0, line cap=round, >=Stealth]
% Mur (plan vertical vu de côté, dessiné en façade)
\fill[popcream] (2.6,-1.8) rectangle (6.4,2.6);
\draw[popmuted] (2.6,-1.8) rectangle (6.4,2.6);
\node[popmuted, font=\small] at (4.5,2.35) {Mur};
% Sol
\draw[popmuted, thick] (-1.5,-1.8) -- (6.4,-1.8);
\node[popmuted, font=\small] at (-0.6,-2.05) {Sol};
% Torche (sommet S)
\fill[popink] (0,0) circle (0.07);
\node[popdark, font=\small, below left] at (0,0) {$S$ (torche)};
% Faisceau : cône stylisé vers le mur (generatrices)
\draw[poporange, thick] (0,0) -- (2.6,1.9); % generatrice sup
\draw[poporange, thick] (0,0) -- (2.6,-0.5); % generatrice inf
\fill[poporangeL, opacity=0.3] (0,0) -- (2.6,1.9) -- (2.6,-0.5) -- cycle;
% Deuxième nappe (vers le bas) qui frappe le sol
\draw[poporange, thick, dashed] (0,0) -- (3.6,-1.8);
\draw[poporange, thick, dashed] (0,0) -- (1.5,-1.8);
\fill[poporangeL, opacity=0.2] (0,0) -- (3.6,-1.8) -- (1.5,-1.8) -- cycle;
% Hyperbole sur le mur : intersection plan x=2.6 avec cone.
% Cone: y^2+z^2 = (k x)^2. Plan x=2.6. Hyperbole verticale.
% En 2D (x horizontal, y vertical), z est la profondeur. Courbe sur le mur : y = +/- sqrt((k*2.6)^2 - z^2) -> Ellipse? Non.
% Si torche à (0,0,0), axe horizontal (x). Cone: y^2+z^2 = (tan(alpha) x)^2.
% Mur x=2.6. Section: y^2+z^2 = C^2 -> Cercle si axe perpendiculaire.
% Ici torche inclinée. Disons axe horizontal. Mur vertical (plan x=cte). Section = Cercle/Ellipse.
% Pour avoir hyperbole sur mur, il faut que le plan du mur coupe les deux nappes.
% Si torche pointe vers le mur (axe horizontal), le cône est horizontal. Le mur (vertical) est perpendiculaire à l'axe -> Cercle.
% Pour avoir hyperbole, il faut que le mur soit // à l'axe (vertical) et que le cône soit incliné.
% Le schéma est une vue de côté (plan xOy). z=0.
% Torche à (0,0). Faisceau : lignes y = mx et y = -mx (generatrices dans le plan z=0).
% Mur : x = 2.6 (vertical). Intersection : points (2.6, 2.6m) et (2.6, -2.6m). Juste deux points.
% L'hyperbole apparait dans la 3eme dimension (z).
% Le dessin 2D ne peut montrer l'hyperbole sur le mur que schématiquement.
% On garde l'esprit du schéma : deux branches, une sur le mur, une au sol.
% Branche sur le mur (x=2.6, y variable, z variable). Projection sur xOy : segment vertical.
% Pour "voir" l'hyperbole, on la dessine décalée.
% Branche 1 (mur) : x = 2.6 + epsilon * f(y). f(y) = sqrt(y^2 - a^2).
\draw[popink, very thick, domain=1.2:2.4, samples=60, smooth, variable=\t]
  plot ({2.6 + 0.3*sqrt(max(0,\t*\t - 1.0))}, {\t});
\draw[popink, very thick, domain=-2.4:-1.2, samples=60, smooth, variable=\t]
  plot ({2.6 + 0.3*sqrt(max(0,\t*\t - 1.0))}, {\t});
\node[popink, font=\small] at (4.9,1.2) {Branche sur le mur};
% Branche au sol (y=-1.8, x variable, z variable). Projection sur xOy : courbe y=-1.8.
% Equation sol: y=-1.8. Cone: (-1.8)^2 + z^2 = (k x)^2 -> z^2 = k^2 x^2 - 3.24. Hyperbole.
% Dans le plan xOy (z=0), on voit l'intersection de la nappe avec z=0 : lignes.
% Pour montrer la branche, on la dessine au-dessus du sol.
\draw[popink, very thick, domain=2.8:5.5, samples=60, smooth, variable=\t]
  plot ({\t}, {-1.8 + 0.8/sqrt(max(0,\t-2.5))}); % Aspect hyperbole
\node[popink, font=\small] at (5.1,-1.35) {Branche au sol};
\end{tikzpicture}
\legende{Le faisceau conique d'une torche rencontre le mur et le sol : chaque surface plane porte une branche d'hyperbole. En inclinant différemment la torche, on obtiendrait une ellipse ou une parabole.}
\end{popfigure}

\begin{exemplebox}[Trajectoires et constructions]
\begin{itemize}
\item L'\textbf{antenne parabolique} de l'installateur de Yaoundé est une surface obtenue par rotation d'une \emph{parabole} autour de son axe : tous les rayons venant du satellite (parallèles à l'axe) se réfléchissent vers un unique point, le foyer, où l'on place le récepteur.
\item La \textbf{trajectoire d'un satellite} autour de la Terre, suivie depuis le port de Douala, est une \emph{ellipse} (première loi de Kepler) ; celle d'une comète qui ne revient jamais est une branche d'\emph{hyperbole}.
\item L'\textbf{arche du pont} de Mbalmayo et les câbles de certains ponts suspendus suivent des arcs de \emph{parabole}, formes réputées pour leur bonne répartition des efforts.
\end{itemize}
\end{exemplebox}

\begin{exoresolu}[Identifier une section conique]
Un cône de révolution a pour demi-angle au sommet $\alpha = 30^{\circ}$. On le coupe par un plan $(P)$ ne passant pas par le sommet. Dans chaque cas, déterminer la nature de la section :
\begin{enumerate}
\item $(P)$ est perpendiculaire à l'axe du cône ;
\item $(P)$ fait un angle de $20^{\circ}$ avec un plan perpendiculaire à l'axe ;
\item $(P)$ est parallèle à une génératrice du cône ;
\item $(P)$ est parallèle à l'axe du cône.
\end{enumerate}
\tcblower
On compare dans chaque cas l'inclinaison $\beta$ du plan (par rapport à la position perpendiculaire à l'axe) au demi-angle $\alpha = 30^{\circ}$.
\begin{enumerate}
\item Le plan est perpendiculaire à l'axe : $\beta = 0$, cas particulier de l'ellipse : la section est un \textbf{cercle}.
\item $\beta = 20^{\circ} < \alpha = 30^{\circ}$ : le plan est moins incliné que les génératrices, il ne coupe qu'une nappe selon une courbe fermée : la section est une \textbf{ellipse} (non réduite à un cercle).
\item Le plan est parallèle à une génératrice : $\beta = \alpha = 30^{\circ}$ exactement : la section est une \textbf{parabole}.
\item Un plan parallèle à l'axe est plus incliné que toutes les génératrices ($\beta = 90^{\circ} > \alpha$) : il coupe les deux nappes, la section est une \textbf{hyperbole}.
\end{enumerate}
\end{exoresolu}

\begin{exercice}[Exercice d'application]
Un cône de révolution a pour demi-angle au sommet $\alpha = 45^{\circ}$ (cône \textbf{rectangle}). Pour chaque plan suivant, préciser la nature de la section, en justifiant par la comparaison des angles :
\begin{enumerate}
\item un plan faisant un angle de $45^{\circ}$ avec l'axe du cône ;
\item un plan faisant un angle de $60^{\circ}$ avec l'axe du cône ;
\item un plan faisant un angle de $30^{\circ}$ avec l'axe du cône.
\end{enumerate}
(Attention : ici les angles sont mesurés par rapport à l'\emph{axe}, et non par rapport au plan perpendiculaire à l'axe.)
\end{exercice}

\begin{corrige}[Exercice d'application]
Si $\theta$ désigne l'angle entre le plan $(P)$ et l'axe du cône, alors l'angle entre $(P)$ et le plan perpendiculaire à l'axe est $\beta = 90^{\circ} - \theta$. Le plan est parallèle à une génératrice lorsque $\theta = \alpha = 45^{\circ}$.
\begin{enumerate}
\item $\theta = 45^{\circ} = \alpha$ : le plan est parallèle à une génératrice, la section est une \textbf{parabole}.
\item $\theta = 60^{\circ} > \alpha$ : le plan est « trop redressé » (plus vertical que la génératrice), il coupe les deux nappes : \textbf{hyperbole}.
\item $\theta = 30^{\circ} < \alpha$ : le plan est « trop couché » (plus horizontal que la génératrice), il n'en coupe qu'une selon une courbe fermée : \textbf{ellipse}.
\end{enumerate}
\end{corrige}

\begin{savaistu}[Apollonius de Perga, le « Grand Géomètre »]
C'est le mathématicien grec \textbf{Apollonius de Perga} (vers 262 -- 190 av. J.-C.) qui, dans son traité \emph{Les Coniques} en huit livres, a baptisé les trois courbes \emph{ellipse}, \emph{parabole} et \emph{hyperbole}. Ces noms viennent de mots grecs signifiant respectivement « manquer », « jeter à côté, comparer » et « dépasser » — une allusion à une construction géométrique de son époque, mais qui s'accorde remarquablement avec la classification moderne par l'excentricité $e$ que nous verrons dans la prochaine section : l'ellipse « manque » ($e < 1$), la parabole est « juste » ($e = 1$), l'hyperbole « dépasse » ($e > 1$). Dix-huit siècles plus tard, Kepler découvrira que les planètes décrivent des ellipses, et Galilée que les projectiles suivent des paraboles : les coniques d'Apollonius gouvernent le cosmos !
\end{savaistu}

\begin{aretenir}[L'essentiel de la section]
\begin{itemize}
\item Un \cle{cône de révolution} est engendré par la rotation d'une droite (génératrice) autour d'un axe, avec un sommet $S$ et un demi-angle au sommet $\alpha$ ; il possède deux nappes.
\item Une \cle{conique} est l'intersection d'un cône de révolution avec un plan.
\item Plan ne passant pas par $S$ : \textbf{cercle} (plan perpendiculaire à l'axe, $\beta=0$), \textbf{ellipse} ($0 < \beta < \alpha$), \textbf{parabole} ($\beta = \alpha$, plan parallèle à une génératrice), \textbf{hyperbole} ($\beta > \alpha$, deux nappes coupées, deux branches).
\item Plan passant par $S$ : coniques \textbf{dégénérées} (point, droite, deux droites sécantes).
\item La parabole est le cas limite entre ellipse et hyperbole : ne pas confondre « incliné » ($\beta < \alpha$) avec « parallèle à une génératrice » ($\beta = \alpha$).
\end{itemize}
\end{aretenir}

Dans la section suivante, nous abandonnerons le cône pour une définition plane et bien plus maniable des coniques : la définition \emph{monofocale} $MF = e \cdot d(M, (D))$, où un unique nombre, l'\cle{excentricité} $e$, permettra de retrouver toute la classification que nous venons de découvrir géométriquement.

\coursec{Définition monofocale et excentricité}

Dans la section précédente, nous avons découvert les coniques comme des \emph{sections planes d'un cône de révolution} : selon l'inclinaison du plan de coupe, on obtient une ellipse, une parabole ou une hyperbole. Cette approche géométrique, héritée des Grecs, est belle mais peu pratique pour calculer. Comment tracer ces courbes avec une règle et un compas ? Comment en écrire des équations ? Toute cette section repose sur une idée simple et puissante : \textbf{une conique peut être fabriquée à partir d'un point (le foyer) et d'une droite (la directrice), en comparant deux distances}. C'est ce qu'on appelle la \cle{définition monofocale}.

\courssub{La définition monofocale}

\textbf{Une expérience pour deviner.}\par
Imagine un poteau téléphonique de MTN planté en un point $F$ d'un champ, et une route rectiligne matérialisée par une droite $D$ qui ne passe pas par le poteau. On cherche les positions $M$ du champ où la distance au poteau vaut \emph{exactement} la distance à la route : $MF = d(M, D)$. En explorant, tu obtiens une courbe qui s'écarte de plus en plus de la route : c'est une \cle{parabole}. Si maintenant on exige que la distance au poteau vaille la \emph{moitié} de la distance à la route ($MF = \frac{1}{2}d(M,D)$), la courbe se referme sur elle-même : c'est une \cle{ellipse}. Et si l'on exige $MF = 2\,d(M,D)$, la courbe se scinde en deux branches qui partent à l'infini : c'est une \cle{hyperbole}. Un seul nombre — le rapport des deux distances — contrôle donc la nature de la courbe. Ce nombre s'appelle l'\cle{excentricité}.

\begin{definition}[Conique de foyer $F$, de directrice $D$ et d'excentricité $e$]
Soit $F$ un point du plan, $D$ une droite \textbf{ne passant pas par} $F$, et $e$ un réel strictement positif. On appelle \cle{conique} de \cle{foyer} $F$, de \cle{directrice} $D$ et d'\cle{excentricité} $e$ l'ensemble des points $M$ du plan tels que :
\[ MF = e \times d(M, D) \]
où $d(M, D)$ est la distance du point $M$ à la droite $D$ (c'est-à-dire $MH$, où $H$ est le projeté orthogonal de $M$ sur $D$).
\begin{itemize}
\item Si $e = 1$, la conique est une \cle{parabole} ;
\item si $0 < e < 1$, la conique est une \cle{ellipse} ;
\item si $e > 1$, la conique est une \cle{hyperbole}.
\end{itemize}
\end{definition}

\begin{attention}[Condition essentielle : $F \notin D$]
La directrice ne doit \textbf{jamais} passer par le foyer. En effet, si $F \in D$, alors pour tout point $M$ de $D$, la condition $MF = e \cdot d(M,D)$ devient $MF = e \times 0 = 0$, c'est-à-dire $M = F$ : la conique \emph{dégénère} (elle se réduit à un point ou à des droites). Dans tout exercice, vérifie d'abord que $F \notin D$ avant de conclure qu'on a affaire à une vraie conique.
\end{attention}

\begin{propriete}[Lien avec les sections du cône (admis)]
Toute conique obtenue comme section plane d'un cône de révolution (autre qu'un cercle) admet une définition monofocale : il existe un point $F$, une droite $D$ et un réel $e > 0$ tels que la conique soit l'ensemble des points $M$ vérifiant $MF = e \cdot d(M,D)$. Réciproquement, toute courbe définie par $MF = e \cdot d(M,D)$ est une conique au sens des sections d'un cône. Ce théorème est \textbf{admis} en Terminale ; il garantit que les deux points de vue (section du cône et foyer-directrice) décrivent exactement les mêmes courbes.
\end{propriete}

\begin{savaistu}[D'où vient le mot « excentricité » ?]
Le mot vient du latin \emph{ex centrum}, « hors du centre ». Plus $e$ est proche de $0$, plus la courbe est « centrée », proche d'un cercle ; plus $e$ grandit, plus la courbe s'excentre et s'ouvre. L'orbite de la Terre autour du Soleil est une ellipse d'excentricité $e \approx 0,0167$ : presque un cercle ! Celle de la comète de Halley vaut $e \approx 0,967$ : une ellipse très allongée. Au-delà de $e = 1$, un corps céleste ne revient jamais : trajectoire parabolique ou hyperbolique.
\end{savaistu}

\courssub{Discussion selon la valeur de l'excentricité}

Reprenons notre champ et faisons varier $e$. Posons $K$ le projeté orthogonal de $F$ sur $D$, et notons $\delta = FK > 0$ la distance du foyer à la directrice. Cherchons, dans chaque cas, les points de la conique situés sur la droite $(FK)$ : un tel point $M$ vérifie $MF = e \cdot MK$ (car son projeté orthogonal sur $D$ est $K$).

\textbf{Cas $e = 1$ : la parabole.} L'équation devient $MF = d(M, D)$. Sur la droite $(FK)$, la condition $MF = MK$ caractérise les points équidistants de $F$ et de $K$ : il n'y en a qu'un, le milieu $S$ du segment $[FK]$, appelé \cle{sommet}. La courbe passe par $S$ et s'ouvre du côté de $F$, en s'écartant indéfiniment : une seule branche, non bornée.

\textbf{Cas $0 < e < 1$ : l'ellipse.} Cherchons les points de la courbe situés sur la droite $(FK)$. Un point $M$ de cette droite, à distance $x$ de $F$ et situé entre $F$ et $K$, est à distance $\delta - x$ de $D$ : la condition $x = e(\delta - x)$ donne $x = \dfrac{e\delta}{1+e}$. De l'autre côté de $F$ (opposé à $D$), $M$ est à distance $\delta + x$ de $D$ : la condition $x = e(\delta + x)$ donne $x = \dfrac{e\delta}{1-e}$, qui est un nombre \emph{fini positif} car $e < 1$. La courbe coupe donc l'axe focal en \textbf{deux} points, tous deux à distance finie de $F$, et l'on peut montrer qu'elle reste prisonnière d'une zone bornée : elle se referme. C'est une courbe fermée, ovale.

\textbf{Cas $e > 1$ : l'hyperbole.} Le premier point existe toujours ($x = \dfrac{e\delta}{1+e}$, entre $F$ et $K$), mais la quantité $\dfrac{e\delta}{1-e}$ est maintenant \emph{négative} : cela signifie qu'il n'existe aucun point de la courbe de l'autre côté de $F$. En revanche, un second point d'intersection apparaît de l'\emph{autre côté de la directrice} : pour un point $M$ situé au-delà de $K$, à distance $x$ de $F$, on a $d(M,D) = x - \delta$ et la condition $x = e(x - \delta)$ donne $x = \dfrac{e\delta}{e-1} > 0$. La courbe possède donc \textbf{deux branches} disjointes, chacune partant à l'infini.

\textbf{Cas limite $e = 0$ (mention).} L'équation deviendrait $MF = 0 \times d(M,D) = 0$, c'est-à-dire $M = F$ : pas de courbe. Mais on peut voir le \cle{cercle} comme une ellipse « d'excentricité nulle » : quand $e$ tend vers $0$, l'ellipse devient de plus en plus ronde et ses deux foyers se confondent en un centre. C'est pourquoi on dit parfois que le cercle est la conique d'excentricité $0$.

\begin{aretenir}[Le réflexe de reconnaissance]
Face à une excentricité : $e = 1 \Rightarrow$ parabole ; $0 < e < 1 \Rightarrow$ ellipse ; $e > 1 \Rightarrow$ hyperbole. L'excentricité mesure « l'aplatissement » ou « l'ouverture » de la conique : c'est le paramètre de \emph{forme}, alors que la distance $\delta = d(F, D)$ est le paramètre de \emph{taille}.
\end{aretenir}

\begin{popfigure}
\begin{tikzpicture}[scale=0.9]
% Directrice
\draw[popink, very thick] (6,-3.2) -- (6,3.2) node[above] {$D$};
% Axe focal
\draw[popmuted, dashed] (-1.5,0) -- (6.6,0);
% Foyer
\fill[poporange] (2,0) circle (2.2pt) node[below left] {$F$};
\fill[popmuted] (6,0) circle (1.5pt) node[below right] {$K$};
% Parallèle d = 3
\draw[popblue, dashed] (3,-3.2) -- (3,3.2);
\draw[popblue] (2,0) circle (1.5);
\fill[popblue] (3,1.118) circle (2pt) node[above right] {$M_1$};
\fill[popblue] (3,-1.118) circle (2pt);
\draw[popblue!60] (3,1.118) -- (6,1.118) node[midway, above, scale=0.8] {$d_1$};
\draw[popblue!60] (2,0) -- (3,1.118) node[midway, left, scale=0.8] {$e\,d_1$};
% Parallèle d = 4
\draw[popgreen, dashed] (2,-3.2) -- (2,3.2);
\fill[popgreen] (2,2) circle (2pt) node[above right] {$M_2$};
\fill[popgreen] (2,-2) circle (2pt);
\draw[popgreen!70] (2,2) -- (6,2) node[midway, above, scale=0.8] {$d_2$};
\draw[popgreen!70] (2,0) -- (2,2) node[midway, right, scale=0.8] {$e\,d_2$};
% Parallèle d = 6
\draw[poppurple, dashed] (0,-3.2) -- (0,3.2);
\fill[poppurple] (0,2.236) circle (2pt) node[above left] {$M_3$};
\fill[poppurple] (0,-2.236) circle (2pt);
\draw[poppurple!70] (0,2.236) -- (6,2.236) node[midway, above, scale=0.8] {$d_3$};
\draw[poppurple!70] (2,0) -- (0,2.236) node[midway, above, scale=0.8] {$e\,d_3$};
\end{tikzpicture}
\legende{Construction point par point (ici $e = \tfrac12$) : pour chaque parallèle à $D$ à distance $d$, on trace le cercle de centre $F$ et de rayon $e\,d$ ; les intersections sont des points de la conique.}
\end{popfigure}

\courssub{Axe focal, sommets et éléments de symétrie}

\begin{definition}[Axe focal]
On appelle \cle{axe focal} de la conique la droite passant par le foyer $F$ et perpendiculaire à la directrice $D$ (c'est la droite $(FK)$, où $K$ est le projeté orthogonal de $F$ sur $D$). Les points d'intersection de la conique avec son axe focal sont appelés les \cle{sommets} de la conique.
\end{definition}

\begin{propriete}[Symétrie par rapport à l'axe focal]
Toute conique de foyer $F$ et de directrice $D$ est symétrique par rapport à son axe focal.
\end{propriete}

\begin{experience}[Démonstration guidée]
On veut montrer que si $M$ appartient à la conique, alors son symétrique $M'$ par rapport à l'axe focal $(FK)$ appartient aussi à la conique.
\begin{enumerate}
\item \textbf{Ce qui est conservé par la symétrie.} La symétrie axiale par rapport à $(FK)$ est une isométrie : elle conserve les distances. Comme $F \in (FK)$, le point $F$ est invariant, donc $M'F = MF$.
\item \textbf{Distance à la directrice.} La droite $D$ est perpendiculaire à l'axe $(FK)$ : elle est donc globalement invariante par la symétrie. Une isométrie conservant les distances, on a $d(M', D) = d(M, D)$.
\item \textbf{Conclusion.} Si $M$ vérifie $MF = e \cdot d(M, D)$, alors :
\[ M'F = MF = e \cdot d(M, D) = e \cdot d(M', D), \]
donc $M'$ appartient à la conique. \textbf{La conique est bien symétrique par rapport à son axe focal.} $\blacksquare$
\end{enumerate}
\end{experience}

\begin{propriete}[Éléments de symétrie de chaque conique]
\begin{itemize}
\item La \textbf{parabole} admet un unique axe de symétrie : son axe focal. Elle possède \textbf{un seul sommet} $S$, milieu du segment $[FK]$.
\item L'\textbf{ellipse} et l'\textbf{hyperbole} admettent, outre l'axe focal, un \textbf{second axe de symétrie} (la médiatrice du segment joignant leurs deux sommets) et un \textbf{centre de symétrie} $O$ (le milieu de ce segment). On les appelle les \cle{coniques à centre}. Chacune possède \textbf{deux sommets} sur l'axe focal, et l'ellipse en possède deux autres sur le second axe.
\end{itemize}
\end{propriete}

\textbf{Pourquoi un centre pour l'ellipse et l'hyperbole ?}\par
Nous avons vu dans la discussion que lorsque $e \neq 1$, la conique coupe l'axe focal en deux points $S$ et $S'$. On peut alors montrer (et nous l'admettrons) qu'il existe un \emph{second} couple foyer-directrice $(F', D')$, symétrique de $(F, D)$ par rapport au milieu $O$ de $[SS']$, qui définit la \emph{même} courbe. La symétrie centrale de centre $O$ échange $F$ et $F'$, $D$ et $D'$, et conserve donc la conique. La parabole, elle, n'a qu'un seul sommet et ne peut pas avoir de centre : c'est une conique « sans centre ».

\courssub{Construction point par point d'une conique}

La définition monofocale fournit immédiatement un procédé de tracé à la règle et au compas : c'est exactement ce procédé qu'on te demandera de mettre en œuvre au Baccalauréat.

\begin{methode}[Construire une conique point par point]
On donne le foyer $F$, la directrice $D$ et l'excentricité $e$.
\begin{enumerate}
\item Tracer l'axe focal (perpendiculaire à $D$ passant par $F$) et placer $K$, projeté orthogonal de $F$ sur $D$.
\item \textbf{Choisir} une distance $d > 0$ et tracer la droite $\Delta_d$, parallèle à $D$, située à la distance $d$ de $D$ (du côté de $F$, et éventuellement de l'autre côté).
\item \textbf{Reporter au compas} la longueur $e \times d$ à partir de $F$ : tracer le cercle $\mathcal{C}(F, e\,d)$.
\item Les points d'intersection de $\Delta_d$ et du cercle (s'ils existent) sont des points de la conique : en effet, pour un tel point $M$, $d(M, D) = d$ et $MF = e\,d$, donc $MF = e \cdot d(M, D)$.
\item Recommencer avec plusieurs valeurs de $d$, puis relier les points obtenus par une courbe régulière, en exploitant la symétrie par rapport à l'axe focal.
\end{enumerate}
\textbf{Remarque :} le cercle et la parallèle ne se coupent que si le rayon $e\,d$ est supérieur ou égal à la distance de $F$ à $\Delta_d$ ; si le cercle est trop petit, il n'y a pas de point pour cette valeur de $d$ — il faut en choisir une autre.
\end{methode}

\begin{exemplebox}[Exemple chiffré : cinq points d'une ellipse]
Soit $F(2\,;\,0)$, $D$ la droite d'équation $x = 8$, et $e = \dfrac{1}{2}$. Pour un point $M(x\,;\,y)$ situé à gauche de $D$, $d(M, D) = 8 - x$, et la condition est $MF = \dfrac{8-x}{2}$. Choisissons plusieurs valeurs de $d = 8 - x$ :
\begin{itemize}
\item $d = 4$ : $x = 4$, $MF = 2$. Or $MF^2 = (4-2)^2 + y^2 = 4$ donne $y = 0$ : point $S_1(4\,;\,0)$, un \textbf{sommet}.
\item $d = 6$ : $x = 2$, $MF = 3$. Alors $0 + y^2 = 9$, soit $y = \pm 3$ : points $(2\,;\,3)$ et $(2\,;\,-3)$.
\item $d = 8$ : $x = 0$, $MF = 4$. Alors $4 + y^2 = 16$, soit $y = \pm 2\sqrt{3} \approx \pm 3,46$ : points $(0\,;\,2\sqrt{3})$ et $(0\,;\,-2\sqrt{3})$.
\item $d = 12$ : $x = -4$, $MF = 6$. Alors $36 + y^2 = 36$, soit $y = 0$ : point $S_2(-4\,;\,0)$, le second sommet.
\end{itemize}
On obtient ainsi cinq points (plus leurs symétriques) d'une ellipse de sommets $S_1(4\,;\,0)$ et $S_2(-4\,;\,0)$, de centre $O(0\,;\,0)$ — le milieu de $[S_1S_2]$.
\end{exemplebox}

\begin{attention}[Pièges fréquents dans la construction]
\begin{itemize}
\item Ne pas confondre $d$ (distance de $M$ à $D$) avec $MF$ : c'est $MF = e \cdot d$, donc le rayon du cercle est $e \times d$, pas $d$.
\item Si $e < 1$, le rayon $e\,d$ est \emph{plus petit} que $d$ : beaucoup d'élèves inversent et tracent un cercle trop grand.
\item La parallèle à $D$ peut être tracée des \emph{deux côtés} de $D$ : pour l'hyperbole, c'est indispensable pour obtenir la deuxième branche.
\item N'oublie pas d'exploiter la symétrie par rapport à l'axe focal : chaque point construit au-dessus de l'axe en donne un gratuitement en dessous.
\end{itemize}
\end{attention}

\courssub{Une même famille, trois visages}

La figure ci-dessous montre trois coniques définies par le \emph{même} foyer $F$ (l'origine) et des directrices de même direction, seule l'excentricité change : $e = 0,6$ pour l'ellipse, $e = 1$ pour la parabole, $e = 1,5$ pour l'hyperbole. Observe comment la courbe « s'ouvre » quand $e$ augmente.

\begin{popfigure}
\begin{tikzpicture}[scale=1.35]
\begin{scope}
\draw[popink, thick] (-1.667,-1.6) -- (-1.667,1.6) node[above, scale=0.75] {$D$};
\draw[popmuted, dashed] (-2.1,0) -- (2.8,0);
\fill[poporange] (0,0) circle (1.6pt) node[below, scale=0.8] {$F$};
\draw[popblue, very thick, domain=0:360, samples=120, smooth] plot ({0.9375+1.5625*cos(\x)}, {1.25*sin(\x)});
\node[popblue, scale=0.85] at (0.4,-1.9) {Ellipse : $e = 0,6$};
\end{scope}
\begin{scope}[xshift=5.4cm]
\draw[popink, thick] (-1,-1.6) -- (-1,1.6) node[above, scale=0.75] {$D$};
\draw[popmuted, dashed] (-1.4,0) -- (2.1,0);
\fill[poporange] (0,0) circle (1.6pt) node[below, scale=0.8] {$F$};
\draw[popgreen, very thick, domain=-0.49:1.6, samples=60, smooth] plot (\x, {sqrt(max(0,2*\x+1))});
\draw[popgreen, very thick, domain=-0.49:1.6, samples=60, smooth] plot (\x, {-sqrt(max(0,2*\x+1))});
\node[popgreen!70!black, scale=0.85] at (0.3,-1.9) {Parabole : $e = 1$};
\end{scope}
\begin{scope}[xshift=10.2cm]
\draw[popink, thick] (-0.667,-1.6) -- (-0.667,1.6) node[above, scale=0.75] {$D$};
\draw[popmuted, dashed] (-2.6,0) -- (0.7,0);
\fill[poporange] (0,0) circle (1.6pt) node[below, scale=0.8] {$F$};
\draw[poppurple, very thick, domain=-1.55:1.55, samples=60, smooth] plot ({-1.2+0.8*cosh(\x)}, {0.8944*sinh(\x)});
\draw[poppurple, very thick, domain=-1.55:1.55, samples=60, smooth] plot ({-1.2-0.8*cosh(\x)}, {0.8944*sinh(\x)});
\node[poppurple, scale=0.85] at (-1,-1.9) {Hyperbole : $e = 1,5$};
\end{scope}
\end{tikzpicture}
\legende{Trois coniques de même foyer $F$ et de directrices verticales : seule l'excentricité change. L'ellipse est fermée, la parabole est ouverte à une branche, l'hyperbole a deux branches.}
\end{popfigure}

\begin{exoresolu}[Identifier et construire]
\textbf{Énoncé.} Dans un repère orthonormé, on donne $F(0\,;\,0)$ et la droite $D$ d'équation $x = 4$. On considère l'ensemble $\mathcal{C}$ des points $M(x\,;\,y)$ tels que $MF = \dfrac{3}{2}\,d(M, D)$. (1) Quelle est la nature de $\mathcal{C}$ ? (2) Déterminer les sommets de $\mathcal{C}$ sur l'axe focal. (3) Donner les points de $\mathcal{C}$ d'abscisse $-2$.
\tcblower
\textbf{Solution.} (1) L'excentricité est $e = \dfrac{3}{2} > 1$ et $F \notin D$ (car l'abscisse de $F$ est $0 \neq 4$) : $\mathcal{C}$ est une \textbf{hyperbole} de foyer $F$ et de directrice $D$.

(2) L'axe focal est la perpendiculaire à $D$ passant par $F$ : c'est l'axe des abscisses. Un point $M(x\,;\,0)$ de cet axe vérifie $MF = |x|$ et $d(M, D) = |x - 4|$. La condition s'écrit $|x| = \dfrac{3}{2}|x-4|$, soit $2|x| = 3|x-4|$. En élevant au carré (les deux membres sont positifs, il y a équivalence) : $4x^2 = 9(x-4)^2$, d'où $2x = 3(x-4)$ ou $2x = -3(x-4)$.
\begin{align*}
2x &= 3(x-4) \quad\Rightarrow\quad 2x = 3x - 12 \quad\Rightarrow\quad x = 12 ; \\
2x &= -3(x-4) \quad\Rightarrow\quad 5x = 12 \quad\Rightarrow\quad x = \tfrac{12}{5} = 2,4.
\end{align*}
Vérification : pour $x = 12$, $|12| = 12$ et $\dfrac32|12-4| = \dfrac32 \times 8 = 12$ : convient. Pour $x = \dfrac{12}{5}$, $\dfrac{12}{5}$ et $\dfrac32\left|\dfrac{12}{5} - 4\right| = \dfrac32 \times \dfrac85 = \dfrac{12}{5}$ : convient. Les sommets sont donc $S_1\left(\dfrac{12}{5}\,;\,0\right)$ (situé entre $F$ et $D$) et $S_2(12\,;\,0)$ (situé de l'autre côté de la directrice), ce qui confirme l'existence de deux branches.

(3) Pour $x = -2$ : $MF^2 = (-2)^2 + y^2 = 4 + y^2$ et $d(M,D) = |-2 - 4| = 6$, donc la condition donne $MF = \dfrac32 \times 6 = 9$, soit $4 + y^2 = 81$, d'où $y^2 = 77$ et $y = \pm\sqrt{77} \approx \pm 8,77$. Les points $(-2\,;\,\sqrt{77})$ et $(-2\,;\,-\sqrt{77})$ appartiennent à $\mathcal{C}$.
\end{exoresolu}

\begin{exercice}[À toi de jouer]
On donne $F(1\,;\,2)$, $D$ la droite d'équation $y = -1$, et $e = \dfrac{1}{3}$.
\begin{enumerate}
\item Quelle est la nature de la conique de foyer $F$, de directrice $D$ et d'excentricité $e$ ?
\item Déterminer les coordonnées de ses deux sommets sur l'axe focal.
\item Déterminer quatre autres points de la conique en choisissant deux valeurs de $d$ bien adaptées.
\end{enumerate}
\end{exercice}

\begin{corrige}[Exercice : conique d'excentricité 1/3]
\begin{enumerate}
\item $e = \dfrac{1}{3} < 1$ et $F \notin D$ (car l'ordonnée de $F$ est $2 \neq -1$) : c'est une \textbf{ellipse}.
\item L'axe focal est la perpendiculaire à $D$ passant par $F$ : c'est la droite d'équation $x = 1$. Pour $M(1\,;\,y)$ : $MF = |y - 2|$ et $d(M, D) = |y + 1|$. La condition $|y-2| = \dfrac{1}{3}|y+1|$ s'écrit $3|y-2| = |y+1|$, soit $3(y-2) = y+1$ ou $3(y-2) = -(y+1)$.
\begin{itemize}
\item $3y - 6 = y + 1$ donne $2y = 7$, soit $y = \dfrac{7}{2}$ ;
\item $3y - 6 = -y - 1$ donne $4y = 5$, soit $y = \dfrac{5}{4}$.
\end{itemize}
Les deux valeurs conviennent (vérifier : $\left|\dfrac72 - 2\right| = \dfrac32 = \dfrac13\left|\dfrac72 + 1\right|$ et $\left|\dfrac54 - 2\right| = \dfrac34 = \dfrac13\left|\dfrac54 + 1\right|$). Sommets : $S_1\left(1\,;\,\dfrac{7}{2}\right)$ et $S_2\left(1\,;\,\dfrac{5}{4}\right)$.
\item \textbf{Choix $d = 3$} (du côté de $F$) : la parallèle à $D$ à distance $3$ est la droite $y = 2$, et le rayon du cercle est $e\,d = 1$. Le cercle de centre $F(1\,;\,2)$ et de rayon $1$ coupe la droite $y = 2$ aux points vérifiant $(x-1)^2 = 1$, soit $x = 0$ ou $x = 2$ : points $(0\,;\,2)$ et $(2\,;\,2)$.

\textbf{Choix $d = 4$} (du côté de $F$) : la parallèle est la droite $y = 3$, et le rayon est $e\,d = \dfrac{4}{3}$. Les points cherchés vérifient $(x-1)^2 + (3-2)^2 = \dfrac{16}{9}$, soit $(x-1)^2 = \dfrac{16}{9} - 1 = \dfrac{7}{9}$, d'où $x = 1 \pm \dfrac{\sqrt{7}}{3} \approx 1 \pm 0,88$ : points $\left(1 + \dfrac{\sqrt{7}}{3}\,;\,3\right)$ et $\left(1 - \dfrac{\sqrt{7}}{3}\,;\,3\right)$.

Avec les deux sommets, ces quatre points permettent d'esquisser l'ellipse, de centre $\left(1\,;\,\dfrac{19}{8}\right)$, milieu de $[S_1S_2]$.
\end{enumerate}
\end{corrige}

\begin{aretenir}[L'essentiel de la section]
\begin{itemize}
\item Une conique est l'ensemble des points $M$ tels que $MF = e \cdot d(M, D)$, avec $F \notin D$ et $e > 0$.
\item $e = 1$ : parabole ; $0 < e < 1$ : ellipse ; $e > 1$ : hyperbole ; $e = 0$ correspond au cas limite du cercle.
\item L'axe focal (perpendiculaire à $D$ passant par $F$) est toujours axe de symétrie ; ellipse et hyperbole ont en plus un centre de symétrie : ce sont des coniques à centre.
\item Construction point par point : choisir $d$, tracer la parallèle à $D$ à distance $d$, puis le cercle de centre $F$ et de rayon $e\,d$ ; les intersections sont des points de la conique.
\end{itemize}
Dans la section suivante, nous traduirons cette définition géométrique en \textbf{équations cartésiennes réduites}, ce qui nous permettra de calculer foyers, sommets et directrices avec une efficacité redoutable.
\end{aretenir}

\coursec{Équations réduites et éléments caractéristiques}

Dans la section précédente, nous avons défini une conique par sa \cle{définition monofocale} : étant donné un point $F$ (le \cle{foyer}), une droite $\mathcal{D}$ (la \cle{directrice}) ne passant pas par $F$, et un réel $e > 0$ (l'\cle{excentricité}), la conique est l'ensemble des points $M$ du plan vérifiant
\[ MF = e \cdot d(M, \mathcal{D}). \]
Cette définition est belle, mais peu pratique pour calculer. Toute la puissance de la géométrie analytique consiste à \emph{choisir un bon repère} pour transformer cette condition géométrique en une \cle{équation} simple. C'est exactement ce que fait un ingénieur de CAMTEL lorsqu'il modélise la courbe d'une antenne parabolique : il place son repère « au bon endroit » pour obtenir l'équation la plus simple possible. Cette équation, dite \emph{réduite}, contient alors à elle seule toute l'information sur la conique : foyers, directrices, sommets, excentricité. Apprenons à la lire comme on lit une carte d'identité.

\courssub{Le choix du repère adapté : le repère propre}

L'idée directrice est la suivante : les coniques possèdent des \cle{éléments de symétrie} (vus en section 2), et un repère dont les axes coïncident avec ces axes de symétrie simplifie considérablement les calculs, car les termes en $xy$ et la plupart des termes du premier degré disparaissent.

\begin{definition}[Repère propre et équation réduite]
On appelle \cle{repère propre} d'une conique un repère orthonormé choisi de la manière suivante :
\begin{itemize}
\item pour une \cle{ellipse} ou une \cle{hyperbole} (coniques \emph{à centre}) : l'\textbf{origine est le centre de symétrie} de la conique, et l'axe des abscisses est l'\cle{axe focal} (la droite passant par les deux foyers) ;
\item pour une \cle{parabole} : l'\textbf{origine est le sommet} $S$ de la parabole, et l'axe des abscisses est l'axe de symétrie, orienté du sommet vers le foyer.
\end{itemize}
L'équation d'une conique exprimée dans son repère propre s'appelle son \cle{équation réduite}.
\end{definition}

\begin{attention}[Le repère propre n'est pas toujours le repère de l'énoncé]
Dans les exercices du Baccalauréat, l'équation est parfois donnée dans un repère quelconque : on reconnaît alors des termes en $x$ et $y$ du premier degré. La mise sous forme réduite (par \emph{formes canoniques}) sera l'objet de la section 5. Ici, nous supposons le repère déjà bien choisi.
\end{attention}

\courssub{Équation réduite de l'ellipse}

\courssub{De la définition bifocale à l'équation}

Commençons par une expérience simple que tu peux réaliser chez toi : plante deux punaises $F$ et $F'$ sur une planche, attache une ficelle de longueur fixe (supérieure à $FF'$) entre les deux punaises, tends la ficelle avec la pointe d'un crayon et fais tourner le crayon. La courbe tracée est une \cle{ellipse} : c'est l'ensemble des points $M$ tels que $MF + MF'$ est constante. C'est la \cle{définition bifocale} de l'ellipse. Montrons qu'elle coïncide avec notre définition monofocale et aboutit à l'équation réduite annoncée.

\begin{experience}[Démonstration guidée : de $MF + MF' = 2a$ à l'équation réduite]
Soient $F(c\,;\,0)$ et $F'(-c\,;\,0)$ deux points fixes (avec $c > 0$), et $a$ un réel tel que $a > c$. On considère l'ensemble $(\mathcal{E})$ des points $M(x\,;\,y)$ tels que $MF + MF' = 2a$.

\textbf{Étape 1 : écrire les distances.}
\[ MF = \sqrt{(x-c)^2 + y^2} \qquad \text{et} \qquad MF' = \sqrt{(x+c)^2 + y^2}. \]

\textbf{Étape 2 : calculer la différence des carrés.}
\begin{align*}
MF'^2 - MF^2 &= \left[(x+c)^2 + y^2\right] - \left[(x-c)^2 + y^2\right] = 4cx.
\end{align*}
Or $MF'^2 - MF^2 = (MF' - MF)(MF' + MF) = (MF' - MF) \times 2a$. On en déduit :
\[ MF' - MF = \frac{4cx}{2a} = \frac{2cx}{a}. \]

\textbf{Étape 3 : isoler $MF$.} En combinant $MF' + MF = 2a$ et $MF' - MF = \dfrac{2cx}{a}$, on obtient par soustraction :
\[ MF = a - \frac{cx}{a}. \]

\textbf{Étape 4 : élever au carré.}
\[ (x-c)^2 + y^2 = \left(a - \frac{cx}{a}\right)^2 = a^2 - 2cx + \frac{c^2 x^2}{a^2}. \]
En développant le membre de gauche : $x^2 - 2cx + c^2 + y^2 = a^2 - 2cx + \dfrac{c^2 x^2}{a^2}$, soit après simplification :
\[ x^2\left(1 - \frac{c^2}{a^2}\right) + y^2 = a^2 - c^2
\quad\Longleftrightarrow\quad
\frac{x^2 (a^2 - c^2)}{a^2} + y^2 = a^2 - c^2. \]

\textbf{Étape 5 : poser $b^2 = a^2 - c^2$.} Comme $a > c$, le réel $b = \sqrt{a^2 - c^2}$ est bien défini. En divisant par $a^2 - c^2 = b^2$ :
\[ \boxed{\frac{x^2}{a^2} + \frac{y^2}{b^2} = 1.} \]
Réciproquement, on vérifie que tout point satisfaisant cette équation vérifie $MF + MF' = 2a$ (les étapes se remontent, les quantités manipulées étant positives). \hfill$\blacksquare$
\end{experience}

\begin{propriete}[Équation réduite et éléments caractéristiques de l'ellipse]
Dans son repère propre, l'ellipse de \cle{demi-grand axe} $a$ et de \cle{demi-petit axe} $b$ (avec $a > b > 0$) a pour équation réduite :
\[ \frac{x^2}{a^2} + \frac{y^2}{b^2} = 1. \]
Ses éléments caractéristiques sont :
\begin{itemize}
\item \cle{centre} : $O(0\,;\,0)$, centre de symétrie ;
\item \cle{sommets} : $A(a\,;\,0)$, $A'(-a\,;\,0)$ (sur le grand axe) et $B(0\,;\,b)$, $B'(0\,;\,-b)$ (sur le petit axe) ;
\item \cle{distance focale} : $c$ défini par $c^2 = a^2 - b^2$ ;
\item \cle{foyers} : $F(c\,;\,0)$ et $F'(-c\,;\,0)$ ;
\item \cle{excentricité} : $e = \dfrac{c}{a}$, avec $0 < e < 1$ ;
\item \cle{directrices} : droites d'équations $x = \dfrac{a^2}{c}$ et $x = -\dfrac{a^2}{c}$ ;
\item \cle{propriété bifocale} : pour tout point $M$ de l'ellipse, $MF + MF' = 2a$.
\end{itemize}
\end{propriete}

\begin{aretenir}[Retenir les relations fondamentales]
\begin{center}
\begin{tabularx}{\linewidth}{|l|X|X|}
\hline
\rowcolor{popblueL} & \textbf{Ellipse} & \textbf{Hyperbole} \\
\hline
Équation & $\dfrac{x^2}{a^2} + \dfrac{y^2}{b^2} = 1$ & $\dfrac{x^2}{a^2} - \dfrac{y^2}{b^2} = 1$ \\
\hline
Relation & $c^2 = a^2 - b^2$ & $c^2 = a^2 + b^2$ \\
\hline
Excentricité & $e = \dfrac{c}{a} < 1$ & $e = \dfrac{c}{a} > 1$ \\
\hline
\end{tabularx}
\end{center}
\end{aretenir}

\begin{attention}[L'erreur classique sur la relation entre $a$, $b$ et $c$]
Pour l'\textbf{ellipse}, $c^2 = a^2 - b^2$ (soustraction, car $c < a$ : le foyer est \emph{à l'intérieur} de l'ellipse). Pour l'\textbf{hyperbole}, $c^2 = a^2 + b^2$ (addition, car $c > a$ : le foyer est \emph{au-delà} du sommet). Inverser ces deux relations est la faute la plus fréquente au Baccalauréat. Moyen mnémotechnique : dans l'ellipse, $a$ est le \emph{plus grand} des trois nombres ; dans l'hyperbole, c'est $c$.
\end{attention}

\begin{exemplebox}[Lecture complète d'une équation d'ellipse]
Considérons la courbe d'équation $\dfrac{x^2}{25} + \dfrac{y^2}{9} = 1$.

On lit : $a^2 = 25$ et $b^2 = 9$, donc $a = 5$ et $b = 3$. Comme $a > b$ et que $a^2$ est sous $x^2$, le grand axe est l'axe des abscisses. Alors :
\[ c^2 = a^2 - b^2 = 25 - 9 = 16 \quad\Longrightarrow\quad c = 4, \qquad e = \frac{c}{a} = \frac{4}{5}. \]
Foyers : $F(4\,;\,0)$ et $F'(-4\,;\,0)$. Sommets : $(\pm 5\,;\,0)$ et $(0\,;\,\pm 3)$. Directrices : $x = \pm \dfrac{a^2}{c} = \pm \dfrac{25}{4}$. Vérification bifocale : pour tout point $M$ de la courbe, $MF + MF' = 2a = 10$.
\end{exemplebox}

\begin{popfigure}
\begin{center}
\begin{tikzpicture}[scale=0.85,>=Stealth]
% directrices
\draw[popmuted,dashed,thick] (6.25,-2.6) -- (6.25,2.6);
\draw[popmuted,dashed,thick] (-6.25,-2.6) -- (-6.25,2.6);
\node[popmuted,below] at (6.25,-2.6) {\small $x=\frac{a^2}{c}$};
\node[popmuted,below] at (-6.25,-2.6) {\small $x=-\frac{a^2}{c}$};
% axes
\draw[->,popline] (-6.9,0) -- (6.9,0) node[below left] {$x$};
\draw[->,popline] (0,-2.8) -- (0,2.8) node[below left] {$y$};
% ellipse
\draw[popblue,very thick] (0,0) ellipse (5 and 3);
% foyers
\fill[popink] (4,0) circle (2.2pt) node[below right=1pt] {\small $F(c;0)$};
\fill[popink] (-4,0) circle (2.2pt) node[below left=1pt] {\small $F'(-c;0)$};
% sommets
\fill[poporange] (5,0) circle (2.2pt) node[above right=0pt] {\small $A(a;0)$};
\fill[poporange] (-5,0) circle (2.2pt) node[above left=0pt] {\small $A'$};
\fill[poporange] (0,3) circle (2.2pt) node[above right=0pt] {\small $B(0;b)$};
\fill[poporange] (0,-3) circle (2.2pt) node[below right=0pt] {\small $B'$};
% a, b, c
\draw[<->,popgreen,thick] (0,0.25) -- (5,0.25) node[midway,above] {\small $a$};
\draw[<->,popgreen,thick] (0.25,0) -- (0.25,3) node[midway,right] {\small $b$};
\draw[<->,poppurple,thick] (0,-0.35) -- (4,-0.35) node[midway,below] {\small $c$};
\node[popdark] at (0,-3.3) {\small Centre $O$};
\end{tikzpicture}
\end{center}
\legende{L'ellipse $\frac{x^2}{a^2}+\frac{y^2}{b^2}=1$ : sommets, foyers, directrices et paramètres $a$, $b$, $c$ (ici $a=5$, $b=3$, $c=4$).}
\end{popfigure}

\courssub{Équation réduite de l'hyperbole}

Pour l'hyperbole, la définition bifocale fait intervenir une \emph{différence} de distances au lieu d'une somme : c'est l'ensemble des points $M$ tels que $|MF - MF'| = 2a$, avec cette fois $2a < FF' = 2c$, donc $c > a$. Un calcul parfaitement analogue à celui de l'ellipse (que nous admettons) conduit à l'équation réduite.

\begin{propriete}[Équation réduite et éléments caractéristiques de l'hyperbole]
Dans son repère propre, l'hyperbole a pour équation réduite :
\[ \frac{x^2}{a^2} - \frac{y^2}{b^2} = 1 \qquad (a > 0,\ b > 0). \]
Ses éléments caractéristiques sont :
\begin{itemize}
\item \cle{centre} : $O(0\,;\,0)$ ;
\item \cle{sommets} : $A(a\,;\,0)$ et $A'(-a\,;\,0)$ (l'axe des ordonnées ne rencontre pas la courbe) ;
\item \cle{distance focale} : $c$ défini par $c^2 = a^2 + b^2$ ;
\item \cle{foyers} : $F(c\,;\,0)$ et $F'(-c\,;\,0)$ ;
\item \cle{excentricité} : $e = \dfrac{c}{a} > 1$ ;
\item \cle{directrices} : $x = \pm \dfrac{a^2}{c}$ ;
\item \cle{asymptotes} : droites d'équations $y = \dfrac{b}{a}\,x$ et $y = -\dfrac{b}{a}\,x$ ;
\item \cle{propriété bifocale} : pour tout point $M$ de l'hyperbole, $|MF - MF'| = 2a$.
\end{itemize}
\end{propriete}

D'où viennent les asymptotes ? Écrivons l'équation sous la forme $y^2 = b^2\left(\dfrac{x^2}{a^2} - 1\right)$, soit $y = \pm \dfrac{b}{a}\sqrt{x^2 - a^2}$. Quand $x$ devient très grand, $\sqrt{x^2 - a^2}$ se comporte comme $|x|$, donc $y \approx \pm \dfrac{b}{a} x$ : la courbe se rapproche indéfiniment de ces deux droites sans jamais les toucher. Le \cle{rectangle fondamental}, de sommets $(\pm a\,;\,\pm b)$, a précisément ces asymptotes pour diagonales : c'est l'outil de tracé par excellence.

\begin{popfigure}
\begin{center}
\begin{tikzpicture}[scale=0.8,>=Stealth]
\draw[->,popline] (-6.2,0) -- (6.2,0) node[below left] {$x$};
\draw[->,popline] (0,-3.6) -- (0,3.6) node[below left] {$y$};
% rectangle fondamental a=2, b=1.5
\draw[popmuted,dashed] (-2,-1.5) rectangle (2,1.5);
% asymptotes
\draw[poporange,dashed,thick,domain=-5.6:5.6,samples=2] plot (\x,{0.75*\x});
\draw[poporange,dashed,thick,domain=-5.6:5.6,samples=2] plot (\x,{-0.75*\x});
\node[poporange] at (4.6,3.9) {\small $y=\frac{b}{a}x$};
\node[poporange] at (4.6,-3.9) {\small $y=-\frac{b}{a}x$};
% hyperbole x^2/4 - y^2/2.25 = 1
\draw[popblue,very thick,domain=-1.55:1.55,samples=80] plot ({2*cosh(\x)},{1.5*sinh(\x)});
\draw[popblue,very thick,domain=-1.55:1.55,samples=80] plot ({-2*cosh(\x)},{1.5*sinh(\x)});
% sommets et foyers (c = sqrt(4+2.25)=2.5)
\fill[poporange] (2,0) circle (2.2pt) node[below right=0pt] {\small $A(a;0)$};
\fill[poporange] (-2,0) circle (2.2pt) node[below left=0pt] {\small $A'$};
\fill[popink] (2.5,0) circle (2.2pt) node[below=2pt] {\small $F$};
\fill[popink] (-2.5,0) circle (2.2pt) node[below=2pt] {\small $F'$};
% directrices x = ±a²/c = ±1.6
\draw[popmuted,dashed,thick] (1.6,-3.3) -- (1.6,3.3);
\draw[popmuted,dashed,thick] (-1.6,-3.3) -- (-1.6,3.3);
\node[popmuted,rotate=90,above] at (1.6,2.6) {\small $x=\frac{a^2}{c}$};
\node[popmuted,rotate=90,above] at (-1.6,2.6) {\small $x=-\frac{a^2}{c}$};
\node[popmuted] at (0,-1.85) {\small rectangle fondamental};
\end{tikzpicture}
\end{center}
\legende{L'hyperbole $\frac{x^2}{a^2}-\frac{y^2}{b^2}=1$ (ici $a=2$, $b=1.5$, $c=2.5$), ses asymptotes, son rectangle fondamental, ses foyers et ses directrices.}
\end{popfigure}

\begin{exemplebox}[Lecture complète d'une équation d'hyperbole]
Soit la courbe d'équation $\dfrac{x^2}{9} - \dfrac{y^2}{16} = 1$.

On lit $a^2 = 9$ et $b^2 = 16$, donc $a = 3$, $b = 4$. Alors :
\[ c^2 = a^2 + b^2 = 9 + 16 = 25 \quad\Longrightarrow\quad c = 5, \qquad e = \frac{c}{a} = \frac{5}{3} > 1. \]
Foyers : $F(5\,;\,0)$, $F'(-5\,;\,0)$. Sommets : $(\pm 3\,;\,0)$. Directrices : $x = \pm \dfrac{9}{5}$. Asymptotes : $y = \pm \dfrac{4}{3}x$. Propriété bifocale : $|MF - MF'| = 2a = 6$.
\end{exemplebox}

\courssub{Équation réduite de la parabole}

La parabole n'a pas de centre de symétrie : le repère propre a pour origine le \cle{sommet} $S$, milieu du segment joignant le foyer à sa projection sur la directrice. Notons $p$ la distance du foyer à la directrice, appelée \cle{paramètre} de la parabole.

\begin{propriete}[Équation réduite de la parabole]
Dans son repère propre (origine au sommet, axe des abscisses = axe focal orienté vers le foyer), la parabole de \cle{paramètre} $p > 0$ a pour équation réduite :
\[ y^2 = 2px. \]
Ses éléments caractéristiques sont :
\begin{itemize}
\item \cle{sommet} : $O(0\,;\,0)$ ;
\item \cle{foyer} : $F\left(\dfrac{p}{2}\,;\,0\right)$ ;
\item \cle{directrice} : droite d'équation $x = -\dfrac{p}{2}$ ;
\item \cle{excentricité} : $e = 1$.
\end{itemize}
\end{propriete}

\begin{experience}[Vérification rapide]
Vérifions que l'équation traduit bien la définition monofocale $MF = d(M, \mathcal{D})$. Pour $M(x\,;\,y)$ :
\[ MF^2 = \left(x - \frac{p}{2}\right)^2 + y^2 \qquad \text{et} \qquad d(M,\mathcal{D})^2 = \left(x + \frac{p}{2}\right)^2. \]
En remplaçant $y^2$ par $2px$ :
\[ MF^2 = x^2 - px + \frac{p^2}{4} + 2px = x^2 + px + \frac{p^2}{4} = \left(x + \frac{p}{2}\right)^2 = d(M,\mathcal{D})^2. \]
Les deux distances étant positives, $MF = d(M, \mathcal{D})$ : la définition monofocale est bien satisfaite avec $e = 1$. \hfill$\blacksquare$
\end{experience}

\begin{popfigure}
\begin{center}
\begin{tikzpicture}[scale=0.9,>=Stealth]
\draw[->,popline] (-3.2,0) -- (5.6,0) node[below left] {$x$};
\draw[->,popline] (0,-3.4) -- (0,3.4) node[below left] {$y$};
% parabole y^2 = 4x (p=2)
\draw[popblue,very thick,domain=-3:3,samples=60] plot ({\x*\x/4},\x);
% directrice x = -1
\draw[popmuted,dashed,thick] (-1,-3.2) -- (-1,3.2);
\node[popmuted,rotate=90,above] at (-1,2.4) {\small $\mathcal{D} : x=-\frac{p}{2}$};
% foyer
\fill[popink] (1,0) circle (2.2pt) node[below right=0pt] {\small $F\left(\frac{p}{2};0\right)$};
% sommet
\fill[poporange] (0,0) circle (2.2pt) node[below left=1pt] {\small $S$};
% paramètre p
\draw[<->,poppurple,thick] (-1,2.9) -- (1,2.9) node[midway,above] {\small $p$};
% point M et égalité des distances
\fill[popgreen] (2.25,3) circle (2pt) node[right] {\small $M$};
\draw[popgreen,dashed] (2.25,3) -- (-1,3) node[midway,above] {\small $d(M,\mathcal{D})$};
\draw[popgreen,dashed] (2.25,3) -- (1,0) node[midway,right=1pt] {\small $MF$};
\end{tikzpicture}
\end{center}
\legende{La parabole $y^2 = 2px$ (ici $p = 2$, donc $y^2 = 4x$) : sommet, foyer, directrice, et l'égalité $MF = d(M,\mathcal{D})$.}
\end{popfigure}

\begin{attention}[Bien lire le paramètre dans $y^2 = 2px$]
Dans l'équation $y^2 = 2px$, le coefficient de $x$ est $2p$, \emph{pas} $p$. Ainsi, pour $y^2 = 8x$, on a $2p = 8$ donc $p = 4$, le foyer est $F(2\,;\,0)$ et la directrice $x = -2$. De même, si l'équation est $x^2 = 2py$, la parabole est \emph{verticale} : foyer $\left(0\,;\,\frac{p}{2}\right)$, directrice $y = -\frac{p}{2}$.
\end{attention}

\courssub{Lecture inverse : de l'équation réduite aux éléments caractéristiques}

Savoir « décoder » une équation réduite est une compétence exigible au Baccalauréat. Voici la démarche complète, puis le tableau de synthèse.

\begin{methode}[Passer de l'équation réduite aux éléments caractéristiques]
\begin{enumerate}
\item \textbf{Identifier le type} : un seul carré (par exemple $y^2 = 2px$) $\Rightarrow$ parabole ; deux carrés de même signe $\Rightarrow$ ellipse ; deux carrés de signes contraires $\Rightarrow$ hyperbole.
\item \textbf{Lire $a^2$ et $b^2$} : ce sont les dénominateurs. Pour l'ellipse, $a^2$ est le \emph{plus grand} des deux ; pour l'hyperbole, $a^2$ est le dénominateur du terme \emph{positif}.
\item \textbf{Calculer $c$} : $c = \sqrt{a^2 - b^2}$ (ellipse) ou $c = \sqrt{a^2 + b^2}$ (hyperbole). Pour la parabole, lire $2p$ et en tirer $p$.
\item \textbf{En déduire} : excentricité $e = \dfrac{c}{a}$, foyers $(\pm c\,;\,0)$, directrices $x = \pm \dfrac{a^2}{c}$, sommets, et asymptotes $y = \pm \dfrac{b}{a}x$ pour l'hyperbole.
\item \textbf{Conclure par une phrase} précisant la nature de la conique et son axe focal.
\end{enumerate}
\end{methode}

\begin{aretenir}[Tableau récapitulatif des trois coniques]
\begin{center}
\begin{tabularx}{\linewidth}{|l|X|X|X|}
\hline
\rowcolor{popblueL} & \textbf{Ellipse} & \textbf{Hyperbole} & \textbf{Parabole} \\
\hline
Équation réduite & $\dfrac{x^2}{a^2}+\dfrac{y^2}{b^2}=1$ & $\dfrac{x^2}{a^2}-\dfrac{y^2}{b^2}=1$ & $y^2 = 2px$ \\
\hline
Excentricité & $0 < e < 1$ & $e > 1$ & $e = 1$ \\
\hline
Relation & $c^2 = a^2 - b^2$ & $c^2 = a^2 + b^2$ & --- \\
\hline
Foyers & $(\pm c\,;\,0)$ & $(\pm c\,;\,0)$ & $\left(\frac{p}{2}\,;\,0\right)$ \\
\hline
Directrices & $x = \pm\dfrac{a^2}{c}$ & $x = \pm\dfrac{a^2}{c}$ & $x = -\dfrac{p}{2}$ \\
\hline
Sommets & $(\pm a\,;\,0)$, $(0\,;\,\pm b)$ & $(\pm a\,;\,0)$ & $(0\,;\,0)$ \\
\hline
Asymptotes & aucune & $y = \pm\dfrac{b}{a}x$ & aucune \\
\hline
Propriété bifocale & $MF+MF' = 2a$ & $|MF - MF'| = 2a$ & --- \\
\hline
\end{tabularx}
\end{center}
\end{aretenir}

\begin{attention}[Ellipse : où est l'axe focal ?]
Pour l'ellipse $\dfrac{x^2}{a^2} + \dfrac{y^2}{b^2} = 1$, la convention est $a > b$ : le \emph{plus grand dénominateur} indique l'axe focal. Ainsi :
\begin{itemize}
\item $\dfrac{x^2}{25} + \dfrac{y^2}{9} = 1$ : $a^2 = 25$ sous $x^2$ $\Rightarrow$ axe focal \textbf{horizontal}, foyers $(\pm 4\,;\,0)$ ;
\item $\dfrac{x^2}{9} + \dfrac{y^2}{25} = 1$ : le plus grand dénominateur est sous $y^2$ $\Rightarrow$ axe focal \textbf{vertical}, $a = 5$, $b = 3$, $c = 4$, foyers $(0\,;\,\pm 4)$, directrices $y = \pm \dfrac{25}{4}$.
\end{itemize}
Oublier cette distinction fait échouer toute la question !
\end{attention}

\begin{exoresolu}[Détermination complète des éléments caractéristiques]
On considère la courbe $(\mathcal{C})$ d'équation $\dfrac{x^2}{16} - \dfrac{y^2}{9} = 1$ dans un repère orthonormé.
\begin{enumerate}
\item Déterminer la nature de $(\mathcal{C})$ et préciser son axe focal.
\item Déterminer ses sommets, ses foyers, son excentricité, ses directrices et ses asymptotes.
\item Soit $M$ un point de $(\mathcal{C})$. Que vaut $|MF - MF'|$ ?
\end{enumerate}
\tcblower
\textbf{1. Nature.} L'équation comporte deux carrés de signes contraires séparés par un signe moins : c'est une \cle{hyperbole}. Le terme positif porte sur $x^2$, donc l'axe focal est l'axe des abscisses.

\textbf{2. Éléments caractéristiques.} On lit $a^2 = 16$ et $b^2 = 9$, donc $a = 4$ et $b = 3$. Pour une hyperbole :
\[ c^2 = a^2 + b^2 = 16 + 9 = 25 \quad\Longrightarrow\quad c = 5. \]
\begin{itemize}
\item Sommets : $A(4\,;\,0)$ et $A'(-4\,;\,0)$ ;
\item Foyers : $F(5\,;\,0)$ et $F'(-5\,;\,0)$ ;
\item Excentricité : $e = \dfrac{c}{a} = \dfrac{5}{4} > 1$ (cohérent avec une hyperbole) ;
\item Directrices : $x = \pm \dfrac{a^2}{c} = \pm \dfrac{16}{5}$ ;
\item Asymptotes : $y = \pm \dfrac{b}{a}x = \pm \dfrac{3}{4}x$.
\end{itemize}

\textbf{3. Propriété bifocale.} Pour tout point $M$ de l'hyperbole :
\[ |MF - MF'| = 2a = 8. \]
\end{exoresolu}

\begin{exoresolu}[Une parabole et son paramètre]
Dans un repère orthonormé, on considère la parabole $(\mathcal{P})$ d'équation $y^2 = 12x$.
\begin{enumerate}
\item Déterminer son paramètre $p$, son foyer et sa directrice.
\item Le point $M_0(3\,;\,y_0)$ avec $y_0 > 0$ appartient-il à $(\mathcal{P})$ ? Si oui, calculer $M_0F$ et vérifier la définition monofocale.
\end{enumerate}
\tcblower
\textbf{1. Éléments.} L'équation est de la forme $y^2 = 2px$ avec $2p = 12$, donc $p = 6$. Le foyer est $F\left(\dfrac{p}{2}\,;\,0\right) = F(3\,;\,0)$ et la directrice a pour équation $x = -\dfrac{p}{2} = -3$.

\textbf{2. Appartenance et vérification.} Pour $x_0 = 3$ : $y_0^2 = 12 \times 3 = 36$, donc $y_0 = 6$ (car $y_0 > 0$). Le point $M_0(3\,;\,6)$ appartient bien à $(\mathcal{P})$. Alors :
\[ M_0F = \sqrt{(3-3)^2 + (6-0)^2} = 6, \qquad d(M_0, \mathcal{D}) = |3 - (-3)| = 6. \]
On a bien $M_0F = d(M_0, \mathcal{D})$, ce qui confirme que l'excentricité vaut $e = 1$ : la définition monofocale est vérifiée.
\end{exoresolu}

\begin{savaistu}[Des orbites képlériennes aux antennes de Yaoundé]
Les équations réduites que tu viens d'apprendre décrivent le monde qui t'entoure. Les planètes décrivent des \textbf{ellipses} dont le Soleil occupe un foyer (première loi de Kepler) : pour la Terre, $e \approx 0{,}017$, une ellipse presque circulaire. Les comètes non périodiques suivent des trajectoires \textbf{hyperboliques} ($e > 1$) : elles passent une fois près du Soleil puis s'éloignent pour toujours. Quant aux antennes paraboliques que l'on voit sur les toits de Douala et de Yaoundé, leur miroir est un \textbf{paraboloïde} : toutes les ondes arrivant parallèlement à l'axe se réfléchissent vers le foyer, où est placé le récepteur. Une seule famille de courbes, trois géométries, et des applications immenses.
\end{savaistu}

\begin{exercice}[À toi de jouer]
Pour chacune des coniques suivantes, déterminer la nature, l'axe focal, les sommets, les foyers, l'excentricité et les directrices (ainsi que les asymptotes le cas échéant) :
\[ (\mathcal{C}_1) : \frac{x^2}{100} + \frac{y^2}{36} = 1 \qquad ; \qquad
(\mathcal{C}_2) : \frac{x^2}{4} - \frac{y^2}{5} = 1 \qquad ; \qquad
(\mathcal{C}_3) : y^2 = 10x. \]
\end{exercice}

\begin{corrige}[Exercice : À toi de jouer]
\textbf{$(\mathcal{C}_1)$ :} deux carrés de même signe $\Rightarrow$ ellipse. $a^2 = 100$ (le plus grand, sous $x^2$), $b^2 = 36$, donc $a = 10$, $b = 6$, axe focal horizontal. $c^2 = 100 - 36 = 64$, $c = 8$, $e = \dfrac{8}{10} = \dfrac{4}{5}$. Sommets $(\pm 10\,;\,0)$ et $(0\,;\,\pm 6)$ ; foyers $(\pm 8\,;\,0)$ ; directrices $x = \pm \dfrac{100}{8} = \pm \dfrac{25}{2}$.

\textbf{$(\mathcal{C}_2)$ :} signes contraires $\Rightarrow$ hyperbole. $a^2 = 4$, $b^2 = 5$, donc $a = 2$, $b = \sqrt{5}$. $c^2 = 4 + 5 = 9$, $c = 3$, $e = \dfrac{3}{2}$. Sommets $(\pm 2\,;\,0)$ ; foyers $(\pm 3\,;\,0)$ ; directrices $x = \pm \dfrac{4}{3}$ ; asymptotes $y = \pm \dfrac{\sqrt{5}}{2}x$.

\textbf{$(\mathcal{C}_3)$ :} un seul carré $\Rightarrow$ parabole. $2p = 10$ donc $p = 5$ ; foyer $F\left(\dfrac{5}{2}\,;\,0\right)$ ; directrice $x = -\dfrac{5}{2}$ ; sommet $O(0\,;\,0)$ ; $e = 1$.
\end{corrige}

Tu sais désormais lire l'équation réduite d'une conique comme une fiche d'identité complète. Reste un problème : que faire lorsque l'équation n'est \emph{pas} donnée sous forme réduite, mais sous la forme générale $ax^2 + by^2 + cx + dy + e = 0$ ? C'est l'objet de la section 5, après que nous aurons étudié les tangentes.

\coursec{Tangentes à une conique}

Dans la section précédente, nous avons appris à reconnaître chaque conique par son équation réduite et à en déterminer tous les éléments caractéristiques : foyers, directrices, sommets, axes. Nous savons donc désormais \emph{où se trouve} une conique dans le plan et \emph{quelle forme} elle a. Il reste une question fondamentale, celle qui relie la géométrie à l'analyse et à la physique : \textbf{comment une droite peut-elle « toucher » une conique sans la traverser ?} C'est le problème des tangentes. Tu connais déjà la tangente à une courbe de fonction en un point ; ici, la conique n'est pas toujours le graphe d'une fonction (pense à l'ellipse, qui est « fermée »), il faut donc une approche purement géométrique et algébrique. Cette approche — la condition de racine double — est un grand classique du Baccalauréat C et E. Et à la clé, tu découvriras pourquoi les antennes paraboliques de ta maison captent si bien les signaux satellite.

\courssub{Notion de tangente à une conique}

\textbf{Intuition : de la sécante à la tangente}\par

Considérons une conique $\mathcal{C}$ et un point $M_0$ situé \emph{sur} cette conique. Prenons un second point $M$ de la conique, distinct de $M_0$ : la droite $(M_0M)$ est une \cle{sécante}, elle coupe la conique en deux points. Faisons maintenant glisser le point $M$ le long de la conique, en le rapprochant de plus en plus de $M_0$. La sécante pivote autour de $M_0$ et se rapproche d'une position limite : une droite qui « effleure » la conique en $M_0$ sans la traverser au voisinage de ce point. Cette position limite, c'est la tangente.

\begin{definition}[Tangente à une conique en un point]
Soit $\mathcal{C}$ une conique et $M_0$ un point de $\mathcal{C}$. La \cle{tangente} à $\mathcal{C}$ en $M_0$ est la position limite de la sécante $(M_0M)$ lorsque le point $M$ de $\mathcal{C}$ tend vers $M_0$ (notion admise). On dit que la tangente et la conique ont en $M_0$ un \emph{contact double} : $M_0$ est en quelque sorte un point d'intersection « compté deux fois ».
\end{definition}

Cette dernière remarque est la clé algébrique de toute la section : dire qu'une droite est tangente à une conique, c'est dire que le système « droite $\cap$ conique » admet une solution \emph{double}. Or une conique a une équation du second degré, et une droite une équation du premier degré : en substituant, on obtient une équation du second degré dont la solution double correspond exactement à un \textbf{discriminant nul}. Voilà le pont entre la géométrie et le calcul.

\courssub{Méthode algébrique : la condition de racine double ($\Delta = 0$)}

\begin{methode}[Déterminer une tangente par la condition $\Delta = 0$ — méthode exigible au Bac]
Pour déterminer l'équation de la tangente en $M_0(x_0\,;\,y_0)$ à une conique $\mathcal{C}$ d'équation connue :
\begin{enumerate}
\item \textbf{Vérifier} que $M_0 \in \mathcal{C}$ (sinon la question n'a pas de sens direct, voir l'attention plus bas).
\item Écrire l'équation d'une droite $(T)$ passant par $M_0$, de pente $m$ : $y - y_0 = m(x - x_0)$. (Traiter à part le cas d'une tangente verticale $x = x_0$ si le contexte le suggère.)
\item \textbf{Substituer} l'expression de $y$ (ou de $x$) dans l'équation de la conique : on obtient une équation du second degré en $x$ (ou en $y$), dont les coefficients dépendent de $m$.
\item Dire que $(T)$ est tangente si et seulement si cette équation admet une \textbf{racine double}, c'est-à-dire si son discriminant est nul : $\Delta = 0$.
\item Résoudre l'équation $\Delta = 0$ d'inconnue $m$, puis conclure en donnant l'équation de la tangente. Vérifier enfin que la racine double obtenue vaut bien $x_0$.
\end{enumerate}
\end{methode}

\begin{exoresolu}[Tangente à une parabole par la condition $\Delta = 0$]
\textbf{Énoncé.} Soit $\mathcal{P}$ la parabole d'équation $y^2 = 8x$. Déterminer l'équation de la tangente à $\mathcal{P}$ au point $M_0(2\,;\,4)$.
\tcblower
\textbf{Solution détaillée.}

\textbf{Étape 1 : appartenance.} $y_0^2 = 16$ et $8x_0 = 16$ : on a bien $M_0 \in \mathcal{P}$.

\textbf{Étape 2 : droite passant par $M_0$.} Une droite de pente $m$ passant par $M_0$ a pour équation $y - 4 = m(x - 2)$, soit $y = mx - 2m + 4$.

\textbf{Étape 3 : substitution.} Remplaçons $y$ dans $y^2 = 8x$ :
\begin{align*}
(mx - 2m + 4)^2 &= 8x \\
m^2x^2 + 2m(4 - 2m)x + (4 - 2m)^2 - 8x &= 0 \\
m^2x^2 + \bigl(8m - 4m^2 - 8\bigr)x + (4 - 2m)^2 &= 0.
\end{align*}

\textbf{Étape 4 : discriminant nul.} Cette équation en $x$ admet une racine double si et seulement si $\Delta = 0$ :
\begin{align*}
\Delta &= \bigl(8m - 4m^2 - 8\bigr)^2 - 4m^2(4 - 2m)^2 \\
&= 16\bigl(2m - m^2 - 2\bigr)^2 - 16m^2(2 - m)^2 \\
&= 16\Bigl[\bigl(m^2 - 2m + 2\bigr)^2 - m^2(m - 2)^2\Bigr].
\end{align*}
Posons $A = m^2 - 2m + 2$ et $B = m(m-2) = m^2 - 2m$. Alors $\Delta = 16(A^2 - B^2) = 16(A - B)(A + B) = 16 \times 2 \times (2m^2 - 4m + 2) = 64(m^2 - 2m + 1) = 64(m - 1)^2$.

La condition $\Delta = 0$ donne donc $(m-1)^2 = 0$, c'est-à-dire $m = 1$.

\textbf{Étape 5 : conclusion.} La tangente a pour équation $y - 4 = 1 \times (x - 2)$, soit
\[ y = x + 2, \quad \text{que l'on peut écrire } 4y = 4(x + 2). \]
Vérification : avec $m = 1$, l'équation du second degré devient $x^2 - 4x + 4 = 0$, soit $(x - 2)^2 = 0$ : la racine double est bien $x_0 = 2$. Le contact est double en $M_0$, comme prévu.
\end{exoresolu}

\courssub{Les formules de dédoublement}

La méthode précédente est sûre mais parfois longue. Il existe un raccourci remarquable : les \cle{formules de dédoublement}. L'idée est simple à retenir : dans l'équation réduite de la conique, on « dédouble » chaque terme quadratique selon les règles
\[ x^2 \longrightarrow x\,x_0, \qquad y^2 \longrightarrow y\,y_0, \qquad x \longrightarrow \frac{x + x_0}{2}, \qquad y \longrightarrow \frac{y + y_0}{2}. \]

\begin{propriete}[Formules de dédoublement — admises]
Soit $M_0(x_0\,;\,y_0)$ un point d'une conique $\mathcal{C}$ donnée par son équation réduite. La tangente à $\mathcal{C}$ en $M_0$ a pour équation :
\begin{itemize}
\item \textbf{Ellipse} $\dfrac{x^2}{a^2} + \dfrac{y^2}{b^2} = 1$ : \quad $\displaystyle \frac{x\,x_0}{a^2} + \frac{y\,y_0}{b^2} = 1$ ;
\item \textbf{Hyperbole} $\dfrac{x^2}{a^2} - \dfrac{y^2}{b^2} = 1$ : \quad $\displaystyle \frac{x\,x_0}{a^2} - \frac{y\,y_0}{b^2} = 1$ ;
\item \textbf{Parabole} $y^2 = 2px$ : \quad $\displaystyle y\,y_0 = p\,(x + x_0)$.
\end{itemize}
Ces formules sont admises ; leur vérification sur un exemple (via $\Delta = 0$) peut être demandée.
\end{propriete}

\begin{experience}[Démonstration guidée sur un exemple : la parabole $y^2 = 2px$]
Montrons que la formule de dédoublement redonne exactement le résultat de la méthode $\Delta = 0$ pour la parabole $y^2 = 2px$ en un point $M_0(x_0\,;\,y_0)$ de la parabole (donc $y_0^2 = 2px_0$), avec $y_0 \neq 0$.
\begin{enumerate}
\item La droite $y\,y_0 = p(x + x_0)$ donne $x = \dfrac{y\,y_0}{p} - x_0$. Substituons dans $y^2 = 2px$ :
\[ y^2 = 2p\left(\frac{y\,y_0}{p} - x_0\right) = 2y\,y_0 - 2px_0. \]
\item D'où $y^2 - 2y_0\,y + 2px_0 = 0$. Or $2px_0 = y_0^2$, donc l'équation s'écrit $y^2 - 2y_0\,y + y_0^2 = 0$, c'est-à-dire $(y - y_0)^2 = 0$.
\item L'équation admet la racine \textbf{double} $y = y_0$ : la droite et la parabole ont un contact double en $M_0$. C'est bien la tangente. $\square$
\end{enumerate}
Retrouve-toi dans le calcul : applique la formule à $y^2 = 8x$ (donc $2p = 8$, $p = 4$) au point $M_0(2\,;\,4)$ : on obtient $4y = 4(x + 2)$, soit $y = x + 2$ — exactement le résultat de l'exercice résolu précédent, obtenu en une ligne au lieu d'une page !
\end{experience}

\begin{attention}[Le dédoublement exige que $M_0$ appartienne à la conique]
La formule de dédoublement n'est valable que si $M_0(x_0\,;\,y_0)$ \textbf{appartient} à la conique. Avant tout calcul, vérifie que les coordonnées de $M_0$ satisfont l'équation de la conique. Si $M_0$ est extérieur, la droite obtenue par dédoublement existe toujours (on l'appelle la \emph{polaire} de $M_0$), mais ce n'est \textbf{pas} une tangente passant par $M_0$ : c'est une erreur classique et éliminatoire au Bac. Autre piège : pour la parabole $y^2 = 2px$, n'oublie pas le facteur $p$ devant $(x + x_0)$ — et identifie correctement $p$ : pour $y^2 = 8x$, on a $p = 4$ et non $p = 8$.
\end{attention}

\begin{exemplebox}[Tangente à une ellipse par dédoublement]
Soit $\mathcal{E}$ l'ellipse d'équation $\dfrac{x^2}{25} + \dfrac{y^2}{9} = 1$ et $M_0(3\,;\,\tfrac{12}{5})$. Vérifions l'appartenance : $\dfrac{9}{25} + \dfrac{144/25}{9} = \dfrac{9}{25} + \dfrac{16}{25} = 1$. Bien. La tangente en $M_0$ est :
\[ \frac{3x}{25} + \frac{(12/5)\,y}{9} = 1 \iff \frac{3x}{25} + \frac{4y}{15} = 1 \iff 9x + 20y = 75. \]
\end{exemplebox}

\courssub{Tangentes de pente donnée}

Un deuxième type de problème, très fréquent au Baccalauréat, consiste à chercher les tangentes à une conique dont on impose la \emph{pente} (par exemple les tangentes parallèles à une droite donnée). La stratégie reste la condition $\Delta = 0$, mais les rôles sont inversés : la pente $m$ est connue, et c'est l'ordonnée à l'origine qui devient l'inconnue.

\begin{methode}[Tangentes de pente $m$ imposée]
\begin{enumerate}
\item Écrire la droite $(D) : y = mx + q$ avec $m$ connu et $q$ inconnu.
\item Substituer dans l'équation de la conique pour obtenir une équation du second degré en $x$, dont les coefficients dépendent de $q$.
\item Écrire $\Delta = 0$ : on obtient une équation en $q$, qui donne en général \textbf{deux} valeurs $q_1$ et $q_2$ — une conique possède deux tangentes de pente donnée, symétriques l'une de l'autre.
\item Conclure par les deux équations de tangentes.
\end{enumerate}
\end{methode}

\begin{exoresolu}[Tangentes de pente donnée à une ellipse]
\textbf{Énoncé.} Déterminer les tangentes à l'ellipse $\mathcal{E} : \dfrac{x^2}{9} + \dfrac{y^2}{4} = 1$ qui sont parallèles à la droite d'équation $y = 2x$.
\tcblower
\textbf{Solution détaillée.}

\textbf{Étape 1.} On cherche les droites $y = 2x + q$ tangentes à $\mathcal{E}$.

\textbf{Étape 2.} Substituons : $\dfrac{x^2}{9} + \dfrac{(2x+q)^2}{4} = 1$. Multiplions par $36$ :
\begin{align*}
4x^2 + 9(2x + q)^2 &= 36 \\
4x^2 + 9(4x^2 + 4qx + q^2) &= 36 \\
40x^2 + 36q\,x + 9q^2 - 36 &= 0.
\end{align*}

\textbf{Étape 3.} Condition de tangence : $\Delta = 0$.
\[ \Delta = (36q)^2 - 4 \times 40 \times (9q^2 - 36) = 1296q^2 - 1440q^2 + 5760 = -144q^2 + 5760. \]
$\Delta = 0 \iff q^2 = \dfrac{5760}{144} = 40 \iff q = \pm 2\sqrt{10}$.

\textbf{Étape 4.} Il existe deux tangentes de pente $2$ :
\[ y = 2x + 2\sqrt{10} \qquad \text{et} \qquad y = 2x - 2\sqrt{10}. \]
Elles sont symétriques par rapport au centre $O$ de l'ellipse, ce qui est cohérent avec la symétrie centrale de $\mathcal{E}$.
\end{exoresolu}

\begin{aretenir}[Résultat utile : tangentes de pente $m$ à l'ellipse et à l'hyperbole]
Pour l'ellipse $\dfrac{x^2}{a^2} + \dfrac{y^2}{b^2} = 1$, les tangentes de pente $m$ sont $y = mx \pm \sqrt{a^2m^2 + b^2}$ : elles existent pour toute pente $m$. Pour l'hyperbole $\dfrac{x^2}{a^2} - \dfrac{y^2}{b^2} = 1$, on obtient $y = mx \pm \sqrt{a^2m^2 - b^2}$, qui n'existent que si $|m| \geq \dfrac{b}{a}$ : c'est l'ombre des asymptotes, dont les pentes sont précisément $\pm \dfrac{b}{a}$.
\end{aretenir}

\courssub{Propriétés optiques des coniques}

\textbf{La propriété optique de la parabole}\par

Voici l'une des plus belles rencontres entre la géométrie et le monde réel.

\begin{propriete}[Propriété optique de la parabole — admise]
Soit $\mathcal{P}$ une parabole de foyer $F$ et d'axe $\Delta$, et $M_0$ un point de $\mathcal{P}$. La tangente en $M_0$ fait des angles égaux avec la droite $(M_0F)$ et avec la parallèle à l'axe menée par $M_0$. Autrement dit, la normale en $M_0$ est la bissectrice de l'angle formé par le rayon focal $(M_0F)$ et la parallèle à l'axe. Conséquence : \textbf{tout rayon lumineux parallèle à l'axe de la parabole se réfléchit sur celle-ci en passant par le foyer} (et réciproquement, tout rayon issu du foyer se réfléchit parallèlement à l'axe).
\end{propriete}

\begin{popfigure}
\begin{tikzpicture}[scale=0.9]
% Parabole y^2 = 4x, p=2, F(1,0)
\draw[->, popmuted] (-0.5,0) -- (5.5,0) node[below left] {$x$};
\draw[->, popmuted] (0,-3) -- (0,5) node[below left] {$y$};
\draw[popblue, very thick, domain=-3.5:3.5, samples=100, smooth] plot ({0.25*\x*\x}, {\x});
\draw[popline, dashed] (-1,-3) -- (-1,5) node[above, popmuted] {\footnotesize directrice $x=-1$};
% Foyer
\fill[poporange] (1,0) circle (2.5pt) node[below right, poporange] {$F(1\,;\,0)$};
% Point M0 (4,4) sur y^2=4x
\coordinate (M0) at (4,4);
\fill[popink] (M0) circle (2.5pt) node[above right, popink] {$M_0(4\,;\,4)$};
% Tangente en M0 : y*y0 = 2(x+x0) => 4y = 2(x+4) => y = 0.5x + 2
\draw[popgreen, very thick, domain=-1:5] plot (\x, {0.5*\x + 2});
\node[popgreen, anchor=south west] at (4.5,4.2) {\footnotesize tangente $y=\frac{1}{2}x+2$};
% Rayon incident parallèle à l'axe (horizontal, venant de la gauche)
\draw[->, popgold, very thick] (-0.5,4) -- (M0);
\node[popgold, anchor=south] at (1.5,4.3) {\footnotesize rayon incident $\parallel$ axe};
% Rayon réfléchi vers F
\draw[->, poppink, very thick] (M0) -- (1,0);
\node[poppink, anchor=north west] at (2.5,1.5) {\footnotesize rayon réfléchi vers $F$};
% Normale (perpendiculaire à la tangente, pente -2)
\draw[poppurple, dashed, thick] (M0) -- ++(-1,2) node[above left, poppurple] {\footnotesize normale};
% Angles égaux (marquages)
\draw[popmuted] (3.6,3.8) arc (180:135:0.4);
\draw[popmuted] (3.6,3.8) arc (135:90:0.4);
\draw[popmuted] (3.2,3.2) arc (-45:-90:0.35);
\draw[popmuted] (3.2,3.2) arc (-90:-135:0.35);
\end{tikzpicture}
\legende{Parabole $y^2 = 4x$ de foyer $F(1\,;\,0)$ : le rayon incident, parallèle à l'axe, se réfléchit en $M_0(4\,;\,4)$ vers le foyer $F$. La tangente (verte) fait des angles égaux avec les deux rayons ; la normale (violette, pointillée) est la bissectrice de l'angle $\widehat{FM_0I}$ où $I$ est le projeté de $M_0$ sur la directrice.}
\end{popfigure}

\begin{savaistu}[Antennes, phares et fours solaires : la parabole partout]
La propriété optique de la parabole explique des objets de ta vie quotidienne. L'\textbf{antenne parabolique} que l'on voit sur les toits de Yaoundé ou de Garoua capte des ondes satellite qui arrivent toutes parallèles à l'axe de la parabole : elles convergent toutes au foyer, où l'on place le récepteur (la tête LNB). À l'inverse, dans un \textbf{phare de voiture} ou une lampe torche, la source lumineuse est placée au foyer : les rayons ressortent tous parallèles à l'axe, formant un faisceau puissant et rectiligne. Même principe pour les \textbf{fours solaires} (comme ceux utilisés dans certaines zones rurales pour la cuisson) et les radiotélescopes géants. Une seule propriété géométrique, démontrée par les Grecs il y a plus de deux mille ans, et des applications qui couvrent la planète !
\end{savaistu}

\textbf{Propriétés optiques de l'ellipse et de l'hyperbole}\par

\begin{propriete}[Propriété optique de l'ellipse — mention culturelle]
Soit $\mathcal{E}$ une ellipse de foyers $F$ et $F'$, et $M_0$ un point de $\mathcal{E}$. La tangente en $M_0$ est la bissectrice extérieure de l'angle $\widehat{FM_0F'}$ : elle fait des angles égaux avec les rayons focaux $(M_0F)$ et $(M_0F')$. Par conséquent, \textbf{tout rayon issu d'un foyer se réfléchit sur l'ellipse et passe par l'autre foyer}.
\end{propriete}

\begin{popfigure}
\begin{tikzpicture}[scale=1.0]
% Ellipse a=3, b=2
\draw[->, popmuted] (-3.8,0) -- (3.8,0) node[below left] {$x$};
\draw[->, popmuted] (0,-2.6) -- (0,2.6) node[below left] {$y$};
\draw[popblue, very thick] (0,0) ellipse (3 and 2);
% Foyers c = sqrt(9-4) = sqrt5 ~ 2.236
\fill[poporange] (-2.236,0) circle (2pt) node[below, poporange] {$F$};
\fill[poporange] (2.236,0) circle (2pt) node[below, poporange] {$F'$};
% Point M0 : angle 50° => (3cos50, 2sin50) = (1.928, 1.532)
\coordinate (M0) at (1.928,1.532);
\fill[popink] (M0) circle (2pt) node[above right, popink] {$M_0$};
% Rayons focaux
\draw[popgold, very thick, ->] (-2.236,0) -- (M0);
\draw[poppink, very thick, ->] (M0) -- (2.236,0);
% Tangente : x*x0/9 + y*y0/4 = 1 => 0.2142 x + 0.383 y = 1
% direction de la tangente : (-0.383, 0.2142) normalisée
\draw[popgreen, very thick, domain=-1.4:1.4] plot ({1.928 - 3.1*\x}, {1.532 + 1.72*\x});
\node[popgreen] at (-2.4,2.1) {\footnotesize tangente en $M_0$};
\node[popmuted] at (0,-2.4) {\footnotesize $\dfrac{x^2}{9}+\dfrac{y^2}{4}=1$};
\end{tikzpicture}
\legende{Ellipse $\frac{x^2}{9}+\frac{y^2}{4}=1$ de foyers $F$ et $F'$ : un rayon issu de $F$ se réfléchit en $M_0$ et passe par $F'$. La tangente fait des angles égaux avec $(M_0F)$ et $(M_0F')$.}
\end{popfigure}

\begin{savaistu}[Galeries des échos et lithotripsie]
La propriété optique de l'ellipse a des applications étonnantes. Dans une \textbf{galerie des échos} (certaines salles voûtées en ellipse), deux personnes placées aux foyers peuvent se chuchoter des secrets inaudibles ailleurs : le son émis en $F$ converge en $F'$. En médecine, la \textbf{lithotripsie} détruit les calculs rénaux sans chirurgie : on place le patient de sorte que l'onde de choc émise en un foyer d'une chambre ellipsoïdale converge exactement sur le calcul, placé à l'autre foyer. L'hyperbole possède, elle aussi, une propriété optique : un rayon dirigé vers un foyer se réfléchit comme s'il provenait de l'autre foyer — principe utilisé dans certains télescopes (monture de Cassegrain). Les coniques, étudiées par Apollonios de Perga vers 200 av. J.-C., sont ainsi devenues des outils de haute technologie.
\end{savaistu}

\begin{attention}[Bien distinguer tangente et normale]
La \cle{normale} en $M_0$ est la perpendiculaire à la tangente en $M_0$. Dans les propriétés optiques, c'est souvent la normale qui est bissectrice (intérieure) des rayons focaux, et la tangente qui est bissectrice extérieure. Au Bac, si l'on te demande la normale, calcule d'abord la tangente (par dédoublement), puis utilise le fait que deux droites perpendiculaires ont des pentes dont le produit vaut $-1$ (dans un repère orthonormé).
\end{attention}

\begin{exercice}[Pour t'entraîner]
\begin{enumerate}
\item Soit $\mathcal{H}$ l'hyperbole d'équation $\dfrac{x^2}{16} - \dfrac{y^2}{9} = 1$. Déterminer l'équation de la tangente à $\mathcal{H}$ au point $M_0\left(5\,;\,\dfrac{9}{4}\right)$ après avoir vérifié que $M_0 \in \mathcal{H}$.
\item Déterminer les tangentes à la parabole $y^2 = 4x$ qui sont perpendiculaires à la droite d'équation $y = -x + 3$.
\item Soit $\mathcal{E}$ l'ellipse $\dfrac{x^2}{8} + \dfrac{y^2}{4} = 1$. Déterminer les tangentes à $\mathcal{E}$ de pente $1$ et préciser leurs points de contact.
\end{enumerate}
\end{exercice}

\begin{corrige}[Exercice n°1]
\textbf{1.} Vérification : $\dfrac{25}{16} - \dfrac{81/16}{9} = \dfrac{25}{16} - \dfrac{9}{16} = \dfrac{16}{16} = 1$ : $M_0 \in \mathcal{H}$. Dédoublement : $\dfrac{5x}{16} - \dfrac{(9/4)y}{9} = 1 \iff \dfrac{5x}{16} - \dfrac{y}{4} = 1 \iff 5x - 4y = 16$.

\textbf{2.} Les tangentes cherchées ont pente $m = 1$ (perpendiculaires à une droite de pente $-1$). Soit $y = x + q$ ; substitution dans $y^2 = 4x$ : $(x+q)^2 = 4x \iff x^2 + (2q - 4)x + q^2 = 0$. Condition : $\Delta = (2q-4)^2 - 4q^2 = -16q + 16 = 0 \iff q = 1$. Unique tangente : $y = x + 1$ (une parabole n'a qu'une tangente de pente donnée non nulle). Point de contact : racine double $x = -\dfrac{2q-4}{2} = 1$, $y = 2$ : le point $(1\,;\,2)$.

\textbf{3.} On cherche les tangentes de pente $m=1$ à l'ellipse $\dfrac{x^2}{8} + \dfrac{y^2}{4} = 1$. D'après la formule du cours : $y = x \pm \sqrt{a^2m^2 + b^2} = x \pm \sqrt{8 \times 1 + 4} = x \pm 2\sqrt{3}$.
Les deux tangentes sont :
\[ (T_1) : y = x + 2\sqrt{3} \qquad \text{et} \qquad (T_2) : y = x - 2\sqrt{3}. \]
\textbf{Point de contact de $(T_1)$ :} Substituons $y = x + 2\sqrt{3}$ dans l'équation de l'ellipse :
\[ \frac{x^2}{8} + \frac{(x+2\sqrt{3})^2}{4} = 1 \iff x^2 + 2(x^2 + 4\sqrt{3}x + 12) = 8 \iff 3x^2 + 8\sqrt{3}x + 16 = 0. \]
Le discriminant est nul (par construction), la racine double est $x_0 = -\dfrac{8\sqrt{3}}{2 \times 3} = -\dfrac{4\sqrt{3}}{3}$.
Alors $y_0 = x_0 + 2\sqrt{3} = \dfrac{2\sqrt{3}}{3}$. Point de contact : $\left(-\dfrac{4\sqrt{3}}{3}\,;\,\dfrac{2\sqrt{3}}{3}\right)$.
\textbf{Point de contact de $(T_2)$ :} Par symétrie centrale (centre $O$), le point de contact est l'opposé : $\left(\dfrac{4\sqrt{3}}{3}\,;\,-\dfrac{2\sqrt{3}}{3}\right)$.
\end{corrige}

\begin{aretenir}[L'essentiel de la section]
\begin{itemize}
\item La tangente en $M_0$ est la position limite des sécantes $(M_0M)$ ; algébriquement, elle correspond à un \textbf{contact double}, donc à un discriminant nul.
\item Méthode exigible : droite $y - y_0 = m(x - x_0)$, substitution, condition $\Delta = 0$.
\item Formules de dédoublement (admises) : $x^2 \to xx_0$, $y^2 \to yy_0$, $x \to \frac{x+x_0}{2}$ — valables \textbf{uniquement si $M_0$ est sur la conique}.
\item Tangentes de pente $m$ : on cherche $q$ dans $y = mx + q$ par $\Delta = 0$ ; deux solutions en général pour l'ellipse et l'hyperbole, une seule pour la parabole.
\item Propriété optique de la parabole : tout rayon parallèle à l'axe se réfléchit vers le foyer — c'est le principe des antennes paraboliques et des phares.
\end{itemize}
\end{aretenir}

\coursec{Reconnaissance d'une équation du second degré}

Dans la section précédente, nous avons appris à déterminer l'équation de la tangente à une conique \emph{donnée sous forme réduite} : $\dfrac{x^2}{a^2}+\dfrac{y^2}{b^2}=1$ pour une ellipse, $\dfrac{x^2}{a^2}-\dfrac{y^2}{b^2}=1$ pour une hyperbole, $y^2=2px$ pour une parabole. Mais dans la pratique — et très souvent au Baccalauréat —, l'énoncé ne vous offre pas ce confort : on vous tend une équation « en vrac », du type
\[ 4x^2+9y^2-16x+18y-11=0, \]
et l'on vous demande : \emph{quelle est la nature de cet ensemble de points ? Donnez ses éléments caractéristiques.}

Toute cette section repose sur une idée simple et puissante : une conique déplacée dans le plan (son centre n'est plus à l'origine) voit son équation « se compliquer » par des termes en $x$ et en $y$. Retrouver la forme réduite, c'est donc \cle{défaire} cette translation : on complète les carrés, puis on translate le repère vers le centre (ou le sommet) de la conique. Autrement dit, nous allons apprendre à lire une équation du second degré comme on décode un message.

\courssub{L'équation générale étudiée et la complétion des carrés}

\begin{definition}[Équation cartésienne du second degré sans terme en $xy$]
On appelle \cle{équation générale d'une conique} (sans terme rectangle) toute équation de la forme
\[ ax^2+by^2+cx+dy+e=0, \]
où $a$, $b$, $c$, $d$, $e$ sont des réels, avec $(a,b)\neq(0,0)$, et où l'inconnue est le couple $(x,y)$ des coordonnées d'un point $M$ du plan rapporté à un repère orthonormé $(O\,;\vec{\imath},\vec{\jmath})$.
\end{definition}

\begin{attention}[Pas de terme en $xy$]
Nous nous limitons ici aux équations \textbf{sans terme en $xy$}. Un terme en $xy$ traduirait une \emph{rotation} des axes de la conique : sa prise en compte exige des outils qui dépassent le programme de Terminale. Toutes les coniques de cette section ont donc leurs axes \textbf{parallèles aux axes du repère}.
\end{attention}

L'outil central de toute la section est la \cle{complétion du carré}, que tu as rencontrée dès la classe de Seconde avec le trinôme :
\[ x^2+\alpha x = \left(x+\dfrac{\alpha}{2}\right)^2-\dfrac{\alpha^2}{4}. \]
Appliquée séparément aux termes en $x$ et aux termes en $y$, elle fait apparaître des quantités du type $(x-x_0)^2$ et $(y-y_0)^2$, c'est-à-dire les coordonnées dans un repère translaté, d'origine $\Omega(x_0\,;y_0)$.

\begin{methode}[De l'équation générale à la forme réduite]
Pour reconnaître l'ensemble des points $M(x\,;y)$ vérifiant $ax^2+by^2+cx+dy+e=0$ :
\begin{enumerate}
\item \textbf{Factoriser} $a$ dans le groupe des termes en $x$ et $b$ dans le groupe des termes en $y$ (si $a\neq 0$ et $b\neq 0$) :
\[ a\left(x^2+\dfrac{c}{a}x\right)+b\left(y^2+\dfrac{d}{b}y\right)+e=0. \]
\item \textbf{Compléter les carrés} dans chaque parenthèse :
\[ x^2+\dfrac{c}{a}x=\left(x+\dfrac{c}{2a}\right)^2-\dfrac{c^2}{4a^2}, \qquad y^2+\dfrac{d}{b}y=\left(y+\dfrac{d}{2b}\right)^2-\dfrac{d^2}{4b^2}. \]
\item \textbf{Regrouper les constantes} dans le second membre : on obtient
\[ a(x-x_0)^2+b(y-y_0)^2 = k, \quad \text{avec } x_0=-\dfrac{c}{2a},\ y_0=-\dfrac{d}{2b}. \]
\item \textbf{Translater le repère} : on pose $X=x-x_0$ et $Y=y-y_0$ (repère d'origine $\Omega(x_0\,;y_0)$). L'équation devient $aX^2+bY^2=k$.
\item \textbf{Conclure} selon les signes de $a$, $b$ et $k$ (voir la propriété ci-dessous), puis \textbf{traduire les éléments caractéristiques} (centre, sommets, foyers, directrices) dans le repère initial.
\end{enumerate}
\end{methode}

\begin{attention}[Factoriser avant de compléter !]
L'erreur la plus fréquente consiste à compléter le carré \emph{sans avoir factorisé} le coefficient de $x^2$. Par exemple, face à $4x^2-16x$, écrire directement $(2x-4)^2$ mélange tout. La démarche correcte est :
\[ 4x^2-16x = 4\left(x^2-4x\right)=4\left[(x-2)^2-4\right]=4(x-2)^2-16. \]
Retiens : \textbf{on factorise d'abord, on complète ensuite.}
\end{attention}

\courssub{Discussion selon les signes de $a$, $b$ et du second membre}

Après complétion, tout se ramène à l'étude de l'équation
\[ aX^2+bY^2=k \qquad (\text{ou } aX^2+dY+e'=0 \text{ si } b=0). \]
La nature de l'ensemble dépend des \textbf{signes} de $a$ et $b$, et de la \textbf{valeur} de $k$.

\begin{propriete}[Tableau de reconnaissance — démonstration immédiate par mise sous forme réduite]
Soit $(\mathcal{E})$ l'ensemble des points $M(x\,;y)$ vérifiant $ax^2+by^2+cx+dy+e=0$ avec $(a,b)\neq(0,0)$. Après complétion des carrés, on se ramène à $aX^2+bY^2=k$ (lorsque $ab\neq 0$).
\begin{center}
\begin{tabularx}{\linewidth}{|l|l|X|}
\hline
\rowcolor{popblueL} \textbf{Signes de $a$ et $b$} & \textbf{Second membre} & \textbf{Nature de $(\mathcal{E})$} \\
\hline
$a$ et $b$ de même signe ($ab>0$) & $k$ de même signe que $a$ et $b$ & Ellipse (cercle si $a=b$) \\
\hline
$a$ et $b$ de même signe & $k=0$ & Un point : $\{ \Omega \}$ \\
\hline
$a$ et $b$ de même signe & $k$ de signe opposé & Ensemble vide \\
\hline
$a$ et $b$ de signes contraires ($ab<0$) & $k\neq 0$ & Hyperbole \\
\hline
$a$ et $b$ de signes contraires & $k=0$ & Deux droites sécantes en $\Omega$ \\
\hline
$a=0$, $b\neq 0$ et $c\neq 0$ & --- & Parabole d'axe parallèle à $(Ox)$ \\
\hline
$a\neq 0$, $b=0$ et $d\neq 0$ & --- & Parabole d'axe parallèle à $(Oy)$ \\
\hline
$a=0$, $b\neq 0$, $c=0$ & selon le discriminant & Deux droites parallèles, une droite, ou $\varnothing$ \\
\hline
\end{tabularx}
\end{center}
\end{propriete}

\begin{experience}[Pourquoi ces conclusions ? Une justification rapide]
Supposons $a>0$, $b>0$ et $k>0$. En divisant $aX^2+bY^2=k$ par $k$, il vient
\[ \dfrac{X^2}{k/a}+\dfrac{Y^2}{k/b}=1, \]
qui est la forme réduite d'une \cle{ellipse} de centre $\Omega$, de demi-axes $\sqrt{k/a}$ et $\sqrt{k/b}$ (un \cle{cercle} si $a=b$, car les deux demi-axes sont alors égaux). Si $k=0$, l'équation $aX^2+bY^2=0$ impose $X=Y=0$ : un seul point. Si $k<0$, le premier membre est positif ou nul et le second strictement négatif : impossible, l'ensemble est vide.

Supposons maintenant $a>0$, $b<0$ et $k\neq 0$. En divisant par $k$, on obtient $\dfrac{X^2}{k/a}-\dfrac{Y^2}{k/|b|}=1$ (ou $-1$ selon le signe de $k$) : c'est une \cle{hyperbole}. Si $k=0$, alors $aX^2=|b|Y^2$, c'est-à-dire $Y=\pm\sqrt{a/|b|}\,X$ : deux droites sécantes en $\Omega$ — c'est le cas \emph{dégénéré}, les asymptotes réunies.

Enfin, si $b=0$ et $d\neq 0$, l'équation $aX^2+dY+e'=0$ s'écrit $Y=-\dfrac{a}{d}X^2-\dfrac{e'}{d}$ : une \cle{parabole} d'axe vertical. Si de plus $d=0$, il ne reste que $aX^2+e'=0$, équation en $X$ seul : deux droites verticales, une droite « double », ou l'ensemble vide.
\end{experience}

\begin{attention}[Ne conclus pas trop vite !]
Après complétion des carrés, \textbf{regarde toujours le second membre avant de nommer la conique}. Dire « $a$ et $b$ de même signe, donc ellipse » est faux en général : $2x^2+3y^2+5=0$ ne représente \emph{aucun} point, et $2x^2+3y^2=0$ un seul. De même, $x^2-y^2=0$ n'est pas une hyperbole mais le couple de droites $y=x$ et $y=-x$. La nature exacte se décide \emph{après} la mise sous forme $aX^2+bY^2=k$.
\end{attention}

\begin{popfigure}
\begin{center}
\begin{tikzpicture}[
  grow=right,
  level distance=3.6cm,
  level 1/.style={sibling distance=3.6cm},
  level 2/.style={sibling distance=1.7cm},
  every node/.style={font=\small},
  decision/.style={draw=popblue, fill=popblueL, rounded corners, thick, align=center, inner sep=4pt},
  result/.style={draw=popdark, fill=popcream, rounded corners, align=center, inner sep=3pt, font=\footnotesize},
  edge from parent/.style={draw=popmuted, thick, -{Stealth}}
]
\node[decision] {$ab$ ?}
  child { node[decision, align=center] {$a=0$ ou $b=0$\\ (exclusivement)}
    child { node[result, align=center] {coefficient du terme\\ du $1^{\text{er}}$ degré nul :\\ droites parallèles ou $\varnothing$} }
    child { node[result, align=center] {\textbf{Parabole}\\ (axe parallèle à $(Ox)$ ou $(Oy)$)} }
  }
  child { node[decision] {$ab<0$}
    child { node[result, align=center] {$k=0$ : deux droites\\ sécantes en $\Omega$} }
    child { node[result] {$k\neq 0$ : \textbf{hyperbole}} }
  }
  child { node[decision] {$ab>0$}
    child { node[result] {$k$ de signe opposé : $\varnothing$} }
    child { node[result] {$k=0$ : le point $\Omega$} }
    child { node[result, align=center] {$k$ de même signe :\\ \textbf{ellipse} (cercle si $a=b$)} }
  };
\end{tikzpicture}
\end{center}
\legende{Arbre de décision : la nature de l'ensemble $ax^2+by^2+cx+dy+e=0$ se lit sur les signes de $a$, $b$ et du second membre $k$ obtenu après complétion des carrés.}
\end{popfigure}

\begin{savaistu}[Un critère éclair pour le jour du Bac (hors strict programme)]
Une fois l'équation mise sous la forme $aX^2+bY^2=k$ avec $k\neq 0$ et du « bon » signe, le \textbf{signe du produit $a\times b$} suffit à trancher :
\begin{itemize}
\item $ab>0$ : ellipse (cercle si $a=b$) ;
\item $ab<0$ : hyperbole ;
\item $ab=0$ : parabole (ou cas dégénéré).
\end{itemize}
Ce réflexe de quelques secondes te permet de \emph{vérifier} ta conclusion détaillée. Attention : il ne dispense jamais de la rédaction complète (complétion des carrés, translation du repère), qui est ce que le correcteur attend.
\end{savaistu}

\courssub{Exemple chiffré complet : étude de $4x^2+9y^2-16x+18y-11=0$}

Traitons intégralement un exemple, en suivant la méthode pas à pas, comme tu devras le faire en copie.

\begin{exoresolu}[Reconnaissance complète d'une ellipse]
\textbf{Énoncé.} Le plan est rapporté à un repère orthonormé $(O\,;\vec{\imath},\vec{\jmath})$. On considère l'ensemble $(\mathcal{E})$ des points $M(x\,;y)$ tels que
\[ 4x^2+9y^2-16x+18y-11=0. \]
\begin{enumerate}
\item Montrer que $(\mathcal{E})$ est une ellipse dont on précisera le centre $\Omega$.
\item Déterminer ses demi-axes, ses sommets, ses foyers, ses directrices et son excentricité.
\end{enumerate}
\tcblower
\textbf{Solution détaillée.}

\textbf{1.} On factorise les coefficients de $x^2$ et $y^2$, puis on complète les carrés :
\begin{align*}
4x^2-16x &= 4\left(x^2-4x\right)=4\left[(x-2)^2-4\right]=4(x-2)^2-16,\\
9y^2+18y &= 9\left(y^2+2y\right)=9\left[(y+1)^2-1\right]=9(y+1)^2-9.
\end{align*}
L'équation devient :
\[ 4(x-2)^2-16+9(y+1)^2-9-11=0
\iff 4(x-2)^2+9(y+1)^2=36. \]
En divisant par $36$ :
\[ \boxed{\ \dfrac{(x-2)^2}{9}+\dfrac{(y+1)^2}{4}=1\ } \]
On pose $X=x-2$ et $Y=y+1$ : dans le repère d'origine $\Omega(2\,;-1)$, l'équation $\dfrac{X^2}{9}+\dfrac{Y^2}{4}=1$ est la forme réduite d'une \textbf{ellipse} de centre $\Omega(2\,;-1)$.

\textbf{2.} On lit $a^2=9$ et $b^2=4$, donc $a=3$ (grand axe, porté par l'axe des $X$) et $b=2$ (petit axe). Comme $a>b$, on a
\[ c=\sqrt{a^2-b^2}=\sqrt{9-4}=\sqrt{5}, \qquad e=\dfrac{c}{a}=\dfrac{\sqrt{5}}{3}. \]
Dans le repère $(\Omega\,;\vec{\imath},\vec{\jmath})$, les sommets sont $A(\pm 3\,;0)$ et $B(0\,;\pm 2)$, les foyers $F(\pm\sqrt{5}\,;0)$, et les directrices ont pour équations $X=\pm\dfrac{a^2}{c}=\pm\dfrac{9}{\sqrt{5}}=\pm\dfrac{9\sqrt{5}}{5}$.

On traduit dans le repère initial par $x=X+2$, $y=Y-1$ :
\begin{itemize}
\item sommets du grand axe : $A_1(5\,;-1)$ et $A_2(-1\,;-1)$ ;
\item sommets du petit axe : $B_1(2\,;1)$ et $B_2(2\,;-3)$ ;
\item foyers : $F_1(2+\sqrt{5}\,;-1)$ et $F_2(2-\sqrt{5}\,;-1)$ ;
\item directrices : $x=2+\dfrac{9\sqrt{5}}{5}$ et $x=2-\dfrac{9\sqrt{5}}{5}$ ;
\item excentricité : $e=\dfrac{\sqrt{5}}{3}\approx 0{,}745$, bien comprise entre $0$ et $1$ : cohérent pour une ellipse.
\end{itemize}
\end{exoresolu}

\begin{popfigure}
\begin{center}
\begin{tikzpicture}[scale=0.85]
\begin{axis}[
  width=0.9\linewidth,
  axis lines=middle,
  axis line style={-{Stealth}, popmuted},
  xlabel={$x$}, ylabel={$y$},
  xlabel style={at={(ticklabel* cs:1)}, anchor=west},
  ylabel style={at={(ticklabel* cs:1)}, anchor=south},
  xmin=-2.2, xmax=6.2, ymin=-4.2, ymax=2.2,
  xtick={-1,1,2,3,4,5}, ytick={-3,-2,-1,1},
  grid=both, grid style={popline!40, very thin},
  tick label style={font=\footnotesize},
  axis equal image
]
\addplot[popcurve, domain=0:360, samples=200] ({2+3*cos(x)},{-1+2*sin(x)});
\addplot[only marks, mark=*, mark size=2pt, poporange] coordinates {(2,-1)};
\node[poporange, anchor=south west, font=\small] at (axis cs:2.1,-0.9) {$\Omega(2\,;-1)$};
\addplot[only marks, mark=*, mark size=1.6pt, poppurple] coordinates {(4.236,-1) (-0.236,-1)};
\node[poppurple, anchor=south, font=\footnotesize] at (axis cs:4.236,-1) {$F_1$};
\node[poppurple, anchor=south, font=\footnotesize] at (axis cs:-0.236,-1) {$F_2$};
\addplot[popgreen, dashed, thick] coordinates {(2,-3) (2,1)};
\addplot[popgreen, dashed, thick] coordinates {(-1,-1) (5,-1)};
\end{axis}
\end{tikzpicture}
\end{center}
\legende{L'ellipse d'équation $4x^2+9y^2-16x+18y-11=0$, de centre $\Omega(2\,;-1)$, de demi-grand axe $a=3$ et demi-petit axe $b=2$ ; les axes de l'ellipse sont en pointillés verts, les foyers en violet.}
\end{popfigure}

\begin{attention}[Vérifications de fin de calcul]
Trois réflexes qui sauvent des points :
\begin{enumerate}
\item \textbf{Développe} ta forme réduite finale : tu dois retomber sur l'équation de départ. Ici, $4(x-2)^2+9(y+1)^2=36$ redonne bien $4x^2+9y^2-16x+18y-11=0$.
\item \textbf{Contrôle l'excentricité} : $0<e<1$ pour une ellipse, $e>1$ pour une hyperbole. Un résultat contraire signale une erreur.
\item \textbf{Place mentalement le centre} : les termes $-16x$ et $+18y$ « tirent » le centre vers $x=2$ et $y=-1$ ; si ton centre est ailleurs, reprends la complétion.
\end{enumerate}
\end{attention}

\courssub{Les autres cas : hyperbole, parabole et cas dégénérés}

\begin{exemplebox}[Une hyperbole reconnue]
Soit $(\mathcal{H})$ : $x^2-4y^2-2x+8y-7=0$. On complète :
\[ (x^2-2x)-4(y^2-2y)-7=0 \iff (x-1)^2-1-4\left[(y-1)^2-1\right]-7=0, \]
soit $(x-1)^2-4(y-1)^2=4$, puis
\[ \dfrac{(x-1)^2}{4}-(y-1)^2=1. \]
C'est une \textbf{hyperbole} de centre $\Omega(1\,;1)$, avec $a=2$, $b=1$, $c=\sqrt{a^2+b^2}=\sqrt{5}$, $e=\dfrac{\sqrt{5}}{2}>1$, d'axe focal horizontal. Si le second membre avait été $0$, on aurait obtenu $(x-1)=\pm 2(y-1)$ : deux droites sécantes en $\Omega$.
\end{exemplebox}

\begin{exemplebox}[Une parabole reconnue]
Soit $(\mathcal{P})$ : $y^2-4x-2y+9=0$. Ici $a=0$ : on complète seulement en $y$ :
\[ (y-1)^2-1-4x+9=0 \iff (y-1)^2=4x-8=4(x-2). \]
En posant $X=x-2$, $Y=y-1$ : $Y^2=4X$, parabole de \textbf{sommet} $S(2\,;1)$, de paramètre $p=2$, de foyer $F(3\,;1)$ et de directrice $x=1$.
\end{exemplebox}

\begin{exemplebox}[Cas dégénérés à repérer]
\begin{itemize}
\item $x^2+2y^2+2x+4y+5=0 \iff (x+1)^2+2(y+1)^2=-2$ : \textbf{ensemble vide}.
\item $x^2+2y^2+2x+4y+3=0 \iff (x+1)^2+2(y+1)^2=0$ : le \textbf{point} $(-1\,;-1)$.
\item $y^2-6y+5=0 \iff (y-1)(y-5)=0$ : \textbf{deux droites parallèles} d'équations $y=1$ et $y=5$.
\end{itemize}
\end{exemplebox}

\begin{aretenir}[L'essentiel de la section]
\begin{itemize}
\item Toute équation $ax^2+by^2+cx+dy+e=0$ se ramène, par \textbf{factorisation puis complétion des carrés}, à $aX^2+bY^2=k$ dans un repère translaté d'origine $\Omega\left(-\dfrac{c}{2a}\,;-\dfrac{d}{2b}\right)$.
\item La nature de l'ensemble se lit sur les \textbf{signes de $a$, $b$ et $k$} : $ab>0$ et $k$ de même signe $\Rightarrow$ ellipse (cercle si $a=b$) ; $ab<0$ et $k\neq 0$ $\Rightarrow$ hyperbole ; $ab=0$ $\Rightarrow$ parabole (ou cas dégénéré).
\item Un second membre nul ou de « mauvais » signe donne un \textbf{point}, \textbf{deux droites} ou l'\textbf{ensemble vide} : toujours examiner $k$ avant de conclure.
\item Les éléments caractéristiques (foyers, sommets, directrices) se calculent dans le repère réduit, puis se \textbf{traduient dans le repère initial} par $x=X+x_0$, $y=Y+y_0$.
\end{itemize}
\end{aretenir}

\begin{exercice}[Pour t'entraîner]
Déterminer la nature et les éléments caractéristiques (centre ou sommet, foyers, excentricité) des ensembles d'équations :
\begin{enumerate}
\item $9x^2+4y^2+18x-16y-11=0$ ;
\item $4x^2-y^2+8x+4y-4=0$ ;
\item $x^2+6x+4y+1=0$ ;
\item $x^2+y^2-2x+4y+5=0$.
\end{enumerate}
\end{exercice}

\begin{corrige}[Exercice]
\textbf{1.} $9(x+1)^2+4(y-2)^2=36 \iff \dfrac{(x+1)^2}{4}+\dfrac{(y-2)^2}{9}=1$ : ellipse de centre $\Omega(-1\,;2)$, $a=3$ (axe vertical), $b=2$, $c=\sqrt{5}$, $e=\dfrac{\sqrt{5}}{3}$.

\textbf{2.} $4(x+1)^2-(y-2)^2=4 \iff (x+1)^2-\dfrac{(y-2)^2}{4}=1$ : hyperbole de centre $\Omega(-1\,;2)$, $a=1$, $b=2$, $c=\sqrt{5}$, $e=\sqrt{5}$.

\textbf{3.} $(x+3)^2=-4(y-2)$ : parabole de sommet $S(-3\,;2)$, d'axe vertical, tournée vers le bas, $p=2$, foyer $F(-3\,;1)$, directrice $y=3$.

\textbf{4.} $(x-1)^2+(y+2)^2=0$ : l'ensemble est réduit au point $(1\,;-2)$ (cercle « de rayon nul »).
\end{corrige}

\coursec{Coniques, similitudes et situations concrètes}

Dans la section précédente, tu as appris à reconnaître une conique à partir de son équation du second degré : en complétant les carrés, tu ramènes toute équation $ax^2+by^2+cx+dy+e=0$ à une équation réduite dont tu identifies la nature. Nous allons maintenant répondre à deux questions naturelles. D'abord : que devient une conique lorsqu'on la transforme par une similitude directe ? Ensuite : où se cachent les coniques dans le monde qui nous entoure ? Tu vas découvrir que les antennes paraboliques qui captent Canal+ à Yaoundé, les orbites des satellites et les arches de ponts sont gouvernées par les mêmes mathématiques que celles que tu viens d'étudier.

\courssub{Image d'une conique par une similitude directe}

\textbf{Rappel sur les similitudes directes.} Une \cle{similitude directe} du plan est une transformation $s$ de rapport $k>0$ qui \textbf{multiplie toutes les distances par $k$} et \textbf{conserve les angles orientés}. Pour tous points $A$ et $B$ d'images $A'$ et $B'$ :
\[ A'B' = k \cdot AB. \]
Les similitudes directes de rapport $k$ sont les composées d'une homothétie de rapport $k$ et d'une isométrie directe (translation ou rotation). Les homothéties de rapport positif en sont le cas particulier le plus simple.

\textbf{Intuition.} Prenons une photo d'une ellipse et zoomons dessus : la forme ne change pas, seule la taille change. Une ellipse « aplatie » reste exactement aussi aplatie ; un cercle reste un cercle. Ce que le zoom conserve, c'est la \emph{forme} de la conique. Or nous avons un nombre qui mesure précisément cette forme : l'\cle{excentricité} $e$. Le théorème suivant rend cette intuition rigoureuse.

\begin{propriete}[Image d'une conique par une similitude directe]
Soit $\mathcal{C}$ une conique de foyer $F$, de directrice $\mathcal{D}$ et d'excentricité $e$, et soit $s$ une similitude directe de rapport $k>0$. Alors l'image $\mathcal{C}' = s(\mathcal{C})$ est une conique :
\begin{itemize}
\item de foyer $F' = s(F)$ ;
\item de directrice $\mathcal{D}' = s(\mathcal{D})$ ;
\item de \textbf{même excentricité} $e$, donc de \textbf{même nature} (parabole, ellipse ou hyperbole).
\end{itemize}
Ses dimensions (distances focales, axes, paramètre) sont multipliées par $k$.
\end{propriete}

\begin{experience}[Démonstration guidée du théorème]
Soit $\mathcal{C}$ la conique de foyer $F$, de directrice $\mathcal{D}$ et d'excentricité $e$. Notons $s$ une similitude directe de rapport $k$, et pour tout point $P$, notons $P' = s(P)$.

\textbf{Étape 1 : traduire l'appartenance à la conique.} Par définition monofocale :
\[ M \in \mathcal{C} \iff MF = e \cdot d(M, \mathcal{D}). \]

\textbf{Étape 2 : utiliser la conservation des distances à une droite.} Notons $H$ le projeté orthogonal de $M$ sur $\mathcal{D}$, de sorte que $d(M,\mathcal{D}) = MH$. Comme $s$ conserve les angles, elle conserve l'orthogonalité : l'image $H'$ de $H$ est le projeté orthogonal de $M'$ sur la droite $\mathcal{D}' = s(\mathcal{D})$. Par conséquent :
\[ d(M', \mathcal{D}') = M'H' = k \cdot MH = k \cdot d(M, \mathcal{D}). \]

\textbf{Étape 3 : conclure.} Supposons $M \in \mathcal{C}$. Alors :
\begin{align*}
M'F' &= k \cdot MF & &\text{(la similitude multiplie les distances par } k\text{)}\\
&= k \cdot e \cdot d(M, \mathcal{D}) & &\text{(car } M \in \mathcal{C}\text{)}\\
&= e \cdot \big(k \cdot d(M, \mathcal{D})\big) & &\\
&= e \cdot d(M', \mathcal{D}') & &\text{(étape 2).}
\end{align*}
Donc $M'$ appartient à la conique de foyer $F'$, de directrice $\mathcal{D}'$ et d'excentricité $e$. Réciproquement, en appliquant le même raisonnement à la similitude réciproque $s^{-1}$ (de rapport $1/k$), tout point de cette conique image provient d'un point de $\mathcal{C}$. L'image $\mathcal{C}'$ est donc exactement la conique de foyer $F'$, de directrice $\mathcal{D}'$ et d'excentricité $e$. \hfill$\blacksquare$
\end{experience}

\begin{aretenir}[L'excentricité, invariant des similitudes]
L'\cle{excentricité} est un \textbf{invariant} des similitudes : elle caractérise la \emph{forme} de la conique, indépendamment de sa taille et de sa position.
\begin{itemize}
\item $e = 1$ : l'image d'une \textbf{parabole} est une \textbf{parabole} ;
\item $0 < e < 1$ : l'image d'une \textbf{ellipse} est une \textbf{ellipse} de même excentricité (un cercle, cas particulier d'ellipse, a pour image un cercle) ;
\item $e > 1$ : l'image d'une \textbf{hyperbole} est une \textbf{hyperbole} de même excentricité.
\end{itemize}
Deux coniques de même nature et de même excentricité sont donc \textbf{semblables} : on passe de l'une à l'autre par une similitude. En particulier, \emph{toutes les paraboles sont semblables entre elles} !
\end{aretenir}

\begin{attention}[L'excentricité est conservée, pas les dimensions]
Ne confonds pas « même forme » et « mêmes dimensions ». Une similitude de rapport $k$ conserve $e$, mais \textbf{multiplie toutes les longueurs par $k$} : distance focale $c$, demi-axes $a$ et $b$, paramètre $p$ de la parabole. Par exemple, l'image d'une ellipse de demi-grand axe $a=5$ par une homothétie de rapport $2$ a un demi-grand axe $a'=10$, mais la même excentricité $e = c/a = c'/a'$.
\end{attention}

\begin{exemplebox}[Image de la parabole $y^2 = 4x$ par l'homothétie de centre $O$ et de rapport $3$]
La parabole $\mathcal{P}$ d'équation $y^2 = 4x$ a pour paramètre $p = 2$, pour foyer $F(1\,;\,0)$ et pour directrice $\mathcal{D}$ d'équation $x = -1$.

L'homothétie $h$ de centre $O$ et de rapport $3$ envoie tout point $M(x\,;\,y)$ sur $M'(3x\,;\,3y)$. Donc $x = \dfrac{x'}{3}$ et $y = \dfrac{y'}{3}$. En substituant dans l'équation :
\[ \left(\frac{y'}{3}\right)^2 = 4 \cdot \frac{x'}{3} \iff \frac{y'^2}{9} = \frac{4x'}{3} \iff y'^2 = 12x'. \]
L'image est la parabole d'équation $y^2 = 12x$, de paramètre $p' = 6 = 3p$, de foyer $F'(3\,;\,0) = h(F)$ et de directrice $x = -3$, image de $\mathcal{D}$. L'excentricité reste $e = 1$, comme prévu.
\end{exemplebox}

\begin{popfigure}
\begin{center}
\begin{tikzpicture}[scale=0.55, >=Stealth]
% axes
\draw[->, popmuted] (-8.5,0) -- (8.5,0) node[below right] {$x$};
\draw[->, popmuted] (0,-5.5) -- (0,5.5) node[above left] {$y$};
% ellipse initiale
\draw[popblue, thick] (0,0) ellipse (3 and 2);
% ellipse image (homothétie rapport 2)
\draw[poporange, thick] (0,0) ellipse (6 and 4);
% foyers initiaux : c = sqrt(9-4) = sqrt5 ≈ 2,24
\fill[popblue] (2.24,0) circle (2.5pt) node[below left=1pt, font=\small] {$F_1$};
\fill[popblue] (-2.24,0) circle (2.5pt) node[below right=1pt, font=\small] {$F_2$};
% foyers images
\fill[poporange] (4.47,0) circle (2.5pt) node[below left=1pt, font=\small] {$F'_1$};
\fill[poporange] (-4.47,0) circle (2.5pt) node[below right=1pt, font=\small] {$F'_2$};
% directrices : a²/c = 9/2,24 ≈ 4,03 et 8,05
\draw[popblue, dashed] (4.03,-2.6) -- (4.03,2.6) node[above, font=\small] {$\mathcal{D}$};
\draw[poporange, dashed] (8.05,-5) -- (8.05,5) node[above, font=\small] {$\mathcal{D}'$};
% centre
\fill[popdark] (0,0) circle (2pt) node[below left, font=\small] {$O$};
% légendes courbes
\node[popblue, font=\small] at (-2.1,2.35) {$\mathcal{E}$};
\node[poporange, font=\small] at (-4.2,4.6) {$\mathcal{E}' = h(\mathcal{E})$};
\end{tikzpicture}
\end{center}
\legende{L'ellipse $\mathcal{E}$ de demi-axes $a=3$, $b=2$ (excentricité $e=\frac{\sqrt{5}}{3}$) et son image $\mathcal{E}'$ par l'homothétie de centre $O$ et de rapport $2$ : demi-axes $6$ et $4$, foyers et directrices doublés en distance, même excentricité.}
\end{popfigure}

\courssub{Les coniques dans la nature et la technologie}

\textbf{Pourquoi les coniques sont-elles partout ?} Parce qu'elles possèdent des \cle{propriétés focales} remarquables, héritées de leur définition et des lois de la physique (réflexion de la lumière et des ondes, gravitation). En voici les plus spectaculaires.

\courssub{L'antenne parabolique : la propriété focale de la parabole}

\begin{propriete}[Propriété focale de la parabole — admise]
Tout rayon parallèle à l'axe d'une parabole se réfléchit sur celle-ci en passant par le \textbf{foyer}. Réciproquement, tout rayon issu du foyer se réfléchit parallèlement à l'axe.
\end{propriete}

C'est exactement le principe de l'\textbf{antenne parabolique} : les ondes émises par le satellite arrivent en faisceau parallèle à l'axe de la parabole, frappent la surface réfléchissante, et convergent toutes au foyer, où l'on place le récepteur (la tête LNB). Ainsi, même un signal très faible, capté sur toute la surface de la parabole, est concentré en un seul point.

\begin{popfigure}
\begin{center}
\begin{tikzpicture}[scale=0.9, >=Stealth]
% parabole x = y²/4 (p=2, foyer (1,0)), tracée pour y in [-3,3]
\draw[popblue, very thick, domain=-3:3, samples=100] plot ({0.25*\x*\x}, {\x});
% foyer
\fill[poporange] (1,0) circle (3pt) node[below=2pt, font=\small] {$F$ (récepteur)};
% rayons parallèles entrants
\draw[->, popgreen, thick] (8,2.25) -- (1.27,2.25);
\draw[->, popgreen, thick] (8,0) -- (0,0);
\draw[->, popgreen, thick] (8,-2.25) -- (1.27,-2.25);
% rayons réfléchis vers le foyer
\draw[->, poppink, thick] (1.27,2.25) -- (1.03,0.09);
\draw[->, poppink, thick] (1.27,-2.25) -- (1.03,-0.09);
% axe
\draw[popmuted, dashed] (-0.5,0) -- (8.3,0) node[below, font=\small] {axe};
\node[popgreen, font=\small, align=center] at (5.2,2.9) {ondes parallèles\\du satellite};
\node[popblue, font=\small] at (2.6,-2.6) {parabole};
\end{tikzpicture}
\end{center}
\legende{Antenne parabolique : les ondes parallèles à l'axe se réfléchissent sur la surface parabolique et convergent toutes au foyer $F$, où est placé le récepteur.}
\end{popfigure}

\begin{exoresolu}[Où placer le récepteur d'une antenne ?]
Une antenne parabolique a la forme d'un paraboloïde obtenu par rotation d'une parabole autour de son axe. Sa section méridienne est une parabole d'équation $y^2 = 0{,}8\,x$ (unité : le mètre). À quelle distance du sommet de l'antenne faut-il placer le récepteur ?
\tcblower
L'équation réduite d'une parabole d'axe $(Ox)$ est $y^2 = 2px$, où $p$ est le paramètre et le foyer est $F\left(\dfrac{p}{2}\,;\,0\right)$.

Ici $2p = 0{,}8$, donc $p = 0{,}4$ et $\dfrac{p}{2} = 0{,}2$.

Le foyer est $F(0{,}2\,;\,0)$ : le récepteur doit être placé sur l'axe de l'antenne, à $\SI{0,20}{m}$ (soit \SI{20}{cm}) du sommet, à l'intérieur de la coupole.
\end{exoresolu}

\courssub{Les orbites des planètes et des satellites : des ellipses}

\begin{savaistu}[La première loi de Kepler (1609)]
En analysant pendant des années les observations de l'astronome danois Tycho Brahe, Johannes Kepler énonça en 1609 sa première loi : \emph{chaque planète décrit une ellipse dont le Soleil occupe l'un des foyers}. C'est l'une des plus belles victoires des mathématiques sur l'astronomie : des siècles de modèles compliqués à base de cercles s'effondraient devant une simple ellipse ! Plus tard, Newton montra que cette loi découle de la gravitation universelle : tout corps en orbite autour d'un astre suit une conique dont l'astre est un foyer — ellipse si l'orbite est fermée (planètes, satellites), parabole ou hyperbole si le corps s'échappe (certaines comètes). L'excentricité de l'orbite de la Terre vaut environ $0{,}017$ : elle est presque circulaire. Celle de la comète de Halley vaut environ $0{,}967$ : une ellipse très allongée. Les satellites de télécommunication qui couvrent le Cameroun suivent eux aussi des orbites elliptiques (ou circulaires, cas limite $e=0$) dont la Terre est un foyer.
\end{savaistu}

\courssub{Le pont à arche parabolique : une modélisation complète}

Les arches de nombreux ponts ont une forme parabolique : cette courbe répartit remarquablement bien les efforts de compression. Modéliser une telle arche est un problème classique du Baccalauréat.

\begin{popfigure}
\begin{center}
\begin{tikzpicture}[scale=0.14, >=Stealth]
% axes
\draw[->, popmuted] (-24,0) -- (24,0) node[below right, font=\small] {$x$ (m)};
\draw[->, popmuted] (0,-2) -- (0,14) node[above, font=\small] {$y$ (m)};
% arche y = 10 - x²/40
\draw[popblue, very thick, domain=-20:20, samples=100] plot (\x, {10 - 0.025*\x*\x});
% tablier
\draw[popdark, thick] (-24,10) -- (24,10);
\fill[pattern=north east lines, pattern color=popmuted] (-24,10) rectangle (24,10.8);
% cotes
\draw[<->, poporange, thick] (-20,-3) -- (20,-3) node[midway, below, font=\small] {portée : $40$ m};
\draw[<->, poporange, thick] (22,0) -- (22,10) node[midway, right, font=\small] {flèche : $10$ m};
% points remarquables
\fill[popdark] (0,10) circle (8pt) node[above left=2pt, font=\small] {$S(0\,;\,10)$};
\fill[popdark] (-20,0) circle (8pt) node[below left, font=\small] {$A(-20\,;\,0)$};
\fill[popdark] (20,0) circle (8pt) node[below right, font=\small] {$B(20\,;\,0)$};
\node[popdark, font=\small] at (-14,11.8) {tablier};
\end{tikzpicture}
\end{center}
\legende{Arche parabolique d'un pont : portée $AB = \SI{40}{m}$, flèche (hauteur au sommet) $\SI{10}{m}$. L'arche est la parabole d'équation $y = 10 - \dfrac{x^2}{40}$ dans le repère indiqué.}
\end{popfigure}

\begin{methode}[Modéliser une arche parabolique]
Pour déterminer l'équation d'une arche parabolique à partir de ses dimensions :
\begin{enumerate}
\item \textbf{Choisir un repère adapté} : origine au milieu de la portée, axe $(Oy)$ vertical passant par le sommet de l'arche. L'arche a alors une équation de la forme $y = ax^2 + c$ (parabole de sommet sur l'axe des ordonnées).
\item \textbf{Traduire les données} : le sommet $S(0\,;\,c)$ donne la flèche $c$ ; les pieds de l'arche $A$ et $B$, symétriques, donnent la portée.
\item \textbf{Déterminer $a$} en écrivant que $A$ (ou $B$) appartient à la parabole.
\item \textbf{Conclure} et vérifier la cohérence des signes (ici $a < 0$ car la parabole est tournée vers le bas).
\end{enumerate}
\end{methode}

\begin{exoresolu}[Équation de l'arche d'un pont]
Un pont a une arche parabolique de portée $\SI{40}{m}$ et de flèche $\SI{10}{m}$ (voir la figure ci-dessus). Dans le repère où l'origine est le milieu de la portée et $(Oy)$ la verticale du sommet, déterminer l'équation de l'arche, puis la hauteur de l'arche à $\SI{10}{m}$ du centre.
\tcblower
\textbf{Équation de l'arche.} L'arche est une parabole d'axe $(Oy)$, tournée vers le bas : son équation est de la forme $y = ax^2 + c$ avec $a < 0$.

Le sommet est $S(0\,;\,10)$, donc $c = 10$. La portée vaut $\SI{40}{m}$, donc les pieds de l'arche sont $A(-20\,;\,0)$ et $B(20\,;\,0)$. Le point $B$ appartient à la parabole :
\[ 0 = a \times 20^2 + 10 \iff 400a = -10 \iff a = -\frac{1}{40}. \]
L'équation de l'arche est donc :
\[ y = 10 - \frac{x^2}{40}, \qquad x \in [-20\,;\,20]. \]

\textbf{Hauteur à $\SI{10}{m}$ du centre.} Pour $x = 10$ :
\[ y = 10 - \frac{100}{40} = 10 - 2{,}5 = 7{,}5. \]
À $\SI{10}{m}$ du centre, l'arche se trouve à $\SI{7,5}{m}$ au-dessus de la base.
\end{exoresolu}

\begin{attention}[Le choix du repère change l'équation, pas la parabole]
Si tu places l'origine du repère au pied $A$ de l'arche plutôt qu'au centre, tu obtiens une équation différente (par exemple $y = -\dfrac{1}{40}(x-20)^2 + 10$), mais c'est \textbf{la même parabole}. Au Baccalauréat, lis attentivement le repère imposé par l'énoncé avant d'écrire la moindre équation, et vérifie toujours que les points remarquables (sommet, pieds) satisfont ton équation finale.
\end{attention}

\courssub{Sections planes et ombres : le retour aux sources}

Pour terminer, revenons à l'origine géométrique des coniques. L'\textbf{ombre} d'un objet éclairé par une source ponctuelle est une \textbf{projection conique} : les rayons lumineux forment un cône, et l'ombre sur un plan est la section de ce cône par ce plan. Ainsi :
\begin{itemize}
\item l'ombre d'un ballon (sphère) ou d'un cerceau (cercle) sur un sol horizontal, avec le soleil au zénith, est un cercle ;
\item si le plan est incliné (ombre portée sur un talus, sur un escalier), le cercle devient une \textbf{ellipse} ;
\item en inclinant davantage le plan, l'ombre s'allonge jusqu'à devenir une \textbf{parabole} (plan parallèle à une génératrice du cône de lumière), puis une \textbf{hyperbole}.
\end{itemize}
C'est exactement la classification des sections planes d'un cône de révolution vue au début de la leçon : la nature de la conique dépend uniquement de l'inclinaison du plan par rapport au cône.

\begin{savaistu}[Les coniques, vingt-trois siècles d'histoire]
Les coniques ont été étudiées pour elles-mêmes par les Grecs : Ménæchme (vers 350 av. J.-C.), puis Apollonios de Perga, qui leur donna leurs noms — \emph{ellipse} (« il manque »), \emph{parabole} (« à côté, égal »), \emph{hyperbole} (« au-delà ») — en référence à des propriétés d'aires. Pendant près de deux mille ans, ces courbes restèrent de pures curiosités géométriques... jusqu'à ce que Kepler découvre que les planètes suivent des ellipses, et Galilée que les projectiles décrivent des paraboles. Aujourd'hui, les coniques sont partout : antennes, phares de voiture, ponts, orbites des satellites qui font fonctionner ton téléphone, et même certains fours solaires. Peu de chapitres des mathématiques relient aussi directement la géométrie pure au monde réel.
\end{savaistu}

\begin{exercice}[Pour t'entraîner]
\begin{enumerate}
\item Soit $\mathcal{E}$ l'ellipse d'équation réduite $\dfrac{x^2}{25} + \dfrac{y^2}{9} = 1$. Déterminer son excentricité, puis l'équation de son image par l'homothétie de centre $O$ et de rapport $\dfrac{1}{2}$. Vérifier que l'excentricité est conservée.
\item Une antenne parabolique a $\SI{1,2}{m}$ de diamètre et $\SI{0,3}{m}$ de profondeur. En choisissant un repère d'origine le sommet de la parabole, déterminer l'équation de sa section méridienne, puis la position du foyer où placer le récepteur.
\item L'orbite d'un satellite est une ellipse de demi-grand axe $a = \SI{7200}{km}$ et d'excentricité $e = 0{,}2$. La Terre occupe un foyer. Calculer les distances minimale et maximale du satellite au centre de la Terre.
\end{enumerate}
\end{exercice}

\begin{corrige}[Exercice]
\textbf{1.} On a $a = 5$, $b = 3$, donc $c = \sqrt{a^2 - b^2} = \sqrt{25-9} = 4$ et $e = \dfrac{c}{a} = \dfrac{4}{5} = 0{,}8$.

L'homothétie de rapport $\dfrac{1}{2}$ envoie $M(x\,;\,y)$ sur $M'\left(\dfrac{x}{2}\,;\,\dfrac{y}{2}\right)$, donc $x = 2x'$ et $y = 2y'$. En substituant :
\[ \frac{4x'^2}{25} + \frac{4y'^2}{9} = 1 \iff \frac{x'^2}{25/4} + \frac{y'^2}{9/4} = 1. \]
L'image est l'ellipse de demi-axes $a' = \dfrac{5}{2}$ et $b' = \dfrac{3}{2}$. Alors $c' = \sqrt{\dfrac{25}{4} - \dfrac{9}{4}} = 2$ et $e' = \dfrac{c'}{a'} = \dfrac{2}{5/2} = \dfrac{4}{5} = e$. L'excentricité est bien conservée.

\textbf{2.} Plaçons le sommet de la parabole à l'origine, l'axe $(Oy)$ vertical vers le haut. La parabole a une équation de la forme $y = ax^2$. Le bord de l'antenne correspond à $x = 0{,}6$ (demi-diamètre) et $y = 0{,}3$ (profondeur) :
\[ 0{,}3 = a \times 0{,}36 \iff a = \frac{0{,}3}{0{,}36} = \frac{5}{6}. \]
L'équation est $y = \dfrac{5}{6}x^2$, soit $x^2 = \dfrac{6}{5}y = 1{,}2\,y$. Or $x^2 = 2py$, donc $2p = 1{,}2$, $p = 0{,}6$, et le foyer est à $\dfrac{p}{2} = 0{,}3$ du sommet sur l'axe. Le récepteur se place à $\SI{0,30}{m}$ au-dessus du sommet (exactement au niveau du bord de l'antenne).

\textbf{3.} Pour une ellipse orbitale dont le centre de la Terre occupe un foyer, la distance au foyer varie entre $a - c$ (périgée) et $a + c$ (apogée), avec $c = ae = 7200 \times 0{,}2 = \SI{1440}{km}$.

Distance minimale : $a - c = 7200 - 1440 = \SI{5760}{km}$.

Distance maximale : $a + c = 7200 + 1440 = \SI{8640}{km}$.
\end{corrige}

\newpage

\coursec{Activité d'intégration}

\courssub{Objectifs de l'activité}

À l'issue de cette activité d'intégration, tu seras capable de :

\begin{itemize}
\item modéliser une situation concrète (antenne parabolique, arche de pont, ombre d'un panneau) par une conique et choisir un repère adapté ;
\item déterminer l'\cle{équation réduite} d'une parabole à partir de données géométriques et en déduire la position du \cle{foyer} et de la \cle{directrice} ;
\item déterminer l'équation de la \cle{tangente} à une conique en un point et l'interpréter concrètement ;
\item reconnaître et mettre sous \cle{forme réduite} une équation du type $ax^2+by^2+cx+dy+e=0$, puis déterminer centre, sommets, foyers et \cle{excentricité} ;
\item vérifier qu'une \cle{homothétie} (cas particulier de similitude directe) transforme une ellipse en une ellipse de même excentricité.
\end{itemize}

\courssub{Situation}

\textbf{Partie A — L'antenne parabolique de Ngaoundéré.}

Dans un quartier de Ngaoundéré, un technicien installe une antenne parabolique pour capter les signaux d'un satellite. L'antenne a la forme d'un paraboloïde de révolution : toute coupe passant par l'axe de l'antenne est une \cle{parabole}. La propriété remarquable de la parabole est que tous les rayons arrivant parallèlement à son axe se réfléchissent en un même point : le \cle{foyer}. C'est donc exactement là qu'il faut fixer le récepteur (la tête LNB), maintenu par un bras.

L'antenne a un \textbf{diamètre de \SI{1,2}{m}} et une \textbf{profondeur de \SI{20}{cm}} (distance entre le fond de la parabole et le plan du bord). On modélise la coupe méridienne de l'antenne dans un repère orthonormé $(O\,;\vec{i},\vec{j})$ dont l'origine $O$ est le sommet de la parabole (le fond de l'antenne), l'axe des abscisses étant l'axe de la parabole, dirigé vers l'ouverture. L'unité est le mètre. Le bord de l'antenne correspond alors aux points $A(0,2\,;\,0,6)$ et $A'(0,2\,;\,-0,6)$.

\textbf{Partie B — Le pont sur la Sanaga et l'ombre d'un panneau.}

Sur l'axe Yaoundé--Bafoussam, l'arche principale d'un pont franchissant un affluent de la Sanaga est modélisée par une \textbf{parabole de portée \SI{40}{m}} (distance entre les deux points d'appui au sol) et de \textbf{flèche \SI{10}{m}} (hauteur maximale de l'arche au-dessus du sol). On choisit un repère orthonormé dont l'origine est le milieu des points d'appui, l'axe des abscisses étant horizontal au niveau du sol, l'axe des ordonnées étant vertical ascendant.

À proximité du chantier, un panneau de signalisation \textbf{circulaire}, éclairé par le soleil rasant du matin, projette au sol une \textbf{ombre elliptique}. Dans un repère orthonormé du sol (unité : le mètre), cette ombre a pour équation :
\[ 4x^2+9y^2-16x+18y-11=0. \]
Le chef de chantier affirme qu'en doublant toutes les dimensions (homothétie de rapport $2$), l'ombre resterait une ellipse « de même forme ». Il s'agit de vérifier cette affirmation.

\courssub{Consignes}

\textbf{Partie A}

\begin{enumerate}
\item Justifier que, dans le repère choisi, la parabole admet une équation de la forme $y^2=2px$ avec $p>0$.
\item En utilisant les coordonnées du point $A$, déterminer la valeur de $p$ et l'équation réduite de la parabole.
\item En déduire les coordonnées du foyer $F$ et l'équation de la directrice. À quelle distance du fond de l'antenne le technicien doit-il fixer le récepteur ?
\item Déterminer l'équation de la tangente à la parabole au point $A(0,2\,;\,0,6)$, puis calculer son coefficient directeur. Interpréter : quel angle cette tangente fait-elle approximativement avec l'axe de l'antenne, et pourquoi cette information est-elle utile pour régler le bras de fixation ?
\end{enumerate}

\textbf{Partie B}

\begin{enumerate}
\item Déterminer l'équation de la parabole modélisant l'arche du pont dans le repère indiqué.
\item Montrer que l'équation $4x^2+9y^2-16x+18y-11=0$ est celle d'une ellipse, en la mettant sous forme réduite à l'aide de formes canoniques.
\item Déterminer le centre $\Omega$, les demi-axes $a$ et $b$, les sommets, les foyers et l'excentricité $e$ de cette ellipse.
\item Écrire l'équation de l'image de cette ellipse par l'homothétie $h$ de centre $O$ et de rapport $2$. Montrer que c'est encore une ellipse, déterminer son excentricité $e'$ et conclure quant à l'affirmation du chef de chantier.
\end{enumerate}

\courssub{Ressources à mobiliser}

\begin{aretenir}[Boîte à outils de l'activité]
\begin{itemize}
\item \textbf{Parabole} de sommet $O$, d'axe $(Ox)$, ouverte vers les $x$ positifs : équation réduite $y^2=2px$ avec $p>0$ ; foyer $F\left(\dfrac{p}{2}\,;\,0\right)$ ; directrice d'équation $x=-\dfrac{p}{2}$ ; pour tout point $M$ de la parabole, $MF=d(M,\mathcal{D})$ (excentricité $e=1$).
\item \textbf{Tangente} à la parabole $y^2=2px$ au point $M_0(x_0\,;\,y_0)$ : $y\,y_0=p(x+x_0)$.
\item \textbf{Ellipse} de centre $\Omega(\alpha\,;\,\beta)$, d'axe focal horizontal : $\dfrac{(x-\alpha)^2}{a^2}+\dfrac{(y-\beta)^2}{b^2}=1$ avec $a>b>0$ ; distance focale $c=\sqrt{a^2-b^2}$ ; foyers de coordonnées $(\alpha\pm c\,;\,\beta)$ ; excentricité $e=\dfrac{c}{a}\in\,]0\,;\,1[$.
\item \textbf{Forme canonique} : pour $a\neq 0$, $ax^2+cx=a\left(x+\dfrac{c}{2a}\right)^2-\dfrac{c^2}{4a}$.
\item \textbf{Similitude directe} : toute similitude transforme une conique en une conique de \emph{même nature} et de \emph{même excentricité} ; une homothétie de rapport $k$ multiplie les longueurs par $|k|$.
\end{itemize}
\end{aretenir}

\begin{attention}[Avant de te lancer]
\begin{itemize}
\item Dans $y^2=2px$, le paramètre $p$ est la \emph{distance foyer--directrice} ; le foyer est à l'abscisse $\dfrac{p}{2}$, pas $p$. C'est l'erreur la plus fréquente au Bac.
\item La formule de la tangente $yy_0=p(x+x_0)$ n'est valable que si $M_0(x_0\,;\,y_0)$ appartient effectivement à la parabole : vérifie-le toujours avant de l'appliquer.
\item Pour reconnaître une conique, la mise sous forme canonique doit être menée \emph{séparément} sur $x$ et sur $y$, en factorisant d'abord le coefficient du terme carré.
\end{itemize}
\end{attention}

\courssub{Démarche et corrigé commenté}

\begin{exoresolu}[Corrigé complet de l'activité]
\textbf{Partie A.} 1) Choisir le repère au sommet. 2) Calculer $p$. 3) Localiser le foyer. 4) Tangente en $A$.
\tcblower
\textbf{A.1.} La parabole a pour sommet $O$, origine du repère, et pour axe de symétrie l'axe des abscisses ; elle est de plus ouverte vers les $x$ positifs (vers l'ouverture de l'antenne). Son équation réduite est donc bien de la forme $y^2=2px$ avec $p>0$.

\textbf{A.2.} Le point $A(0,2\,;\,0,6)$ appartient à la parabole, donc ses coordonnées vérifient l'équation :
\[ (0,6)^2=2p\times 0,2 \iff 0,36=0,4\,p \iff p=\frac{0,36}{0,4}=0,9. \]
L'équation réduite de la parabole est $\boxed{y^2=1,8\,x}$.

\textbf{A.3.} Le foyer est $F\left(\dfrac{p}{2}\,;\,0\right)$, soit $F(0,45\,;\,0)$. La directrice a pour équation $x=-\dfrac{p}{2}$, c'est-à-dire $x=-0,45$. Le récepteur doit donc être fixé sur l'axe de l'antenne, à \textbf{\SI{45}{cm} du fond}, c'est-à-dire \SI{25}{cm} devant le plan du bord (car $45-20=25$). C'est la conséquence directe de la définition monofocale $MF=d(M,\mathcal{D})$ (ici $e=1$) : tout rayon parallèle à l'axe se réfléchit sur la parabole et converge en $F$.

\textbf{A.4.} Vérifions d'abord que $A$ est bien sur la parabole : $y_A^2=0,36$ et $1,8\,x_A=1,8\times 0,2=0,36$ : c'est bon. La tangente en $A(x_0\,;\,y_0)=(0,2\,;\,0,6)$ a pour équation $y\,y_0=p(x+x_0)$ :
\[ 0,6\,y=0,9(x+0,2) \iff 0,6\,y=0,9\,x+0,18 \iff y=1,5\,x+0,3. \]
Le coefficient directeur est $1,5$ : la tangente fait avec l'axe de l'antenne un angle $\theta$ tel que $\tan\theta=1,5$, soit $\theta\approx 56,3^{\circ}$. Le bras de fixation doit donc être orienté en tenant compte de cette inclinaison au bord de l'antenne pour que la tête LNB soit positionnée exactement en $F$.

\medskip
\textbf{Partie B.}

\textbf{B.1.} L'arche a pour sommet $S(0\,;\,10)$ (la flèche vaut \SI{10}{m}) et passe par les points d'appui $(-20\,;\,0)$ et $(20\,;\,0)$ (la portée vaut \SI{40}{m}). Par symétrie par rapport à l'axe des ordonnées, on cherche une équation de la forme $y=\alpha x^2+10$. La condition $0=400\alpha+10$ donne $\alpha=-\dfrac{10}{400}=-\dfrac{1}{40}$. D'où :
\[ \boxed{y=10-\frac{x^2}{40}} \qquad\text{soit encore } x^2=-40(y-10), \]
parabole d'axe vertical, de sommet $S$, de paramètre $p=20$ (car $2p=40$), ouverte vers le bas.

\textbf{B.2.} Regroupons les termes en $x$ et en $y$ et passons aux formes canoniques :
\begin{align*}
4x^2-16x &= 4\left(x^2-4x\right)=4\left[(x-2)^2-4\right]=4(x-2)^2-16,\\
9y^2+18y &= 9\left(y^2+2y\right)=9\left[(y+1)^2-1\right]=9(y+1)^2-9.
\end{align*}
L'équation devient :
\[ 4(x-2)^2+9(y+1)^2-16-9-11=0 \iff 4(x-2)^2+9(y+1)^2=36. \]
En divisant les deux membres par $36$ :
\[ \boxed{\frac{(x-2)^2}{9}+\frac{(y+1)^2}{4}=1} \]
Les coefficients de $(x-2)^2$ et $(y+1)^2$ sont de même signe et le second membre est strictement positif : c'est bien l'équation réduite d'une \textbf{ellipse}.

\textbf{B.3.} Centre : $\Omega(2\,;\,-1)$. Comme $9>4$, le grand axe est horizontal : demi-grand axe $a=\sqrt{9}=3$ ; demi-petit axe $b=\sqrt{4}=2$. Distance focale :
\[ c=\sqrt{a^2-b^2}=\sqrt{9-4}=\sqrt{5}. \]
Sommets du grand axe : $(2\pm 3\,;\,-1)$, soit $(5\,;\,-1)$ et $(-1\,;\,-1)$ ; sommets du petit axe : $(2\,;\,-1\pm 2)$, soit $(2\,;\,1)$ et $(2\,;\,-3)$. Foyers :
\[ F_1\left(2+\sqrt{5}\,;\,-1\right), \qquad F_2\left(2-\sqrt{5}\,;\,-1\right). \]
Excentricité : $e=\dfrac{c}{a}=\dfrac{\sqrt{5}}{3}\approx 0,745$. Comme $0<e<1$, le résultat est cohérent avec une ellipse.

\textbf{B.4.} L'homothétie $h$ de centre $O$ et de rapport $2$ envoie $M(x\,;\,y)$ sur $M'(2x\,;\,2y)$. Un point $M'(X\,;\,Y)$ appartient à l'image si et seulement si son antécédent $M\left(\dfrac{X}{2}\,;\,\dfrac{Y}{2}\right)$ appartient à l'ellipse :
\[ \frac{\left(\frac{X}{2}-2\right)^2}{9}+\frac{\left(\frac{Y}{2}+1\right)^2}{4}=1
\iff \frac{(X-4)^2}{36}+\frac{(Y+2)^2}{16}=1, \]
car $\left(\dfrac{X}{2}-2\right)^2=\dfrac{(X-4)^2}{4}$ et de même pour l'autre terme. C'est une ellipse de centre $\Omega'(4\,;\,-2)$, de demi-axes $a'=\sqrt{36}=6=2a$ et $b'=\sqrt{16}=4=2b$. Sa distance focale est :
\[ c'=\sqrt{a'^2-b'^2}=\sqrt{36-16}=\sqrt{20}=2\sqrt{5}=2c, \]
d'où :
\[ e'=\frac{c'}{a'}=\frac{2\sqrt{5}}{6}=\frac{\sqrt{5}}{3}=e. \]
\textbf{Conclusion :} l'image est bien une ellipse de \emph{même excentricité} : le chef de chantier a raison, l'ombre agrandie conserve exactement la même « forme ». C'est un cas particulier du théorème du cours : toute similitude directe conserve la nature et l'excentricité d'une conique.
\end{exoresolu}

\begin{popfigure}
\begin{center}
\begin{tikzpicture}[scale=0.85]
% Ellipse initiale
\draw[popblue, thick] (2,-1) ellipse (3 and 2);
\fill[popblue] (2,-1) circle (1.5pt) node[below left] {\small $\Omega$};
\fill[popink] (4.236,-1) circle (1.5pt) node[above right] {\small $F_1$};
\fill[popink] (-0.236,-1) circle (1.5pt) node[above left] {\small $F_2$};
% Ellipse image
\draw[poporange, thick, dashed] (4,-2) ellipse (6 and 4);
\fill[poporange] (4,-2) circle (1.5pt) node[below right] {\small $\Omega'$};
% Centre O
\fill[popdark] (0,0) circle (1.5pt) node[above left] {\small $O$};
\draw[popmuted, ->] (0,0) -- (2,-1);
\draw[popmuted, ->] (0,0) -- (4,-2);
\end{tikzpicture}
\end{center}
\legende{L'ellipse de l'ombre (bleu) et son image par l'homothétie de centre $O$ et de rapport $2$ (orange, pointillés) : les dimensions doublent, l'excentricité est conservée.}
\end{popfigure}

\courssub{Bilan de l'activité}

Cette activité a permis de mettre en évidence plusieurs idées essentielles de la leçon :

\begin{itemize}
\item \textbf{Le choix du repère est stratégique.} Placer l'origine au sommet de la parabole (Partie A) ou au milieu des appuis (Partie B.1) conduit directement à une équation réduite, sans calcul superflu. C'est une compétence de modélisation à part entière.
\item \textbf{La définition monofocale a un sens physique.} La position du foyer à \SI{45}{cm} du fond de l'antenne n'est pas un simple nombre : c'est l'endroit exact où l'énergie du signal se concentre, parce que $MF=d(M,\mathcal{D})$ avec $e=1$.
\item \textbf{La tangente relie le calcul à la géométrie.} L'équation $yy_0=p(x+x_0)$ donne immédiatement l'inclinaison du bord de l'antenne, information directement exploitable par le technicien.
\item \textbf{La reconnaissance d'une équation du second degré} repose sur la mise sous forme canonique : signes des coefficients de $x^2$ et $y^2$ (ici de même signe : ellipse), puis lecture directe du centre, des axes, des foyers et de l'excentricité.
\item \textbf{Les similitudes préservent la « forme » des coniques.} L'homothétie de rapport $2$ double $a$, $b$ et $c$, donc laisse invariant le quotient $e=\dfrac{c}{a}$ : l'excentricité est un invariant de forme, exactement comme les angles pour les figures usuelles.
\end{itemize}

Tu as ainsi mobilisé l'ensemble des savoir-faire de la leçon dans deux situations authentiques : identifier la nature d'une conique, déterminer son équation réduite et ses éléments caractéristiques, calculer une tangente et étudier l'image d'une conique par une similitude. Ces compétences sont exactement celles exigées au Baccalauréat dans les problèmes de modélisation géométrique.

\newpage

\coursec{Méthodes et exercices résolus}

\coursec{Leçon M15 : Coniques}

\coursec{Coniques comme sections d'un cône}

\begin{definition}[Cône de révolution et sections planes]
Soit $(\mathcal{C})$ un cône de révolution de sommet $S$, d'axe $(\Delta)$ et d'angle au sommet $2\alpha$ ($0 < \alpha < \frac{\pi}{2}$). On coupe ce cône par un plan $(\mathcal{P})$ ne passant pas par $S$. Selon l'angle $\beta$ que forme $(\mathcal{P})$ avec l'axe $(\Delta)$, on obtient quatre types de courbes planes appelées \cle{coniques} :
\begin{itemize}
\item \textbf{Cercle} : $\beta = \frac{\pi}{2}$ (plan perpendiculaire à l'axe).
\item \textbf{Ellipse} : $\alpha < \beta < \frac{\pi}{2}$ (plan incliné coupant une seule nappe).
\item \textbf{Parabole} : $\beta = \alpha$ (plan parallèle à une génératrice du cône).
\item \textbf{Hyperbole} : $0 \le \beta < \alpha$ (plan coupant les deux nappes).
\end{itemize}
\end{definition}

\begin{popfigure}
\begin{tikzpicture}[scale=1.2, >=Stealth]
% Cone lines
\draw[popblue, thick] (-2, -3) -- (0, 0) -- (2, -3);
\draw[popblue, thick] (-2, 3) -- (0, 0) -- (2, 3);
% Axis
\draw[popmuted, dashed] (0, -3.5) -- (0, 3.5) node[above] {$(\Delta)$};
% Vertex
\fill[popblue] (0,0) circle (1.5pt) node[right] {$S$};
% Angle alpha
\draw[poporange, thick] (0.3, 0) arc (0:30:0.3) node[midway, right] {$\alpha$};
% Cutting planes (lines in cross-section)
% Circle (horizontal)
\draw[popgreen, thick] (-1.5, -1.5) -- (1.5, -1.5) node[right] {\textcolor{popgreen}{Cercle ($\beta=90^\circ$)}};
% Ellipse (slanted)
\draw[poporange, thick] (-2.5, 0.5) -- (2.5, -0.5) node[right] {\textcolor{poporange}{Ellipse ($\alpha<\beta<90^\circ$)}};
% Parabola (parallel to generator)
\draw[poppurple, thick] (-2.5, 1.5) -- (2.5, -1.5) node[right] {\textcolor{poppurple}{Parabole ($\beta=\alpha$)}};
% Hyperbola (vertical/steep)
\draw[poppink, thick] (-1, 2.5) -- (-1, -2.5) node[below] {\textcolor{poppink}{Hyperbole ($\beta<\alpha$)}};
% Generator line for reference
\draw[popmuted, thin] (0,0) -- (2, 2);
\draw[popmuted, thin] (0,0) -- (2, -2);
\draw[poporange, thick] (0.5, 0) arc (0:45:0.5) node[midway, above right] {$\beta$};
\end{tikzpicture}
\legende{Sections planes d'un cône de révolution : les quatre types de coniques.}
\end{popfigure}

\begin{aretenir}[Lien avec la définition monofocale]
Toute conique (sauf le cercle) admet une \cle{définition monofocale} (ou bifocale équivalente) qui ne fait pas intervenir le cône : c'est l'ensemble des points $M$ du plan tels que la distance à un point fixe $F$ (foyer) est proportionnelle à la distance à une droite fixe $D$ (directrice), $F \notin D$. Le coefficient de proportionnalité $e > 0$ est l'\cle{excentricité}. Cette définition unifiée permet l'étude algébrique et métrique des coniques.
\end{aretenir}

\coursec{Définition monofocale et excentricité}

\begin{methode}[Reconnaître une conique grâce à son excentricité]
Une conique de foyer $F$, de directrice $D$ ($F \notin D$) et d'excentricité $e > 0$ est l'ensemble des points $M$ tels que $MF = e \cdot d(M, D)$. Pour identifier la nature de la conique :
\begin{enumerate}
\item Repérer dans l'énoncé le foyer $F$, la directrice $D$ et l'excentricité $e$ (ou les données qui permettent de calculer $e$).
\item Comparer $e$ à $1$ :
\begin{itemize}
\item si $e = 1$ : c'est une \cle{parabole} ;
\item si $0 < e < 1$ : c'est une \cle{ellipse} ;
\item si $e > 1$ : c'est une \cle{hyperbole}.
\end{itemize}
\item Si l'énoncé donne une équation réduite, retrouver $e$ par les formules :
\begin{itemize}
\item ellipse $\dfrac{x^2}{a^2} + \dfrac{y^2}{b^2} = 1$ ($a > b > 0$, grand axe horizontal) : $c = \sqrt{a^2 - b^2}$ et $e = \dfrac{c}{a}$ ;
\item hyperbole $\dfrac{x^2}{a^2} - \dfrac{y^2}{b^2} = 1$ : $c = \sqrt{a^2 + b^2}$ et $e = \dfrac{c}{a}$.
\end{itemize}
\item Conclure par une phrase complète nommant la conique et justifiant par la valeur de $e$.
\end{enumerate}
\textbf{Réflexe :} pour une ellipse, $c^2 = a^2 - b^2$ (soustraction) ; pour une hyperbole, $c^2 = a^2 + b^2$ (addition). Ne jamais confondre ces deux relations. Pour l'ellipse, $e<1$ car $c<a$ ; pour l'hyperbole, $e>1$ car $c>a$.
\end{methode}

\begin{popfigure}
\begin{tikzpicture}[scale=0.9, >=Stealth, domain=-3:3]
% Axes
\draw[->] (-3.5,0) -- (3.5,0) node[right] {$x$};
\draw[->] (0,-2.5) -- (0,2.5) node[above] {$y$};
% Directrix
\draw[popblue, thick] (2, -2.5) -- (2, 2.5) node[above] {$D$};
\node[popblue] at (2, -2.8) {$x = \frac{a^2}{c}$};
% Focus
\fill[popred] (1,0) circle (2pt) node[below] {$F(c;0)$};
\fill[popred] (-1,0) circle (2pt) node[below] {$F'(-c;0)$};
% Ellipse
\draw[popgreen, thick, samples=200] plot (\x, {1.5*sqrt(max(0,1-\x*\x/4))});
\draw[popgreen, thick, samples=200] plot (\x, {-1.5*sqrt(max(0,1-\x*\x/4))});
\node[popgreen] at (3, 1.8) {Ellipse ($e<1$)};
% Point M
\fill[popdark] (1.5, {1.5*sqrt(1-2.25/4)}) circle (1.5pt) node[above right] {$M$};
\draw[popmuted, dashed] (1.5, {1.5*sqrt(max(0,1-2.25/4))}) -- (1,0) node[midway, above] {$MF$};
\draw[popmuted, dashed] (1.5, {1.5*sqrt(max(0,1-2.25/4))}) -- (2, {1.5*sqrt(max(0,1-2.25/4))}) node[midway, above] {$d(M,D)$};
\end{tikzpicture}
\legende{Définition monofocale de l'ellipse : $MF = e \cdot d(M,D)$ avec $e=c/a < 1$.}
\end{popfigure}

\begin{exoresolu}[Type d'une conique donnée par son excentricité]
\textbf{Énoncé.} Dans le plan muni d'un repère orthonormé, on considère le point $F(3\,;\,0)$ et la droite $D$ d'équation $x = \dfrac{25}{3}$. Soit $(\mathcal{C})$ l'ensemble des points $M$ du plan tels que $MF = \dfrac{3}{5}\, d(M, D)$.
\begin{enumerate}
\item Déterminer la nature de la conique $(\mathcal{C})$.
\item L'orbite d'un satellite autour de la Terre est modélisée par la conique de foyer $F'(4\,;\,0)$, de directrice $D' : x = 9$ et d'excentricité $e'$. Sachant que le point $A(5\,;\,0)$ appartient à cette orbite, calculer $e'$ et conclure.
\end{enumerate}
\tcblower
\textbf{Solution détaillée.}
\begin{enumerate}
\item L'énoncé donne directement la définition monofocale : $MF = e \cdot d(M, D)$ avec $e = \dfrac{3}{5}$.

Or $0 < \dfrac{3}{5} < 1$. D'après la classification des coniques par l'excentricité, $(\mathcal{C})$ est une \textbf{ellipse} de foyer $F(3\,;\,0)$, de directrice $D : x = \dfrac{25}{3}$ et d'excentricité $e = \dfrac{3}{5}$.

\item Le point $A(5\,;\,0)$ appartient à la conique, donc il vérifie la relation $AF' = e' \cdot d(A, D')$.

Calculons chaque distance :
\[ AF' = |5 - 4| = 1 \qquad \text{et} \qquad d(A, D') = \left| 5 - 9 \right| = 4. \]
On en déduit :
\[ e' = \frac{AF'}{d(A, D')} = \frac{1}{4}. \]
Comme $0 < e' = \dfrac{1}{4} < 1$, l'orbite du satellite est une \textbf{ellipse}, ce qui est cohérent avec les lois de Kepler : les orbites des satellites sont des ellipses dont la Terre occupe un foyer.
\end{enumerate}
\textbf{Conclusion :} la conique $(\mathcal{C})$ est une ellipse d'excentricité $\dfrac{3}{5}$, et l'orbite du satellite est une ellipse d'excentricité $\dfrac{1}{4}$.
\end{exoresolu}

\courssub{Construction point par point d'une conique}

\begin{methode}[Construction point par point à partir de la définition monofocale]
On connaît le foyer $F$, la directrice $D$ et l'excentricité $e$. Pour obtenir des points de la conique :
\begin{enumerate}
\item Tracer la directrice $D$, placer le foyer $F$, et tracer la droite $(\Delta)$ perpendiculaire à $D$ passant par $F$ : c'est l'\cle{axe focal}, axe de symétrie de la conique.
\item Choisir une droite $(d)$ parallèle à $D$. Tout point $M$ de $(d)$ vérifie $d(M, D) = $ distance entre les deux droites, que l'on note $h$.
\item Reporter la longueur $e \cdot h$ (à la règle ou au compas) : les points $M$ de la conique situés sur $(d)$ sont les intersections de $(d)$ avec le cercle de centre $F$ et de rayon $e \cdot h$.
\item Renouveler avec plusieurs droites parallèles à $D$, puis relier les points obtenus par une courbe régulière.
\item Exploiter la symétrie : tout point $M$ obtenu a son symétrique par rapport à l'axe focal qui appartient aussi à la conique.
\end{enumerate}
\textbf{Réflexe :} les sommets situés sur l'axe focal s'obtiennent directement : ce sont les points $S$ de $(\Delta)$ tels que $SF = e \cdot d(S, D)$, ce qui se résout par un simple calcul sur les distances.
\end{methode}

\begin{popfigure}
\begin{tikzpicture}[scale=0.8, >=Stealth]
% Axes
\draw[->] (-1,0) -- (5,0) node[right] {$x$};
\draw[->] (0,-3) -- (0,3) node[above] {$y$};
% Directrix
\draw[popblue, thick] (-1, -3) -- (-1, 3) node[above] {$D$};
\node[popblue] at (-1, -3.3) {$x=-2$};
% Focus
\fill[popred] (2,0) circle (2pt) node[below] {$F(2;0)$};
% Axis
\draw[popmuted, dashed] (-1,0) -- (5,0);
% Construction line d: x=3
\draw[poporange, thick] (3, -3) -- (3, 3) node[above] {$(d)$};
\node[poporange] at (3, -3.3) {$x=3$};
% Distance h = 5
\draw[<->, popmuted] (-1, -2.5) -- (3, -2.5) node[midway, below] {$h=5$};
% Circle center F radius 5
\draw[popgreen, dashed, thick] (2,0) circle (5);
% Intersections M1, M2
\fill[popgreen] (3, {sqrt(24)}) circle (2pt) node[right] {$M_1(3;2\sqrt{6})$};
\fill[popgreen] (3, {-sqrt(24)}) circle (2pt) node[right] {$M_2(3;-2\sqrt{6})$};
% Vertex S
\fill[poppurple] (0,0) circle (2pt) node[below left] {$S(0;0)$};
\end{tikzpicture}
\legende{Construction point par point de la parabole $MF=d(M,D)$ : intersections de $(d)$ avec le cercle $\mathcal{C}(F; e\cdot h)$.}
\end{popfigure}

\begin{exoresolu}[Construction d'une parabole point par point]
\textbf{Énoncé.} Dans le plan muni d'un repère orthonormé $(O\,;\vec{\imath},\vec{\jmath})$, on considère le point $F(2\,;\,0)$ et la droite $D$ d'équation $x = -2$.
\begin{enumerate}
\item Quelle est la nature de la conique $(\mathcal{P})$ d'équation $MF = d(M, D)$ ?
\item Déterminer les coordonnées des points de $(\mathcal{P})$ situés sur la droite d'équation $x = 3$.
\item Déterminer le sommet de $(\mathcal{P})$, puis expliquer comment construire géométriquement les points trouvés à la question 2.
\end{enumerate}
\tcblower
\textbf{Solution détaillée.}
\begin{enumerate}
\item La relation $MF = d(M, D)$ correspond à la définition monofocale avec $e = 1$ : $(\mathcal{P})$ est une \textbf{parabole} de foyer $F(2\,;\,0)$ et de directrice $D : x = -2$.

\item Soit $M(3\,;\,y)$ un point de la droite d'équation $x = 3$. On a :
\[ MF = \sqrt{(3-2)^2 + y^2} = \sqrt{1 + y^2} \qquad \text{et} \qquad d(M, D) = |3 - (-2)| = 5. \]
La condition $MF = d(M, D)$ s'écrit :
\[ \sqrt{1 + y^2} = 5 \iff 1 + y^2 = 25 \iff y^2 = 24 \iff y = \pm 2\sqrt{6}. \]
Les deux points cherchés sont $M_1(3\,;\,2\sqrt{6})$ et $M_2(3\,;\,-2\sqrt{6})$. On remarque qu'ils sont symétriques par rapport à l'axe des abscisses, qui est l'axe focal : c'est bien cohérent avec la symétrie de la parabole.

\item Le sommet $S$ est le point de la parabole situé sur l'axe focal (l'axe des abscisses), entre $F$ et $D$. Posons $S(x\,;\,0)$. La condition $SF = d(S, D)$ donne :
\[ |x - 2| = |x + 2| \iff x - 2 = -(x + 2) \text{ (car } -2 < x < 2\text{)} \iff 2x = 0 \iff x = 0. \]
Donc $S(0\,;\,0)$ : le sommet est le milieu du segment joignant $F$ à sa projection orthogonale sur $D$.

\textbf{Construction géométrique des points $M_1$ et $M_2$ :} on trace la droite $(d) : x = 3$, parallèle à $D$. La distance entre $(d)$ et $D$ vaut $5$. On trace le cercle de centre $F$ et de rayon $5$ (car $e \cdot h = 1 \times 5 = 5$) : ses intersections avec $(d)$ sont exactement $M_1$ et $M_2$.
\end{enumerate}
\textbf{Conclusion :} $(\mathcal{P})$ est la parabole de sommet $O$, et ses points d'abscisse $3$ sont $M_1(3\,;\,2\sqrt{6})$ et $M_2(3\,;\,-2\sqrt{6})$, constructibles à la règle et au compas.
\end{exoresolu}

\coursec{Équations réduites et éléments caractéristiques}

\begin{methode}[Équation réduite, foyers, directrices, sommets]
Dans un repère orthonormé bien choisi (origine au centre de symétrie pour ellipse/hyperbole, au sommet pour parabole ; axes = axes de symétrie) :
\begin{enumerate}
\item \textbf{Ellipse} ($0 < e < 1$) : $\dfrac{x^2}{a^2} + \dfrac{y^2}{b^2} = 1$ avec $a > b > 0$, $c = \sqrt{a^2 - b^2}$, $e = \dfrac{c}{a}$.
Foyers $F(c\,;\,0)$, $F'(-c\,;\,0)$ ; directrices $x = \dfrac{a}{e} = \dfrac{a^2}{c}$ et $x = -\dfrac{a^2}{c}$ ; sommets $A(a\,;\,0)$, $A'(-a\,;\,0)$ (grand axe), $B(0\,;\,b)$, $B'(0\,;\,-b)$ (petit axe).
\item \textbf{Hyperbole} ($e > 1$) : $\dfrac{x^2}{a^2} - \dfrac{y^2}{b^2} = 1$, $c = \sqrt{a^2 + b^2}$, $e = \dfrac{c}{a}$.
Foyers $F(c\,;\,0)$, $F'(-c\,;\,0)$ ; directrices $x = \pm \dfrac{a^2}{c}$ ; sommets $A(a\,;\,0)$, $A'(-a\,;\,0)$ ; asymptotes $y = \pm \dfrac{b}{a} x$.
\item \textbf{Parabole} ($e = 1$) : $y^2 = 2px$ ($p > 0$ est le \cle{paramètre}).
Foyer $F\left(\dfrac{p}{2}\,;\,0\right)$ ; directrice $x = -\dfrac{p}{2}$ ; sommet $O(0\,;\,0)$.
\end{enumerate}
\textbf{Démarche :} identifier les données de l'énoncé $\to$ en déduire $a$, $b$, $c$ ou $p$ par les relations ci-dessus $\to$ écrire l'équation $\to$ lister tous les éléments caractéristiques demandés.
\end{methode}

\begin{popfigure}
\begin{tikzpicture}[scale=0.8, >=Stealth, domain=-5:5]
% Ellipse
\begin{scope}[xshift=-7cm]
\draw[->] (-3.5,0) -- (3.5,0) node[right] {$x$};
\draw[->] (0,-2.5) -- (0,2.5) node[above] {$y$};
\draw[popgreen, thick, samples=200] plot (\x, {1.5*sqrt(max(0,1-\x*\x/4))});
\draw[popgreen, thick, samples=200] plot (\x, {-1.5*sqrt(max(0,1-\x*\x/4))});
\fill[popred] (1.32,0) circle (2pt) node[below] {$F$};
\fill[popred] (-1.32,0) circle (2pt) node[below] {$F'$};
\draw[popblue, dashed] (3.03, -2.5) -- (3.03, 2.5); % directrix a^2/c = 4/1.33 = 3
\draw[popblue, dashed] (-3.03, -2.5) -- (-3.03, 2.5);
\node[popgreen] at (0, -2.8) {Ellipse};
\end{scope}
% Hyperbola
\begin{scope}[xshift=7cm]
\draw[->] (-3.5,0) -- (3.5,0) node[right] {$x$};
\draw[->] (0,-2.5) -- (0,2.5) node[above] {$y$};
\draw[poporange, thick, samples=200, domain=-3.5:-1.2] plot (\x, {1.33*sqrt(max(0,\x*\x/4 - 1))});
\draw[poporange, thick, samples=200, domain=-3.5:-1.2] plot (\x, {-1.33*sqrt(max(0,\x*\x/4 - 1))});
\draw[poporange, thick, samples=200, domain=1.2:3.5] plot (\x, {1.33*sqrt(max(0,\x*\x/4 - 1))});
\draw[poporange, thick, samples=200, domain=1.2:3.5] plot (\x, {-1.33*sqrt(max(0,\x*\x/4 - 1))});
\fill[popred] (2.4,0) circle (2pt) node[below] {$F$};
\fill[popred] (-2.4,0) circle (2pt) node[below] {$F'$};
\draw[popblue, dashed] (1.67, -2.5) -- (1.67, 2.5); % directrix a^2/c = 4/2.4 = 1.67
\draw[popblue, dashed] (-1.67, -2.5) -- (-1.67, 2.5);
\draw[popmuted, dashed] (-3.5, -2.33) -- (3.5, 2.33); % asymptote
\draw[popmuted, dashed] (-3.5, 2.33) -- (3.5, -2.33);
\node[poporange] at (0, -2.8) {Hyperbole};
\end{scope}
\end{tikzpicture}
\legende{Éléments caractéristiques : foyers (rouge), directrices (bleu pointillé), asymptotes (gris pointillé pour l'hyperbole).}
\end{popfigure}

\begin{exoresolu}[Équation réduite d'une ellipse et éléments caractéristiques]
\textbf{Énoncé.} Dans un repère orthonormé $(O\,;\vec{\imath},\vec{\jmath})$, une ellipse $(\mathcal{E})$ a pour foyers $F(4\,;\,0)$ et $F'(-4\,;\,0)$ et pour excentricité $e = \dfrac{4}{5}$.
\begin{enumerate}
\item Déterminer l'équation réduite de $(\mathcal{E})$.
\item Préciser ses sommets, ses directrices et les longueurs de ses axes.
\end{enumerate}
\tcblower
\textbf{Solution détaillée.}
\begin{enumerate}
\item Les foyers sont $F(4\,;\,0)$ et $F'(-4\,;\,0)$, donc $c = 4$ et le centre de l'ellipse est $O$.

De $e = \dfrac{c}{a}$, on tire :
\[ a = \frac{c}{e} = \frac{4}{\frac{4}{5}} = 4 \times \frac{5}{4} = 5. \]
Pour une ellipse, $c^2 = a^2 - b^2$, donc :
\[ b^2 = a^2 - c^2 = 25 - 16 = 9 \quad \text{d'où} \quad b = 3. \]
L'équation réduite de $(\mathcal{E})$ est donc :
\[ \boxed{\frac{x^2}{25} + \frac{y^2}{9} = 1} \]

\item \textbf{Sommets :} $A(5\,;\,0)$, $A'(-5\,;\,0)$ sur le grand axe ; $B(0\,;\,3)$, $B'(0\,;\,-3)$ sur le petit axe.

\textbf{Directrices :} ce sont les droites d'équations $x = \pm \dfrac{a^2}{c}$, c'est-à-dire :
\[ x = \frac{25}{4} \qquad \text{et} \qquad x = -\frac{25}{4}. \]

\textbf{Longueurs des axes :} grand axe $2a = 10$ ; petit axe $2b = 6$.

\textbf{Vérification :} pour le sommet $A(5\,;\,0)$, on a $AF = 5 - 4 = 1$ et $d(A, D) = \dfrac{25}{4} - 5 = \dfrac{5}{4}$. Or $\dfrac{AF}{d(A,D)} = \dfrac{1}{\frac{5}{4}} = \dfrac{4}{5} = e$ : la définition monofocale est bien vérifiée.
\end{enumerate}
\textbf{Conclusion :} $(\mathcal{E})$ est l'ellipse d'équation $\dfrac{x^2}{25} + \dfrac{y^2}{9} = 1$, de sommets $(\pm 5\,;\,0)$ et $(0\,;\,\pm 3)$, de directrices $x = \pm \dfrac{25}{4}$, de grand axe $10$ et de petit axe $6$.
\end{exoresolu}

\coursec{Tangentes à une conique}

\begin{methode}[Tangente à une conique en un point $M_0(x_0\,;\,y_0)$]
\begin{enumerate}
\item \textbf{Vérifier d'abord} que $M_0$ appartient bien à la conique (en substituant ses coordonnées dans l'équation). Cette vérification est exigée dans la rédaction.
\item Appliquer la \cle{règle de dédoublement} : dans l'équation de la conique, remplacer
\[ x^2 \to x_0 x, \qquad y^2 \to y_0 y, \qquad x \to \frac{x + x_0}{2}, \qquad y \to \frac{y + y_0}{2}. \]
L'équation obtenue est celle de la tangente en $M_0$.
\item Formules prêtes à l'emploi (à connaître par cœur) :
\begin{itemize}
\item ellipse $\dfrac{x^2}{a^2} + \dfrac{y^2}{b^2} = 1$ : tangente $\dfrac{x_0 x}{a^2} + \dfrac{y_0 y}{b^2} = 1$ ;
\item hyperbole $\dfrac{x^2}{a^2} - \dfrac{y^2}{b^2} = 1$ : tangente $\dfrac{x_0 x}{a^2} - \dfrac{y_0 y}{b^2} = 1$ ;
\item parabole $y^2 = 2px$ : tangente $y_0 y = p(x + x_0)$.
\end{itemize}
\item Réduire l'équation sous la forme $y = mx + p$ si l'énoncé le demande, et conclure.
\end{enumerate}
\end{methode}

\begin{popfigure}
\begin{tikzpicture}[scale=0.9, >=Stealth, domain=-1:5]
% Axes
\draw[->] (-2.5,0) -- (5.5,0) node[right] {$x$};
\draw[->] (0,-3) -- (0,5) node[above] {$y$};
% Parabola y^2 = 8x -> y = +/- sqrt(8x)
\draw[popgreen, thick, samples=200, domain=0:5] plot (\x, {sqrt(max(0,8*\x))});
\draw[popgreen, thick, samples=200, domain=0:5] plot (\x, {-sqrt(max(0,8*\x))});
% Focus
\fill[popred] (2,0) circle (2pt) node[below] {$F(2;0)$};
% Directrix
\draw[popblue, thick] (-2, -3) -- (-2, 5) node[above] {$D$};
\node[popblue] at (-2, -3.3) {$x=-2$};
% Point M0
\fill[poppurple] (2,4) circle (2pt) node[above right] {$M_0(2;4)$};
% Tangent y = x+2
\draw[poporange, thick] (-2,0) -- (5,7) node[above right] {$(T): y=x+2$};
% Intersection I
\fill[poppink] (-2,0) circle (2pt) node[below left] {$I(-2;0)$};
% Segment M0H
\draw[popmuted, dashed] (2,4) -- (-2,4) node[midway, above] {$d(M_0,D)$};
\draw[popmuted, dashed] (2,4) -- (2,0) node[midway, right] {$M_0F$};
\end{tikzpicture}
\legende{Tangente à la parabole $y^2=8x$ en $M_0(2;4)$ : elle coupe l'axe focal en $I$ sur la directrice.}
\end{popfigure}

\begin{exoresolu}[Tangentes à une parabole]
\textbf{Énoncé.} Une antenne parabolique d'un abonné Canal+ à Douala a pour coupe la parabole $(\mathcal{P})$ d'équation $y^2 = 8x$ dans un repère orthonormé.
\begin{enumerate}
\item Déterminer le foyer et la directrice de $(\mathcal{P})$.
\item Vérifier que le point $M_0(2\,;\,4)$ appartient à $(\mathcal{P})$, puis déterminer une équation de la tangente $(T)$ à $(\mathcal{P})$ en $M_0$.
\item La tangente $(T)$ coupe l'axe des abscisses en un point $I$. Calculer les coordonnées de $I$.
\end{enumerate}
\tcblower
\textbf{Solution détaillée.}
\begin{enumerate}
\item L'équation $y^2 = 8x$ est de la forme $y^2 = 2px$ avec $2p = 8$, donc $p = 4$.
Le foyer est $F\left(\dfrac{p}{2}\,;\,0\right) = F(2\,;\,0)$ et la directrice a pour équation $x = -\dfrac{p}{2}$, c'est-à-dire $x = -2$.

\item \textbf{Appartenance :} $y_0^2 = 4^2 = 16$ et $8x_0 = 8 \times 2 = 16$. Donc $y_0^2 = 8x_0$ : le point $M_0(2\,;\,4)$ appartient bien à $(\mathcal{P})$.

\textbf{Tangente :} par la règle de dédoublement, on remplace $y^2$ par $y_0 y$ et $x$ par $\dfrac{x + x_0}{2}$ :
\[ y_0 y = 8 \times \frac{x + x_0}{2} \iff 4y = 4(x + 2) \iff y = x + 2. \]
La tangente $(T)$ a donc pour équation $\boxed{y = x + 2}$.

\item Le point $I$ est l'intersection de $(T)$ avec l'axe des abscisses ($y = 0$) :
\[ 0 = x + 2 \iff x = -2. \]
Donc $I(-2\,;\,0)$. On remarque que $I$ est le point de la directrice situé sur l'axe focal : c'est une propriété classique de la tangente à la parabole (le pied de la directrice sur l'axe).
\end{enumerate}
\textbf{Conclusion :} $(\mathcal{P})$ a pour foyer $F(2\,;\,0)$ et directrice $x = -2$ ; la tangente en $M_0(2\,;\,4)$ a pour équation $y = x + 2$ et coupe l'axe focal en $I(-2\,;\,0)$.
\end{exoresolu}

\coursec{Reconnaissance d'une équation du second degré}

\begin{methode}[Mise sous forme réduite par complétion de carrés]
\begin{enumerate}
\item \textbf{Diagnostic rapide} sur les coefficients de $x^2$ et $y^2$ (notés $a$ et $b$) :
\begin{itemize}
\item $a$ et $b$ de même signe ($ab > 0$) : ellipse (cercle si $a = b$), ou ensemble vide, ou un point ;
\item $a$ et $b$ de signes contraires ($ab < 0$) : hyperbole (ou deux droites sécantes si le second membre s'annule après réduction) ;
\item un seul des deux coefficients non nul : parabole (ou droites parallèles, ou ensemble vide).
\end{itemize}
\item \textbf{Regrouper} les termes en $x$ d'une part, en $y$ d'autre part, et former des \cle{carrés parfaits} :
\[ ax^2 + cx = a\left(x + \frac{c}{2a}\right)^2 - \frac{c^2}{4a}. \]
\item \textbf{Translater} le repère : poser $X = x - x_0$, $Y = y - y_0$ où $(x_0\,;\,y_0)$ est le centre (ellipse/hyperbole) ou le sommet (parabole) trouvé.
\item \textbf{Diviser} par le second membre pour obtenir une équation réduite du type $\dfrac{X^2}{a^2} \pm \dfrac{Y^2}{b^2} = 1$ ou $Y^2 = 2pX$.
\item \textbf{Conclure} : nature de la conique, centre/sommet dans le repère initial, et éléments caractéristiques si demandés.
\end{enumerate}
\textbf{Attention :} si le second membre est négatif pour une somme de carrés, l'ensemble est vide ; s'il est nul, l'ensemble est réduit à un point (cas ellipse) ou à deux droites (cas hyperbole).
\end{methode}

\begin{popfigure}
\begin{tikzpicture}[scale=0.8, >=Stealth]
% Original axes
\draw[->, popmuted] (-1,0) -- (6,0) node[right] {$x$};
\draw[->, popmuted] (0,-3) -- (0,3) node[above] {$y$};
% Hyperbola center (2,1)
\fill[popred] (2,1) circle (2pt) node[above right] {$\Omega(2;1)$};
% Asymptotes
\draw[popmuted, dashed] (-1, -0.33) -- (6, 4.33);
\draw[popmuted, dashed] (-1, 2.33) -- (6, -2.67);
% Hyperbola branches (x-2)^2/9 - (y-1)^2/4 = 1
% Right branch
\draw[poporange, thick, samples=100, domain=2+3:6] plot (\x, {1 + 2*sqrt(max(0,(\x-2)*(\x-2)/9 - 1))});
\draw[poporange, thick, samples=100, domain=2+3:6] plot (\x, {1 - 2*sqrt(max(0,(\x-2)*(\x-2)/9 - 1))});
% Left branch
\draw[poporange, thick, samples=100, domain=-1:2-3] plot (\x, {1 + 2*sqrt(max(0,(\x-2)*(\x-2)/9 - 1))});
\draw[poporange, thick, samples=100, domain=-1:2-3] plot (\x, {1 - 2*sqrt(max(0,(\x-2)*(\x-2)/9 - 1))});
% Vertices
\fill[popgreen] (5,1) circle (2pt) node[right] {$A(5;1)$};
\fill[popgreen] (-1,1) circle (2pt) node[left] {$A'(-1;1)$};
% Translated axes
\draw[->, poppurple, thick] (2,1) -- (5,1) node[right] {$X$};
\draw[->, poppurple, thick] (2,1) -- (2,3) node[above] {$Y$};
\end{tikzpicture}
\legende{Hyperbole $4x^2-9y^2-16x+18y-29=0$ : centre $\Omega(2;1)$, axes translatés $(X,Y)$, asymptotes.}
\end{popfigure}

\begin{exoresolu}[Réduction d'une équation du second degré]
\textbf{Énoncé.} Le plan est muni d'un repère orthonormé $(O\,;\vec{\imath},\vec{\jmath})$. On considère l'ensemble $(\mathcal{C})$ des points $M(x\,;\,y)$ tels que :
\[ 4x^2 - 9y^2 - 16x + 18y - 29 = 0. \]
\begin{enumerate}
\item Montrer que $(\mathcal{C})$ est une hyperbole dont on précisera le centre $\Omega$.
\item Donner son équation réduite, ses sommets, son excentricité et ses asymptotes.
\end{enumerate}
\tcblower
\textbf{Solution détaillée.}
\begin{enumerate}
\item Les coefficients de $x^2$ et $y^2$ sont $4$ et $-9$, de signes contraires : on s'attend à une hyperbole. Procédons par complétion de carrés.

\textbf{Termes en $x$ :}
\[ 4x^2 - 16x = 4(x^2 - 4x) = 4\left[(x - 2)^2 - 4\right] = 4(x-2)^2 - 16. \]

\textbf{Termes en $y$ :}
\[ -9y^2 + 18y = -9(y^2 - 2y) = -9\left[(y - 1)^2 - 1\right] = -9(y-1)^2 + 9. \]

L'équation devient :
\[ 4(x-2)^2 - 16 - 9(y-1)^2 + 9 - 29 = 0 \iff 4(x-2)^2 - 9(y-1)^2 = 36. \]
En divisant par $36$ :
\[ \frac{(x-2)^2}{9} - \frac{(y-1)^2}{4} = 1. \]
C'est l'équation réduite d'une \textbf{hyperbole} de centre $\Omega(2\,;\,1)$.

\item Posons $X = x - 2$ et $Y = y - 1$. L'équation s'écrit $\dfrac{X^2}{9} - \dfrac{Y^2}{4} = 1$, donc $a^2 = 9$ et $b^2 = 4$, soit $a = 3$ et $b = 2$.

\textbf{Sommets} (dans le repère initial) : $A(\Omega_x + a\,;\,\Omega_y) = A(5\,;\,1)$ et $A'(-1\,;\,1)$.

\textbf{Excentricité :} $c = \sqrt{a^2 + b^2} = \sqrt{9 + 4} = \sqrt{13}$, donc :
\[ e = \frac{c}{a} = \frac{\sqrt{13}}{3}. \]

\textbf{Asymptotes :} dans le repère translaté, $Y = \pm \dfrac{b}{a} X$, c'est-à-dire dans le repère initial :
\[ y - 1 = \frac{2}{3}(x - 2) \qquad \text{et} \qquad y - 1 = -\frac{2}{3}(x - 2), \]
soit $y = \dfrac{2}{3}x - \dfrac{1}{3}$ et $y = -\dfrac{2}{3}x + \dfrac{7}{3}$.
\end{enumerate}
\textbf{Conclusion :} $(\mathcal{C})$ est l'hyperbole de centre $\Omega(2\,;\,1)$, d'équation réduite $\dfrac{(x-2)^2}{9} - \dfrac{(y-1)^2}{4} = 1$, de sommets $(5\,;\,1)$ et $(-1\,;\,1)$, d'excentricité $\dfrac{\sqrt{13}}{3}$.
\end{exoresolu}

\coursec{Coniques, similitudes et situations concrètes}

\begin{methode}[Transformer une conique par une similitude directe]
Une similitude directe $s$ de rapport $k > 0$ multiplie toutes les distances par $k$.
\begin{enumerate}
\item \textbf{Image du foyer et de la directrice :} déterminer $F' = s(F)$ et $D' = s(D)$ (l'image d'une droite par une similitude est une droite).
\item \textbf{Excentricité conservée :} pour tout point $M$ d'image $M'$ :
\[ M'F' = k \cdot MF \qquad \text{et} \qquad d(M', D') = k \cdot d(M, D), \]
donc $\dfrac{M'F'}{d(M', D')} = \dfrac{MF}{d(M, D)} = e$ : l'excentricité est \cle{invariante}.
\item \textbf{Conclure sur la nature :} l'image d'une conique d'excentricité $e$ est une conique de \emph{même nature} (ellipse $\to$ ellipse, parabole $\to$ parabole, hyperbole $\to$ hyperbole), de foyer $F'$ et de directrice $D'$.
\item Cas particulier : une \textbf{homothétie} de rapport $k$ transforme la directrice $D$ en une droite $D'$ parallèle à $D$, et le paramètre (ou les demi-axes $a$, $b$) est multiplié par $|k|$.
\end{enumerate}
\textbf{Réflexe :} au Bac, on demande souvent de justifier que l'image est une conique de même excentricité : rédiger proprement l'égalité des rapports de distances.
\end{methode}

\begin{popfigure}
\begin{tikzpicture}[scale=0.7, >=Stealth]
% Axes
\draw[->] (-4,0) -- (8,0) node[right] {$x$};
\draw[->] (0,-4) -- (0,6) node[above] {$y$};
% Origin
\fill[popdark] (0,0) circle (2pt) node[below left] {$O$};
% Original Parabola y^2=4x (p=2)
\draw[popblue, thick, samples=100, domain=0:7] plot (\x, {2*sqrt(max(0,\x))});
\draw[popblue, thick, samples=100, domain=0:7] plot (\x, {-2*sqrt(max(0,\x))});
\node[popblue] at (7, 5.5) {$(\mathcal{P}): y^2=4x$};
% Focus F(1,0)
\fill[popred] (1,0) circle (2pt) node[below] {$F(1;0)$};
% Directrix x=-1
\draw[popblue, dashed] (-1, -4) -- (-1, 6) node[above] {$D$};
% Image Parabola y^2=12x (p=6)
\draw[poporange, thick, samples=100, domain=0:7] plot (\x, {sqrt(max(0,12*\x))});
\draw[poporange, thick, samples=100, domain=0:7] plot (\x, {-sqrt(max(0,12*\x))});
\node[poporange] at (7, -3.5) {$(\mathcal{P}'): y^2=12x$};
% Focus F'(3,0)
\fill[popred] (3,0) circle (2pt) node[below] {$F'(3;0)$};
% Directrix x=-3
\draw[poporange, dashed] (-3, -4) -- (-3, 6) node[above] {$D'$};
% Homothety rays
\draw[popmuted, thin, dashed] (0,0) -- (1,0) -- (3,0);
\draw[popmuted, thin, dashed] (0,0) -- (-1, 2) -- (-3, 6);
\draw[popmuted, thin, dashed] (0,0) -- (-1, -2) -- (-3, -6);
% Ratio
\node[popmuted] at (0.5, 0.3) {$k=3$};
\node[popmuted] at (2, 0.3) {$k=3$};
\end{tikzpicture}
\legende{Image de la parabole $y^2=4x$ par l'homothétie $h(O,3)$ : la parabole $y^2=12x$.}
\end{popfigure}

\begin{exoresolu}[Image d'une parabole par une homothétie]
\textbf{Énoncé.} Dans le plan muni d'un repère orthonormé, soit $(\mathcal{P})$ la parabole de foyer $F(1\,;\,0)$ et de directrice $D$ d'équation $x = -1$. Soit $h$ l'homothétie de centre $O$ et de rapport $3$.
\begin{enumerate}
\item Déterminer le foyer $F'$ et la directrice $D'$ de la conique $(\mathcal{P}')$, image de $(\mathcal{P})$ par $h$.
\item Justifier que $(\mathcal{P}')$ est une parabole et donner son équation.
\end{enumerate}
\tcblower
\textbf{Solution détaillée.}
\begin{enumerate}
\item \textbf{Image du foyer :} $F' = h(F)$. Une homothétie de centre $O$ et de rapport $3$ envoie tout point $M(x\,;\,y)$ sur $M'(3x\,;\,3y)$. Donc :
\[ F' = (3 \times 1\,;\, 3 \times 0) = (3\,;\,0). \]

\textbf{Image de la directrice :} l'image d'une droite par une homothétie est une droite parallèle. La droite $D : x = -1$ a pour image la droite $D' : x = 3 \times (-1) = -3$, c'est-à-dire $x = -3$.

Vérifions que $F' \notin D'$ : l'abscisse de $F'$ est $3 \neq -3$, donc $F' \notin D'$ : la définition monofocale s'applique bien.

\item Soit $M$ un point de $(\mathcal{P})$ et $M' = h(M)$. L'homothétie multiplie les distances par $|k| = 3$ :
\[ M'F' = 3\, MF \qquad \text{et} \qquad d(M', D') = 3\, d(M, D). \]
Comme $M \in (\mathcal{P})$, on a $MF = d(M, D)$, donc :
\[ M'F' = 3\, MF = 3\, d(M, D) = d(M', D'). \]
Réciproquement, tout point vérifiant $M'F' = d(M', D')$ provient d'un point de $(\mathcal{P})$ (car $h$ est bijective). Donc $(\mathcal{P}')$ est l'ensemble des points équidistants de $F'(3\,;\,0)$ et de $D' : x = -3$ : c'est une \textbf{parabole} (excentricité $e = 1$ conservée).

\textbf{Équation :} le sommet est le milieu de $F'$ et de sa projection sur $D'$, soit $O(0\,;\,0)$, et le paramètre vérifie $\dfrac{p'}{2} = 3$, donc $p' = 6$. L'équation de $(\mathcal{P}')$ est :
\[ y^2 = 2p' x \iff \boxed{y^2 = 12x}. \]
On remarque que le paramètre a bien été multiplié par $3$ (la parabole initiale était $y^2 = 4x$).
\end{enumerate}
\textbf{Conclusion :} $(\mathcal{P}')$ est la parabole de foyer $F'(3\,;\,0)$, de directrice $x = -3$ et d'équation $y^2 = 12x$ : l'homothétie a conservé la nature et l'excentricité de la conique.
\end{exoresolu}

\begin{attention}[Les 5 erreurs qui coûtent des points au Bac]
\begin{enumerate}
\item \textbf{Confondre les relations sur $c$.} Pour l'ellipse, $c^2 = a^2 - b^2$ ; pour l'hyperbole, $c^2 = a^2 + b^2$. Inverser ces deux formules fausse l'excentricité, les foyers et les directrices. Astuce : pour l'hyperbole, $c$ est le « plus grand » ($c > a$ car $e > 1$) ; pour l'ellipse, $c < a$.
\item \textbf{Oublier de vérifier l'appartenance du point avant d'écrire la tangente.} La règle de dédoublement n'est valable que si $M_0(x_0\,;\,y_0)$ appartient à la conique. Rédiger explicitement la vérification ($y_0^2 = \ldots = 2px_0$) : c'est un point de méthode noté.
\item \textbf{Se tromper de signe dans la complétion de carrés.} Par exemple, $-9y^2 + 18y = -9(y-1)^2 + 9$ et non $-9(y-1)^2 - 9$. Toujours développer mentalement le résultat pour contrôler, et vérifier le second membre final : un second membre négatif avec une somme de carrés signifie un ensemble vide, pas une ellipse.
\item \textbf{Écrire une excentricité incompatible avec la conclusion.} Affirmer « c'est une ellipse » avec $e = \dfrac{5}{3} > 1$ est une contradiction qui annule la question. Toujours comparer explicitement $e$ à $1$ avant de nommer la conique, et vérifier la cohérence finale.
\item \textbf{Mal placer foyers et directrices par rapport au centre.} Pour l'ellipse, la directrice $x = \dfrac{a^2}{c}$ est \emph{à l'extérieur} de l'ellipse ($\dfrac{a^2}{c} > a$ car $c < a$) ; pour l'hyperbole, elle passe \emph{entre} le centre et le sommet ($\dfrac{a^2}{c} < a$). Un contrôle rapide de ces inégalités détecte la plupart des erreurs de calcul sur $a$ et $c$.
\end{enumerate}
\end{attention}

\newpage

\coursec{Exercices}

\courssub{Je vérifie mes connaissances}

\begin{exercice}[QCM : nature d'une conique]
Pour chaque question, choisir la (ou les) bonne(s) réponse(s) en justifiant brièvement.
\begin{enumerate}
\item Une conique d'excentricité $e = \dfrac{5}{4}$ est :
\begin{enumerate}
\item une ellipse \quad b) une parabole \quad c) une hyperbole \quad d) un cercle
\end{enumerate}
\item L'équation $\dfrac{x^{2}}{25} + \dfrac{y^{2}}{9} = 1$ est celle :
\begin{enumerate}
\item d'une ellipse de demi-grand axe $5$ \quad b) d'une hyperbole \quad c) d'une ellipse de distance focale $2c = 8$ \quad d) d'une ellipse d'excentricité $\dfrac{4}{5}$
\end{enumerate}
\item La parabole d'équation $y^{2} = 8x$ a pour :
\begin{enumerate}
\item foyer $F(2\,;\,0)$ \quad b) directrice d'équation $x = -2$ \quad c) foyer $F(8\,;\,0)$ \quad d) paramètre $p = 8$
\end{enumerate}
\item L'image d'une ellipse par une similitude directe de rapport $k > 0$ est :
\begin{enumerate}
\item toujours un cercle \quad b) une ellipse de même excentricité \quad c) une ellipse d'excentricité multipliée par $k$ \quad d) une conique de même nature
\end{enumerate}
\end{enumerate}
\end{exercice}

\begin{exercice}[Vrai ou faux, justifié]
Dire si chacune des affirmations suivantes est vraie ou fausse, en justifiant la réponse par une démonstration ou un contre-exemple.
\begin{enumerate}
\item Toute section plane d'un cône de révolution est une conique non dégénérée.
\item Une conique d'excentricité $e = 0$ est une parabole.
\item L'hyperbole d'équation $\dfrac{x^{2}}{16} - \dfrac{y^{2}}{9} = 1$ admet pour asymptotes les droites d'équations $y = \dfrac{3}{4}x$ et $y = -\dfrac{3}{4}x$.
\item Si $MF = e \cdot d(M, D)$ avec $e = 1$, alors la conique de foyer $F$ et de directrice $D$ passe par le milieu du segment joignant $F$ à son projeté orthogonal sur $D$.
\end{enumerate}
\end{exercice}

\begin{exercice}[Définitions et vocabulaire]
\begin{enumerate}
\item Donner la définition monofocale d'une conique de foyer $F$, de directrice $D$ et d'excentricité $e$, puis préciser la nature de la conique selon les valeurs de $e$.
\item Citer les éléments de symétrie d'une ellipse et d'une hyperbole (axes, centre), puis ceux d'une parabole.
\item Qu'appelle-t-on \cle{distance focale} d'une ellipse ? Quelle relation lie $a$, $b$ et $c$ pour une ellipse ? Pour une hyperbole ?
\end{enumerate}
\end{exercice}

\begin{exercice}[Calculs directs]
\begin{enumerate}
\item Déterminer l'excentricité de l'ellipse d'équation $\dfrac{x^{2}}{100} + \dfrac{y^{2}}{64} = 1$.
\item Déterminer les sommets et les asymptotes de l'hyperbole d'équation $\dfrac{x^{2}}{4} - \dfrac{y^{2}}{25} = 1$.
\item Déterminer le foyer et la directrice de la parabole d'équation $x^{2} = -12y$.
\item Identifier la nature de la conique d'équation $4x^{2} + 9y^{2} = 36$ et préciser ses éléments caractéristiques.
\end{enumerate}
\end{exercice}

\courssub{Je m'entraîne}

\begin{exercice}[Ellipse à partir de l'excentricité et de la distance focale]
Soit une ellipse $(E)$ de centre $O$, d'axe focal $(Ox)$, d'excentricité $e = \dfrac{1}{2}$ et de distance focale $2c = 4$.
\begin{enumerate}
\item Déterminer $a$, $b$ et $c$.
\item Écrire l'équation réduite de $(E)$.
\item Déterminer les coordonnées des foyers et les équations des directrices.
\item Déterminer les coordonnées des quatre sommets de $(E)$.
\end{enumerate}
\end{exercice}

\begin{exercice}[Tangentes à une hyperbole]
On considère l'hyperbole $(H)$ d'équation $\dfrac{x^{2}}{9} - \dfrac{y^{2}}{4} = 1$.
\begin{enumerate}
\item Déterminer les équations des tangentes à $(H)$ qui sont parallèles à la droite d'équation $y = x$. (On pourra chercher les tangentes sous la forme $y = x + m$ et écrire que l'équation aux abscisses des points d'intersection admet une racine double.)
\item Vérifier que le point $M_{0}\left(5\,;\,\dfrac{8}{3}\right)$ appartient à $(H)$, puis déterminer l'équation de la tangente à $(H)$ en $M_{0}$.
\end{enumerate}
\end{exercice}

\begin{exercice}[Parabole et propriété optique]
On considère la parabole $(P)$ de sommet $O$, d'axe $(Ox)$, passant par le point $A(2\,;\,4)$.
\begin{enumerate}
\item Déterminer l'équation réduite de $(P)$, son foyer $F$ et sa directrice $D$.
\item Déterminer l'équation de la tangente $(T)$ à $(P)$ au point $A$.
\item Soit $H$ le projeté orthogonal de $A$ sur $D$. Montrer que la tangente $(T)$ est la médiatrice du segment $[FH]$. Quelle propriété optique de la parabole ce résultat illustre-t-il ? Citer une application concrète de cette propriété.
\end{enumerate}
\end{exercice}

\begin{exercice}[Image d'une ellipse par une similitude]
Le plan est muni d'un repère orthonormé $(O\,;\vec{i},\vec{j})$. On considère l'ellipse $(E)$ d'équation $\dfrac{x^{2}}{16} + \dfrac{y^{2}}{9} = 1$ et la similitude directe $s$ de centre $O$, de rapport $2$ et d'angle $\dfrac{\pi}{2}$.
\begin{enumerate}
\item Rappeler l'excentricité de $(E)$.
\item Déterminer une équation de la courbe $(E')$, image de $(E)$ par $s$. (On pourra exprimer les coordonnées $(x\,;\,y)$ d'un point en fonction de celles $(x'\,;\,y')$ de son image.)
\item Préciser la nature de $(E')$, ses sommets et son excentricité. Commenter.
\end{enumerate}
\end{exercice}

\courssub{Je me prépare au Bac}

\begin{exercice}[Type Bac : construction et étude complète d'une conique]
Le plan est muni d'un repère orthonormé $(O\,;\vec{i},\vec{j})$ (unité : $\SI{1}{cm}$). On donne le point $F(3\,;\,0)$, la droite $(D)$ d'équation $x = \dfrac{25}{3}$ et le réel $e = \dfrac{3}{5}$. On note $(\mathcal{C})$ l'ensemble des points $M$ du plan tels que $MF = e \cdot d(M, (D))$.
\begin{enumerate}
\item Quelle est la nature de la conique $(\mathcal{C})$ ? Justifier.
\item Construire point par point six points de $(\mathcal{C})$ en utilisant la définition monofocale (on décrira précisément la construction à la règle et au compas).
\item Montrer que $(\mathcal{C})$ admet pour équation réduite $\dfrac{x^{2}}{25} + \dfrac{y^{2}}{16} = 1$ dans le repère $(O\,;\vec{i},\vec{j})$.
\item Déterminer les coordonnées des sommets de $(\mathcal{C})$, de son second foyer $F'$ et l'équation de sa seconde directrice $(D')$.
\item Déterminer une équation de la tangente à $(\mathcal{C})$ au point d'abscisse $3$ et d'ordonnée positive.
\end{enumerate}
\end{exercice}

\begin{exercice}[Type Bac : reconnaissance d'une courbe du second degré]
Le plan est muni d'un repère orthonormé $(O\,;\vec{i},\vec{j})$. On considère la courbe $(\Gamma)$ d'équation :
\[ 9x^{2} - 4y^{2} + 36x + 8y - 4 = 0. \]
\begin{enumerate}
\item Montrer qu'il existe un point $\Omega$ et un repère $(\Omega\,;\vec{i},\vec{j})$ dans lequel $(\Gamma)$ admet une équation de la forme $\dfrac{X^{2}}{a^{2}} - \dfrac{Y^{2}}{b^{2}} = 1$. En déduire la nature de $(\Gamma)$.
\item Déterminer, dans le repère $(O\,;\vec{i},\vec{j})$, les coordonnées du centre, des sommets et des foyers de $(\Gamma)$.
\item Déterminer les équations des asymptotes de $(\Gamma)$ dans le repère $(O\,;\vec{i},\vec{j})$.
\item Déterminer l'excentricité de $(\Gamma)$ et les équations de ses directrices.
\item Tracer soigneusement $(\Gamma)$, ses asymptotes et ses axes de symétrie.
\item Déterminer une équation de la tangente à $(\Gamma)$ en son sommet d'abscisse la plus grande.
\end{enumerate}
\end{exercice}

\begin{exercice}[Type Bac : coniques autour de nous]
\textbf{Partie A.} Un ballon sphérique est posé au sol en plein soleil, en fin de journée, de sorte que les rayons du soleil sont quasi parallèles. L'ombre du ballon se projette sur un mur vertical voisin.
\begin{enumerate}
\item Expliquer pourquoi l'ensemble des rayons tangents à la sphère forme un cylindre, et pourquoi l'ombre portée sur le mur est une conique.
\item Le mur étant perpendiculaire à la direction des rayons, quelle est la nature de l'ombre ? Et si le mur est incliné par rapport à cette direction ?
\end{enumerate}
\textbf{Partie B.} L'arche d'un pont sur la route Yaoundé--Douala a la forme d'un arc de parabole de portée $\SI{60}{m}$ (distance entre les deux pieds de l'arche, au niveau du sol) et de hauteur $\SI{15}{m}$ au-dessus du sol. On munit le plan d'un repère orthonormé d'origine au sommet de l'arche, l'axe $(Ox)$ étant horizontal et l'axe $(Oy)$ vertical dirigé vers le bas, l'unité étant le mètre.
\begin{enumerate}
\item Déterminer l'équation de la parabole dans ce repère.
\item En déduire le foyer et la directrice de cette parabole.
\item Un camion de $\SI{4}{m}$ de haut et de $\SI{3}{m}$ de large circule en roulant à $\SI{2}{m}$ du pied droit de l'arche (le bord du camion le plus proche du pied est à $\SI{2}{m}$ de celui-ci). Peut-il passer sous l'arche ? Justifier par un calcul.
\item À quelle distance horizontale du sommet de l'arche la hauteur sous l'arche est-elle égale à $\SI{10}{m}$ ?
\end{enumerate}
\end{exercice}

\newpage

\coursec{Corrigés des exercices}

\begin{corrige}[Exercice 1 — QCM : nature d'une conique]
\begin{enumerate}
\item L'excentricité est $e = \dfrac{5}{4} > 1$, donc la conique est une \textbf{hyperbole} (par définition monofocale). Réponse \textbf{c)}.
\item L'équation $\dfrac{x^{2}}{25} + \dfrac{y^{2}}{9} = 1$ est une ellipse ($a=5$, $b=3$).
  \begin{itemize}
    \item a) $a = 5$ : \textbf{vrai}.
    \item b) C'est une ellipse, pas une hyperbole : \textbf{faux}.
    \item c) $c^{2} = a^{2} - b^{2} = 25 - 9 = 16 \Rightarrow c = 4$, distance focale $2c = 8$ : \textbf{vrai}.
    \item d) $e = \dfrac{c}{a} = \dfrac{4}{5}$ : \textbf{vrai}.
  \end{itemize}
  Réponses correctes : \textbf{a), c), d)}.
\item Parabole $y^{2} = 8x$ : forme $y^{2} = 2px$ avec $2p = 8 \Rightarrow p = 4$.
  \begin{itemize}
    \item a) Foyer $F\left(\dfrac{p}{2}\,;\,0\right) = F(2\,;\,0)$ : \textbf{vrai}.
    \item b) Directrice $x = -\dfrac{p}{2} = -2$ : \textbf{vrai}.
    \item c) (Non listée mais implicite) Axe $(Ox)$, sommet $O$.
    \item d) Si l'option d) affirme $p=8$ ou $F(4\,;\,0)$, elle est \textbf{fausse} car $p=4$.
  \end{itemize}
  Réponses correctes : \textbf{a), b)}.
\item Une similitude directe conserve les rapports de distances, donc l'excentricité $e$ est invariante ; la nature de la conique (ellipse, parabole, hyperbole) est conservée.
  \begin{itemize}
    \item a) Faux : l'image d'une ellipse non circulaire par une similitude reste une ellipse non circulaire (le rapport des axes est conservé).
    \item b) Vrai : la nature est conservée.
    \item c) Faux : $e$ est invariant.
    \item d) Vrai : $e$ est conservée.
  \end{itemize}
  Réponses correctes : \textbf{b), d)}.
\end{enumerate}
\end{corrige}

\begin{corrige}[Exercice 2 — Vrai ou faux, justifié]
\begin{enumerate}
\item \textbf{Faux.} Si le plan de section passe par le sommet du cône, la section est \textbf{dégénérée} : un point (plan perpendiculaire à l'axe), une droite (plan tangent au cône) ou deux droites sécantes (plan passant par le sommet et coupant les deux nappes). Contre-exemple : un plan passant par le sommet et parallèle à une génératrice donne deux génératrices (deux droites sécantes au sommet).
\item \textbf{Faux.} $e = 0$ correspond au \textbf{cercle} (cas limite de l'ellipse où les deux foyers sont confondus, $c=0$). La parabole correspond à $e = 1$.
\item \textbf{Vrai.} Les asymptotes de l'hyperbole $\dfrac{x^{2}}{a^{2}} - \dfrac{y^{2}}{b^{2}} = 1$ sont $y = \pm \dfrac{b}{a}x$. Ici $a=4$, $b=3$, donc $y = \pm \dfrac{3}{4}x$.
\item \textbf{Vrai.} Soit $K$ le projeté orthogonal du foyer $F$ sur la directrice $D$ et $S$ le milieu de $[FK]$. Alors $SF = \dfrac{FK}{2}$ et $d(S, D) = SK = \dfrac{FK}{2}$, donc $SF = d(S, D)$, c'est-à-dire $SF = 1 \cdot d(S, D)$ avec $e = 1$ : $S$ appartient à la parabole (c'est son \cle{sommet}).
\end{enumerate}
\end{corrige}

\begin{corrige}[Exercice 3 — Définitions et vocabulaire]
\begin{enumerate}
\item Soit $F$ un point (foyer), $D$ une droite (directrice) ne passant pas par $F$ et $e > 0$ un réel (excentricité). La conique de foyer $F$, de directrice $D$ et d'excentricité $e$ est l'ensemble des points $M$ du plan tels que $MF = e \cdot d(M, D)$.
  \begin{itemize}
    \item Si $0 < e < 1$ : \cle{ellipse}.
    \item Si $e = 1$ : \cle{parabole}.
    \item Si $e > 1$ : \cle{hyperbole}.
  \end{itemize}
\item \textbf{Éléments de symétrie :}
  \begin{itemize}
    \item Ellipse et hyperbole : deux axes de symétrie orthogonaux (l'\cle{axe focal} et l'\cle{axe non focal}) et un centre de symétrie (le \cle{centre} de la conique, milieu des foyers).
    \item Parabole : un seul axe de symétrie (l'axe focal, perpendiculaire à la directrice et passant par le foyer), pas de centre de symétrie.
  \end{itemize}
\item \textbf{Distance focale} : distance $2c$ entre les deux foyers.
  \begin{itemize}
    \item Ellipse (axe focal horizontal, $a > b$) : $a^{2} = b^{2} + c^{2}$ (donc $c^{2} = a^{2} - b^{2}$).
    \item Hyperbole (axe focal horizontal) : $c^{2} = a^{2} + b^{2}$.
  \end{itemize}
\end{enumerate}
\end{corrige}

\begin{corrige}[Exercice 4 — Calculs directs]
\begin{enumerate}
\item Ellipse $\dfrac{x^{2}}{100} + \dfrac{y^{2}}{64} = 1$ : $a=10$, $b=8$. $c^{2} = a^{2} - b^{2} = 100 - 64 = 36 \Rightarrow c = 6$. Excentricité $e = \dfrac{c}{a} = \dfrac{6}{10} = \dfrac{3}{5}$.
\item Hyperbole $\dfrac{x^{2}}{4} - \dfrac{y^{2}}{25} = 1$ : $a=2$, $b=5$. Sommets sur l'axe focal : $A(2\,;\,0)$ et $A'(-2\,;\,0)$. Asymptotes : $y = \pm \dfrac{b}{a}x = \pm \dfrac{5}{2}x$.
\item Parabole $x^{2} = -12y$ : forme $x^{2} = -2py$ avec $2p = 12 \Rightarrow p = 6$. Axe $(Oy)$, ouverture vers le bas (dans le repère standard). Foyer $F\left(0\,;\,-\dfrac{p}{2}\right) = F(0\,;\,-3)$ ; directrice $y = \dfrac{p}{2} = 3$.
\item $4x^{2} + 9y^{2} = 36 \iff \dfrac{x^{2}}{9} + \dfrac{y^{2}}{4} = 1$ : ellipse de centre $O$, $a=3$, $b=2$, $c = \sqrt{9-4} = \sqrt{5}$.
  \begin{itemize}
    \item Sommets : $(\pm 3\,;\,0)$ et $(0\,;\,\pm 2)$.
    \item Foyers : $(\pm\sqrt{5}\,;\,0)$.
    \item Excentricité : $e = \dfrac{\sqrt{5}}{3}$.
    \item Directrices : $x = \pm \dfrac{a^{2}}{c} = \pm \dfrac{9}{\sqrt{5}} = \pm \dfrac{9\sqrt{5}}{5}$.
  \end{itemize}
\end{enumerate}
\end{corrige}

\begin{corrige}[Exercice 5 — Ellipse à partir de l'excentricité et de la distance focale]
\begin{enumerate}
\item Distance focale $2c = 4 \Rightarrow c = 2$. Excentricité $e = \dfrac{c}{a} = \dfrac{1}{2} \Rightarrow a = \dfrac{c}{e} = 4$. $b^{2} = a^{2} - c^{2} = 16 - 4 = 12 \Rightarrow b = 2\sqrt{3}$.
\item Équation réduite (centre $O$, axe focal $(Ox)$) : $\dfrac{x^{2}}{16} + \dfrac{y^{2}}{12} = 1$.
\item Foyers : $F(2\,;\,0)$ et $F'(-2\,;\,0)$. Directrices : $x = \pm \dfrac{a}{e} = \pm \dfrac{a^{2}}{c} = \pm \dfrac{16}{2} = \pm 8$, soit $x = 8$ et $x = -8$.
\item Sommets : sur l'axe focal $A(4\,;\,0)$, $A'(-4\,;\,0)$ ; sur l'axe non focal $B(0\,;\,2\sqrt{3})$, $B'(0\,;\,-2\sqrt{3})$.
\end{enumerate}
\end{corrige}

\begin{corrige}[Exercice 6 — Tangentes à une hyperbole]
\begin{enumerate}
\item Hyperbole $(H) : \dfrac{x^{2}}{9} - \dfrac{y^{2}}{4} = 1$. Cherchons les tangentes de pente $1$ : $y = x + m$.
  Injection : $\dfrac{x^{2}}{9} - \dfrac{(x+m)^{2}}{4} = 1 \iff 4x^{2} - 9(x^{2}+2mx+m^{2}) = 36$
  $\iff -5x^{2} - 18mx - 9m^{2} - 36 = 0 \iff 5x^{2} + 18mx + 9m^{2} + 36 = 0$.
  Tangence $\iff$ discriminant nul :
  $\Delta = (18m)^{2} - 4 \times 5 \times (9m^{2} + 36) = 324m^{2} - 180m^{2} - 720 = 144m^{2} - 720$.
  $\Delta = 0 \iff m^{2} = 5 \iff m = \sqrt{5}$ ou $m = -\sqrt{5}$.
  Les deux tangentes sont : $y = x + \sqrt{5}$ et $y = x - \sqrt{5}$.
\item Vérification : $M_{0}\left(5\,;\,\dfrac{8}{3}\right)$. $\dfrac{25}{9} - \dfrac{64/9}{4} = \dfrac{25}{9} - \dfrac{16}{9} = 1$ : $M_{0} \in (H)$.
  Tangente en $M_{0}(x_{0}\,;\,y_{0})$ par dédoublement : $\dfrac{x x_{0}}{9} - \dfrac{y y_{0}}{4} = 1$.
  $\dfrac{5x}{9} - \dfrac{(8/3)y}{4} = 1 \iff \dfrac{5x}{9} - \dfrac{2y}{3} = 1 \iff 5x - 6y = 9$.
\end{enumerate}
\textbf{Point de vigilance :} la formule de dédoublement $\dfrac{xx_{0}}{a^{2}} - \dfrac{yy_{0}}{b^{2}} = 1$ n'est valable que si $M_{0}$ appartient à la conique : la vérification préalable est exigée au Bac.
\end{corrige}

\begin{corrige}[Exercice 7 — Parabole et propriété optique]
\begin{enumerate}
\item Parabole d'équation réduite $y^{2} = 2px$ (axe $(Ox)$, ouverture vers $x>0$). $A(2\,;\,4) \in (P) \Rightarrow 16 = 4p \Rightarrow p = 4$.
  Équation : $y^{2} = 8x$. Foyer $F\left(\dfrac{p}{2}\,;\,0\right) = F(2\,;\,0)$. Directrice $D : x = -\dfrac{p}{2} = -2$.
\item Tangente en $A(x_{0}\,;\,y_{0})$ : $y y_{0} = p(x + x_{0}) \Rightarrow 4y = 4(x + 2) \iff y = x + 2$.
\item $H$ projeté orthogonal de $A(2\,;\,4)$ sur $D\,(x=-2)$ : $H(-2\,;\,4)$.
  Milieu de $[FH]$ : $\left(\dfrac{2+(-2)}{2}\,;\,\dfrac{0+4}{2}\right) = (0\,;\,2)$.
  Ce point vérifie $y = x + 2$ ($2 = 0 + 2$) : il est sur la tangente $(T)$.
  $\overrightarrow{FH} = (-4\,;\,4)$ et un vecteur directeur de $(T)$ est $\vec{u}(1\,;\,1)$.
  $\overrightarrow{FH} \cdot \vec{u} = -4 + 4 = 0 \Rightarrow (T) \perp (FH)$.
  $(T)$ passe par le milieu de $[FH]$ et lui est perpendiculaire : c'est la \cle{médiatrice} de $[FH]$.
  (Autre preuve : tout point $M$ de $(T)$ vérifie $MF = MH$.)
  Cela illustre la \cle{propriété optique de la parabole} : la tangente en un point est bissectrice de l'angle formé par la droite parallèle à l'axe passant par ce point et la droite joignant ce point au foyer. Tout rayon incident parallèle à l'axe se réfléchit en passant par le foyer.
  \textbf{Applications au Cameroun :} antennes paraboliques (réception satellite MTN/Orange/CAMTEL), phares de véhicules, fours solaires.
\end{enumerate}
\end{corrige}

\begin{corrige}[Exercice 8 — Image d'une ellipse par une similitude]
\begin{enumerate}
\item Ellipse $(E) : \dfrac{x^{2}}{16} + \dfrac{y^{2}}{9} = 1$. $a=4$, $b=3$, $c = \sqrt{16-9} = \sqrt{7}$, $e = \dfrac{\sqrt{7}}{4}$.
\item Similitude directe $s$ de centre $O$, rapport $k=2$, angle $\dfrac{\pi}{2}$.
  $M(x\,;\,y) \mapsto M'(x'\,;\,y')$ avec rotation $\frac{\pi}{2}$ : $(x,y) \to (-y,x)$, puis homothétie $2$ : $(-2y,\,2x)$.
  Donc $x' = -2y$, $y' = 2x$. Inverse : $x = \dfrac{y'}{2}$, $y = -\dfrac{x'}{2}$.
  $M \in (E) \iff \dfrac{x^{2}}{16} + \dfrac{y^{2}}{9} = 1 \iff \dfrac{y'^{2}/4}{16} + \dfrac{x'^{2}/4}{9} = 1 \iff \dfrac{x'^{2}}{36} + \dfrac{y'^{2}}{64} = 1$.
\item $(E')$ est une ellipse d'équation réduite $\dfrac{x^{2}}{36} + \dfrac{y^{2}}{64} = 1$ (grand axe vertical).
  $a' = 8$, $b' = 6$. Sommets : $(\pm 6\,;\,0)$ et $(0\,;\,\pm 8)$.
  $c' = \sqrt{a'^{2} - b'^{2}} = \sqrt{64-36} = \sqrt{28} = 2\sqrt{7}$.
  Excentricité $e' = \dfrac{c'}{a'} = \dfrac{2\sqrt{7}}{8} = \dfrac{\sqrt{7}}{4} = e$.
  \textbf{Commentaire :} l'excentricité est conservée par la similitude (invariance de la forme), seules les dimensions sont multipliées par $k=2$ et l'ellipse a pivoté de $\dfrac{\pi}{2}$.
\end{enumerate}
\end{corrige}

\begin{corrige}[Exercice 9 — Type Bac : construction et étude complète d'une conique]
\begin{enumerate}
\item $e = \dfrac{3}{5} \in\,]0\,;\,1[$ : $(\mathcal{C})$ est une \textbf{ellipse} de foyer $F(3\,;\,0)$ et de directrice $(D) : x = \dfrac{25}{3}$.
\item \textbf{Construction point par point :}
  \begin{enumerate}
    \item Tracer l'axe focal $(FK)$ perpendiculaire à $(D)$ passant par $F$ (c'est l'axe $(Ox)$).
    \item Choisir un point $P$ sur cet axe. Tracer la perpendiculaire $(\Delta_{P})$ à l'axe en $P$.
    \item Calculer $\ell = e \cdot d(P, (D)) = \dfrac{3}{5} \times \left|\dfrac{25}{3} - x_{P}\right|$.
    \item Tracer le cercle de centre $F$ et de rayon $\ell$. Il coupe $(\Delta_{P})$ en deux points $M_{1}, M_{2}$ qui vérifient $FM_{i} = e \cdot d(M_{i}, (D))$, donc appartiennent à $(\mathcal{C})$.
    \item Répéter pour trois positions de $P$ de part et d'autre de $F$ pour obtenir six points.
  \end{enumerate}
  Les sommets sur l'axe s'obtiennent en résolvant $MF = e\,d(M,(D))$ sur l'axe : $x = 5$ et $x = -5$.
\item \textbf{Équation cartésienne :} Soit $M(x\,;\,y)$.
  $MF = e\,d(M,(D)) \iff MF^{2} = e^{2}\,d(M,(D))^{2}$ (les deux membres sont positifs, équivalence conservée).
  \[ (x-3)^{2} + y^{2} = \dfrac{9}{25}\left(x - \dfrac{25}{3}\right)^{2}. \]
  Développement : $x^{2} - 6x + 9 + y^{2} = \dfrac{9}{25}x^{2} - 2 \times \dfrac{9}{25} \times \dfrac{25}{3}x + \dfrac{9}{25} \times \dfrac{625}{9} = \dfrac{9}{25}x^{2} - 6x + 25$.
  $x^{2} - \dfrac{9}{25}x^{2} + y^{2} = 25 - 9 \iff \dfrac{16}{25}x^{2} + y^{2} = 16$.
  Division par $16$ : $\dfrac{x^{2}}{25} + \dfrac{y^{2}}{16} = 1$.
\item $a = 5$, $b = 4$, $c = \sqrt{25-16} = 3$.
  \begin{itemize}
    \item Sommets : $A(5\,;\,0)$, $A'(-5\,;\,0)$, $B(0\,;\,4)$, $B'(0\,;\,-4)$.
    \item Second foyer $F'(-3\,;\,0)$ (symétrique de $F$ par rapport à $O$).
    \item Seconde directrice : $x = -\dfrac{a^{2}}{c} = -\dfrac{25}{3}$.
  \end{itemize}
\item Point d'abscisse $3$ : $\dfrac{9}{25} + \dfrac{y^{2}}{16} = 1 \Rightarrow \dfrac{y^{2}}{16} = \dfrac{16}{25} \Rightarrow y = \pm \dfrac{16}{5}$.
  Ordonnée positive : $M_{0}\left(3\,;\,\dfrac{16}{5}\right)$.
  Tangente en $M_{0}$ (par dédoublement, $M_{0}$ vérifié sur la conique) :
  $\dfrac{3x}{25} + \dfrac{(16/5)y}{16} = 1 \iff \dfrac{3x}{25} + \dfrac{y}{5} = 1 \iff 3x + 5y = 25$.
\end{enumerate}
\textbf{Barème indicatif (sur 5 pts) :} 1) $0,5$ ; 2) $1$ ; 3) $1,5$ ; 4) $1$ ; 5) $1$.
\textbf{Points de vigilance :} citer $0 < e < 1$ pour justifier la nature ; élever au carré en précisant que les deux membres sont positifs (équivalence conservée) ; pour la tangente, vérifier que $M_{0}$ est sur la conique avant d'appliquer le dédoublement.
\end{corrige}

\begin{corrige}[Exercice 10 — Type Bac : reconnaissance d'une courbe du second degré]
\begin{enumerate}
\item Mise sous forme canonique par complétion des carrés :
  $9x^{2} + 36x = 9(x^{2} + 4x) = 9[(x+2)^{2} - 4] = 9(x+2)^{2} - 36$.
  $-4y^{2} + 8y = -4(y^{2} - 2y) = -4[(y-1)^{2} - 1] = -4(y-1)^{2} + 4$.
  L'équation $9x^{2} + 36x - 4y^{2} + 8y - 4 = 0$ devient :
  \[ 9(x+2)^{2} - 36 - 4(y-1)^{2} + 4 - 4 = 0 \iff 9(x+2)^{2} - 4(y-1)^{2} = 36. \]
  Division par $36$ : $\dfrac{(x+2)^{2}}{4} - \dfrac{(y-1)^{2}}{9} = 1$.
  Posons $\Omega(-2\,;\,1)$, $X = x + 2$, $Y = y - 1$. Dans le repère $(\Omega\,;\vec{i},\vec{j})$, $(\Gamma)$ a pour équation $\dfrac{X^{2}}{4} - \dfrac{Y^{2}}{9} = 1$ : c'est une \textbf{hyperbole} avec $a = 2$, $b = 3$ (signes opposés pour $X^{2}$ et $Y^{2}$).
\item Dans le repère initial $(O\,;\vec{i},\vec{j})$ :
  \begin{itemize}
    \item Centre : $\Omega(-2\,;\,1)$.
    \item Sommets : $X = \pm 2$, $Y = 0 \Rightarrow A(0\,;\,1)$ et $A'(-4\,;\,1)$.
    \item $c = \sqrt{a^{2}+b^{2}} = \sqrt{4+9} = \sqrt{13}$.
    \item Foyers : $F(-2 + \sqrt{13}\,;\,1)$ et $F'(-2 - \sqrt{13}\,;\,1)$.
  \end{itemize}
\item Asymptotes dans $(\Omega\,;\vec{i},\vec{j})$ : $Y = \pm \dfrac{b}{a}X = \pm \dfrac{3}{2}X$.
  Dans $(O\,;\vec{i},\vec{j})$ : $y - 1 = \dfrac{3}{2}(x+2) \iff y = \dfrac{3}{2}x + 4$ et $y - 1 = -\dfrac{3}{2}(x+2) \iff y = -\dfrac{3}{2}x - 2$.
\item Excentricité : $e = \dfrac{c}{a} = \dfrac{\sqrt{13}}{2}$.
  Directrices dans $(\Omega)$ : $X = \pm \dfrac{a^{2}}{c} = \pm \dfrac{4}{\sqrt{13}} = \pm \dfrac{4\sqrt{13}}{13}$.
  Dans $(O)$ : $x = -2 \pm \dfrac{4\sqrt{13}}{13}$.
\item \textbf{Tracé :} placer le centre $\Omega(-2\,;\,1)$, tracer les axes de symétrie $y = 1$ et $x = -2$, le rectangle fondamental de côtés $2a = 4$ (horizontal) et $2b = 6$ (vertical) centré en $\Omega$, les asymptotes (diagonales de ce rectangle), les sommets $(0\,;\,1)$ et $(-4\,;\,1)$, puis les deux branches de l'hyperbole s'approchant des asymptotes.
\item Sommet d'abscisse la plus grande : $A(0\,;\,1)$ (correspond à $X=2, Y=0$).
  La tangente en un sommet d'hyperbole $\dfrac{X^{2}}{a^{2}} - \dfrac{Y^{2}}{b^{2}} = 1$ est perpendiculaire à l'axe focal : c'est la droite $X = a$, soit $X = 2$.
  Dans $(O)$ : $x + 2 = 2 \iff x = 0$ (l'axe $(Oy)$).
  Vérification par dédoublement dans $(\Omega)$ : $\dfrac{X \cdot 2}{4} - \dfrac{Y \cdot 0}{9} = 1 \iff X = 2 \iff x = 0$.
\end{enumerate}
\textbf{Barème indicatif (sur 6 pts) :} 1) $1,5$ ; 2) $1$ ; 3) $1$ ; 4) $1$ ; 5) $1$ ; 6) $0,5$.
\textbf{Points de vigilance :} factoriser correctement les coefficients de $x^{2}$ et $y^{2}$ avant de compléter le carré ; ne pas oublier la translation inverse pour donner les éléments dans le repère initial ; citer la condition $a \times b < 0$ (signes opposés des coefficients de $x^{2}$ et $y^{2}$) qui garantit une hyperbole.
\end{corrige}

\begin{corrige}[Exercice 11 — Type Bac : coniques autour de nous]
\textbf{Partie A. Ombre d'une sphère}
\begin{enumerate}
\item Les rayons du soleil sont parallèles. L'ensemble des rayons tangents à la sphère forme un \textbf{cylindre de révolution} dont l'axe est parallèle à la direction des rayons (chaque génératrice touche la sphère en un point du grand cercle de contact). L'ombre portée sur le mur est l'intersection de ce cylindre avec le plan du mur : or la section plane d'un cylindre (tout comme celle d'un cône) est une \cle{conique}.
\item \begin{itemize}
    \item Si le mur est \textbf{perpendiculaire} à la direction des rayons, la section est un \textbf{cercle}.
    \item Si le mur est \textbf{incliné} (non parallèle aux rayons), la section est une \textbf{ellipse}.
    \item (Si le mur était parallèle aux rayons, l'ombre dégénérerait en bande rectangulaire infinie.)
  \end{itemize}
\end{enumerate}

\textbf{Partie B. Arche parabolique}
Repère : origine au sommet $S$ de l'arche, $(Ox)$ horizontal, $(Oy)$ vertical \textbf{vers le bas}, unité le mètre.
\begin{enumerate}
\item La parabole a pour équation $x^{2} = 2py$ (axe vertical, ouverture vers le bas du repère, c'est-à-dire vers le sol). Les pieds de l'arche sont à $\pm 30$ m horizontalement et $15$ m en dessous du sommet : le point $(30\,;\,15)$ est sur la courbe.
  $30^{2} = 2p \times 15 \iff 900 = 30p \iff p = 30$.
  Équation : $x^{2} = 60y$.
\item Foyer : $F\left(0\,;\,\dfrac{p}{2}\right) = F(0\,;\,15)$, c'est-à-dire à $15$ m sous le sommet, au niveau du sol, au milieu de la portée.
  Directrice : $y = -\dfrac{p}{2} = -15$, droite horizontale située $15$ m au-dessus du sommet.
\item Le camion roule à $2$ m du pied droit. Pied droit en $x=30$.
  Bord extérieur du camion : $x = 30 - 2 = 28$.
  Bord intérieur du camion : $x = 28 - 3 = 25$ (largeur $3$ m).
  Hauteur sous l'arche à l'abscisse $x$ : $h(x) = 15 - y = 15 - \dfrac{x^{2}}{60}$.
  Le point critique est le bord extérieur ($x = 28$, arche la plus basse) :
  $h(28) = 15 - \dfrac{784}{60} = 15 - \num{13,066\ldots} \approx \num{1,93}$ m $< 4$ m.
  \textbf{Le camion ne peut pas passer} sous l'arche à cet endroit.
  (Même à $x = 25$ : $h(25) = 15 - \dfrac{625}{60} \approx \num{4,58}$ m $> 4$ m, mais le coin extérieur du camion, haut de $4$ m en $x = 28$, heurte l'arche.)
\item $h(x) = 10 \iff 15 - \dfrac{x^{2}}{60} = 10 \iff x^{2} = 300 \iff x = \pm\sqrt{300} = \pm 10\sqrt{3} \approx \pm \num{17,32}$.
  La hauteur sous l'arche vaut $10$ m à environ $\num{17,3}$ m de part et d'autre du sommet.
\end{enumerate}
\textbf{Barème indicatif (sur 5 pts) :} Partie A : $1$ pt ; Partie B : 1) $1$ ; 2) $0,5$ ; 3) $1,5$ ; 4) $1$.
\textbf{Points de vigilance :} bien identifier les coordonnées des pieds de l'arche dans le repère choisi (l'axe $(Oy)$ est dirigé vers le bas, donc $y = 15 > 0$ au sol) ; pour le camion, tester le bord le plus proche du pied (arche la plus basse) et non le centre du camion ; conclure par une phrase affirmative ou négative.
\end{corrige}

\newpage

\coursec{Fiche bilan}

\coursec{Coniques : sections d'un cône de révolution}

\begin{savaistu}[Origine historique et contexte camerounais]
Les coniques ont été étudiées systématiquement par \textbf{Apollonius de Perge} (vers 200 av. J.-C.) dans son traité \emph{Les Coniques}. Il a montré qu'ellipse, parabole et hyperbole proviennent d'un même objet : le cône de révolution. Au Cameroun, on les retrouve partout : les antennes paraboliques \textbf{MTN} et \textbf{Orange} (paraboloïdes de révolution) captent les signaux satellites grâce à la propriété focale ; les orbites des satellites (et de la lune) sont des \textbf{ellipses} (lois de Kepler) ; les arches du pont sur le Wouri ou les ombres portées par l'éclairage public (hyperboles) en sont d'autres illustrations.
\end{savaistu}

\courssub{Définition par section plane}
\begin{definition}[Cône de révolution et sections coniques]
Soit une droite $\Delta$ (l'axe) et une droite $\mathcal{G}$ (génératrice) sécante à $\Delta$ en un point $S$ (sommet) et formant un angle $\alpha \in ]0; \frac{\pi}{2}[$. Le \cle{cône de révolution} de sommet $S$, d'axe $\Delta$ et d'angle au sommet $2\alpha$ est la surface engendrée par la rotation de $\mathcal{G}$ autour de $\Delta$. Il est composé de deux nappes opposées.
Une \cle{conique} (ou section conique) est l'intersection de ce cône par un plan $\mathcal{P}$ \textbf{ne passant pas par le sommet $S$}.
\end{definition}

\begin{propriete}[Classification selon l'inclinaison du plan]
Soit $\beta$ l'angle entre le plan sécant $\mathcal{P}$ et l'axe $\Delta$ ($0 \le \beta \le \frac{\pi}{2}$).
\begin{itemize}
    \item \textbf{Ellipse} : $\alpha < \beta \le \frac{\pi}{2}$. Le plan ne coupe qu'une seule nappe et n'est pas parallèle à une génératrice. Cas particulier : $\beta = \frac{\pi}{2}$ (plan perpendiculaire à l'axe) $\rightarrow$ \cle{cercle}.
    \item \textbf{Parabole} : $\beta = \alpha$. Le plan est parallèle à une génératrice du cône. Il ne coupe qu'une seule nappe.
    \item \textbf{Hyperbole} : $0 \le \beta < \alpha$. Le plan coupe les deux nappes du cône.
\end{itemize}
\end{propriete}

\begin{popfigure}
\begin{tikzpicture}[scale=0.9, every node/.style={font=\small}]
% Cone
\draw[popblue, thick] (0,0) -- (-3,-4) -- (3,-4) -- cycle;
\draw[popblue, thick, dashed] (0,0) -- (-3,4) -- (3,4) -- cycle; % Upper nappe (dashed for visibility)
\draw[popblue, thick] (0,0) -- (3,4);
\draw[popblue, thick] (0,0) -- (-3,4);
% Axis
\draw[popmuted, ->] (0,4.5) -- (0,-4.5) node[below] {$\Delta$};
\node[popmuted] at (0.3, 4.2) {$S$};
% Plan Ellipse (horizontal)
\fill[popgreen!30, opacity=0.5] (-2.5, 1.5) -- (2.5, 1.5) -- (2.5, 2) -- (-2.5, 2) -- cycle;
\draw[popgreen, thick] (-2.5, 1.75) -- (2.5, 1.75) node[right] {$\mathcal{P}_E$ (Ellipse)};
% Plan Parabola (parallel to generatrix)
\fill[poporange!30, opacity=0.5] (-3, -4) -- (3, 4) -- (3.5, 4.5) -- (-2.5, -3.5) -- cycle;
\draw[poporange, thick] (-3, -4) -- (3, 4) node[above right] {$\mathcal{P}_P$ (Parabole)};
% Plan Hyperbola (vertical)
\fill[poppurple!30, opacity=0.5] (1.5, -4) -- (1.5, 4) -- (2, 4) -- (2, -4) -- cycle;
\draw[poppurple, thick] (1.75, -4) -- (1.75, 4) node[above] {$\mathcal{P}_H$ (Hyperbole)};
% Angle alpha
\draw[popblue] (0.3, 0) arc (0:53:0.3) node[right, pos=0.5] {$\alpha$};
\end{tikzpicture}
\legende{Les trois types de sections coniques dans un cône de révolution}
\end{popfigure}

\coursec{Définition monofocale et excentricité}

\courssub{Définition unifiée}
\begin{definition}[Définition monofocale (ou focale-directrice)]
Soient $F$ un point (le \cle{foyer}), $\mathcal{D}$ une droite ne passant pas par $F$ (la \cle{directrice}) et $e$ un réel strictement positif (l'\cle{excentricité}). La conique de foyer $F$, de directrice $\mathcal{D}$ et d'excentricité $e$ est l'ensemble des points $M$ du plan tels que :
\[ \boxed{MF = e \cdot d(M,\mathcal{D})} \qquad \text{ou équivalentement} \qquad \frac{MF}{d(M,\mathcal{D})} = e \]
où $d(M,\mathcal{D})$ désigne la distance du point $M$ à la droite $\mathcal{D}$ (longueur de la perpendiculaire).
\end{definition}

\courssub{Discussion selon l'excentricité $e$}
\begin{propriete}[Nature de la conique]
\begin{itemize}
    \item Si $0 < e < 1$ : la conique est une \cle{ellipse} (bornée, fermée).
    \item Si $e = 1$ : la conique est une \cle{parabole} (non bornée, une branche).
    \item Si $e > 1$ : la conique est une \cle{hyperbole} (non bornée, deux branches disjointes).
\end{itemize}
\end{propriete}

\begin{experience}[Construction point par point d'une conique (Méthode du rapporteur / compas)]
\textbf{Objectif :} Tracer des points d'une conique connaissant $F$, $\mathcal{D}$ et $e$.
\textbf{Matériel :} Règle, compas, rapporteur (ou équerre).
\textbf{Procédé (cas ellipse $e=1/2$) :}
\begin{enumerate}
    \item Tracer la directrice $\mathcal{D}$ (verticale) et le foyer $F$ à gauche de $\mathcal{D}$. Tracer l'axe $\Delta$ (perpendiculaire à $\mathcal{D}$ passant par $F$).
    \item Choisir un point $H$ quelconque sur $\mathcal{D}$. Tracer la perpendiculaire à $\mathcal{D}$ en $H$ (droite $\mathcal{L}$ parallèle à l'axe).
    \item Sur $\mathcal{L}$, le point $M$ vérifie $MH = \frac{1}{e} MF = 2 MF$. Construire $M$ comme intersection de $\mathcal{L}$ avec le cercle de centre $F$ et de rayon $MH/2$ (ou utiliser le théorème de Thalès pour diviser $FH$).
    \item Répéter pour plusieurs positions de $H$ de part et d'autre du pied de $F$ sur $\mathcal{D}$.
    \item Relier les points $M$ trouvés ; utiliser la symétrie par rapport à l'axe $\Delta$ (et au centre pour l'ellipse/hyperbole).
\end{enumerate}
\end{experience}

\courssub{Éléments de symétrie}
\begin{propriete}[Symétries des coniques]
\begin{itemize}
    \item \textbf{Parabole} : Un seul axe de symétrie (la droite $\Delta$ passant par $F$ et $\perp$ à $\mathcal{D}$). Pas de centre de symétrie.
    \item \textbf{Ellipse et Hyperbole} : Deux axes de symétrie orthogonaux (l'axe focal $\Delta$ et sa perpendiculaire par le centre $\Omega$). Elles admettent un \cle{centre de symétrie} $\Omega$ (milieu du segment $FF'$ pour le second foyer $F'$). Ce sont des \cle{coniques centrales}.
\end{itemize}
\end{propriete}

\coursec{Équations réduites et éléments caractéristiques}

\courssub{Choix du repère cartésien}
On place le repère orthonormé $(O;\vec{\imath},\vec{\jmath})$ de façon à simplifier l'équation :
\begin{itemize}
    \item Origine $O$ = centre $\Omega$ (ellipse, hyperbole) ou sommet $S$ (parabole).
    \item Axe des abscisses $Ox$ = axe focal (droite $F\mathcal{D}$).
    \item Axe des ordonnées $Oy$ = axe secondaire (perpendiculaire à $Ox$ en $O$).
\end{itemize}

\courssub{Ellipse : $\dfrac{x^2}{a^2}+\dfrac{y^2}{b^2}=1 \quad (a > b > 0)$}
\begin{propriete}[Éléments caractéristiques de l'ellipse]
\begin{center}
\begin{tabularx}{\linewidth}{|l|X|}
\hline
\rowcolor{popblueL} \textbf{Élément} & \textbf{Expression / Coordonnées} \\
\hline
Équation réduite & $\dfrac{x^2}{a^2}+\dfrac{y^2}{b^2}=1$ \\
\hline
Excentricité & $e = \dfrac{c}{a} \quad (0 < e < 1)$ \\
\hline
Relation fondamentale & $c^2 = a^2 - b^2$ \quad ($c$ : demi-distance focale) \\
\hline
Foyers & $F(-c;0)$, $F'(c;0)$ \\
\hline
Directrices & $\mathcal{D}: x = -\dfrac{a}{e} = -\dfrac{a^2}{c}$, \quad $\mathcal{D}': x = \dfrac{a}{e} = \dfrac{a^2}{c}$ \\
\hline
Sommets (axe focal) & $A(-a;0)$, $A'(a;0)$ \\
\hline
Sommets (axe secondaire) & $B(0;-b)$, $B'(0;b)$ \\
\hline
Axes de symétrie & $Ox$ (axe focal), $Oy$ (axe secondaire) \\
\hline
Centre de symétrie & $O(0;0)$ \\
\hline
Propriété bifocale (définition équivalente) & $MF + MF' = 2a$ \\
\hline
\end{tabularx}
\end{center}
\end{propriete}

\begin{attention}[Pièges classiques ellipse]
\begin{itemize}
    \item Ne pas confondre $a$ (grand axe) et $b$ (petit axe) : \textbf{toujours $a > b$} pour l'équation réduite standard $\frac{x^2}{a^2}+\frac{y^2}{b^2}=1$. Si $b > a$, l'axe focal est $Oy$ et l'équation s'écrit $\frac{x^2}{b^2}+\frac{y^2}{a^2}=1$.
    \item Relation $c^2 = a^2 - b^2$ (soustraction) vs Hyperbole $c^2 = a^2 + b^2$ (addition).
    \item L'excentricité $e = c/a$ est \textbf{toujours $< 1$}. Vérifier ce critère en fin de calcul.
\end{itemize}
\end{attention}

\courssub{Parabole : $y^2 = 2px \quad (p \neq 0)$}
\begin{propriete}[Éléments caractéristiques de la parabole]
\begin{center}
\begin{tabularx}{\linewidth}{|l|X|}
\hline
\rowcolor{poporangeL} \textbf{Élément} & \textbf{Expression / Coordonnées} \\
\hline
Équation réduite & $y^2 = 2px$ \quad ($p>0$ : ouverture vers $+x$ ; $p<0$ : vers $-x$) \\
\hline
Excentricité & $e = 1$ \\
\hline
Paramètre & $p = 2 \cdot OF$ (distance foyer-sommet $\times 2$) \quad $\Leftrightarrow$ \quad $OF = \frac{|p|}{2}$ \\
\hline
Foyer & $F\left(\dfrac{p}{2}; 0\right)$ \\
\hline
Directrice & $\mathcal{D}: x = -\dfrac{p}{2}$ \\
\hline
Sommet & $S(0;0)$ (origine du repère) \\
\hline
Axe de symétrie & $Ox$ (axe focal) \\
\hline
Propriété caractéristique & $MF = d(M,\mathcal{D})$ (équidistance foyer-directrice) \\
\hline
\end{tabularx}
\end{center}
\end{propriete}

\begin{attention}[Sens de l'ouverture]
L'équation $y^2 = 2px$ a l'axe de symétrie horizontal ($Ox$).
\begin{itemize}
    \item $p > 0$ : Parabole ouverte vers la droite (vers $x>0$).
    \item $p < 0$ : Parabole ouverte vers la gauche (vers $x<0$).
\end{itemize}
Si l'axe est vertical, l'équation réduite est $x^2 = 2py$ ($p>0$ vers le haut).
\end{attention}

\courssub{Hyperbole : $\dfrac{x^2}{a^2}-\dfrac{y^2}{b^2}=1 \quad (a>0, b>0)$}
\begin{propriete}[Éléments caractéristiques de l'hyperbole]
\begin{center}
\begin{tabularx}{\linewidth}{|l|X|}
\hline
\rowcolor{poppurpleL} \textbf{Élément} & \textbf{Expression / Coordonnées} \\
\hline
Équation réduite & $\dfrac{x^2}{a^2}-\dfrac{y^2}{b^2}=1$ \\
\hline
Excentricité & $e = \dfrac{c}{a} \quad (e > 1)$ \\
\hline
Relation fondamentale & $c^2 = a^2 + b^2$ \\
\hline
Foyers & $F(-c;0)$, $F'(c;0)$ \\
\hline
Directrices & $\mathcal{D}: x = -\dfrac{a}{e} = -\dfrac{a^2}{c}$, \quad $\mathcal{D}': x = \dfrac{a}{e} = \dfrac{a^2}{c}$ \\
\hline
Sommets (axe transverse) & $A(-a;0)$, $A'(a;0)$ \\
\hline
Points conjugués (axe conjugué) & $B(0;-b)$, $B'(0;b)$ \quad (n'appartiennent pas à la courbe) \\
\hline
Asymptotes & $y = \pm \dfrac{b}{a}x$ (droites passant par le centre $O$) \\
\hline
Axes de symétrie & $Ox$ (axe transverse), $Oy$ (axe conjugué) \\
\hline
Centre de symétrie & $O(0;0)$ \\
\hline
Propriété bifocale & $|MF - MF'| = 2a$ \\
\hline
\end{tabularx}
\end{center}
\end{propriete}

\begin{savaistu}[Asymptotes et hyperbole équilatère]
Les asymptotes guident les branches à l'infini. Si $a=b$, l'hyperbole est \textbf{équilatère} (asymptotes perpendiculaires : $y=\pm x$), son excentricité est $e=\sqrt{2}$. Son équation peut s'écrire $x^2-y^2=a^2$ ou $xy=k$ dans un repère tourné de $45^\circ$.
\end{savaistu}

\coursec{Tangente à une conique en un point}

\courssub{Règle du dédoublement (Méthode opératoire)}
\begin{methode}[Équation de la tangente en $M_0(x_0;y_0)$]
\textbf{Prérequis :} Vérifier que $M_0$ appartient à la conique (substituer dans l'équation).
\textbf{Règle du dédoublement} (valable pour toute conique d'équation $Ax^2+Bxy+Cy^2+Dx+Ey+F=0$) :
\begin{enumerate}
    \item $x^2 \to x_0x$ \quad ; \quad $y^2 \to y_0y$ \quad ; \quad $xy \to \frac{x_0y + xy_0}{2}$
    \item $x \to \frac{x+x_0}{2}$ \quad ; \quad $y \to \frac{y+y_0}{2}$
    \item Les constantes ne changent pas.
\end{enumerate}
\end{methode}

\courssub{Formules réduites à connaître par cœur}
\begin{propriete}[Tangentes aux coniques réduites]
Soit $M_0(x_0;y_0)$ un point de la conique.
\begin{itemize}
    \item \textbf{Ellipse} $\dfrac{x^2}{a^2}+\dfrac{y^2}{b^2}=1$ : \quad $\boxed{\dfrac{x_0x}{a^2}+\dfrac{y_0y}{b^2}=1}$
    \item \textbf{Parabole} $y^2=2px$ : \quad $\boxed{y_0y = p(x+x_0)}$
    \item \textbf{Hyperbole} $\dfrac{x^2}{a^2}-\dfrac{y^2}{b^2}=1$ : \quad $\boxed{\dfrac{x_0x}{a^2}-\dfrac{y_0y}{b^2}=1}$
\end{itemize}
\end{propriete}

\begin{exemplebox}[Tangente à une ellipse]
Déterminer l'équation de la tangente à l'ellipse $\mathcal{E}: \frac{x^2}{25}+\frac{y^2}{9}=1$ au point $M_0(4; \frac{9}{5})$.
\textbf{Vérification :} $\frac{16}{25}+\frac{81/25}{9} = \frac{16}{25}+\frac{9}{25}=1$. $M_0 \in \mathcal{E}$.
\textbf{Application :} $\frac{4x}{25}+\frac{\frac{9}{5}y}{9}=1 \iff \frac{4x}{25}+\frac{y}{5}=1 \iff 4x+5y=25$.
\end{exemplebox}

\begin{exemplebox}[Tangente à une parabole]
Tangente à $\mathcal{P}: y^2=8x$ au point $M_0(2;4)$. $p=4$.
\textbf{Vérification :} $4^2=16=8\times 2$. OK.
\textbf{Formule :} $4y = 4(x+2) \iff y = x+2$.
\end{exemplebox}

\begin{attention}[Erreurs fatales au Bac]
\begin{itemize}
    \item \textbf{Oublier la vérification $M_0 \in \mathcal{C}$}. Si $M_0$ n'est pas sur la courbe, la formule donne la \emph{polaire} (droite des points de contact des tangentes issues de $M_0$), pas la tangente.
    \item Confondre $p$ et $p/2$ pour la parabole (foyer à $p/2$).
    \item Signe moins pour l'hyperbole : $\frac{x_0x}{a^2} \mathbf{-} \frac{y_0y}{b^2}=1$.
\end{itemize}
\end{attention}

\coursec{Reconnaissance et mise sous forme réduite}

\courssub{Reconnaissance immédiate (Équation du type $ax^2+by^2+cx+dy+e=0$)}
\begin{propriete}[Critère du produit $ab$]
Soit l'équation $ax^2+by^2+cx+dy+e=0$ avec $a,b \in \mathbb{R}^*$ (pas de terme $xy$, axes parallèles aux axes du repère).
\begin{itemize}
    \item $ab > 0$ ($a$ et $b$ même signe) $\Rightarrow$ \textbf{Ellipse} (ou ensemble vide, ou point unique si dégénérescence).
    \item $ab < 0$ ($a$ et $b$ signes contraires) $\Rightarrow$ \textbf{Hyperbole} (ou deux droites sécantes si dégénérescence).
    \item $a=0$ xor $b=0$ (un seul terme carré) $\Rightarrow$ \textbf{Parabole} (axe horizontal si $y^2$, vertical si $x^2$).
\end{itemize}
\end{propriete}

\courssub{Méthode de réduction (Complétion des carrés)}
\begin{methode}[Mise sous forme canonique]
Pour une équation $ax^2+by^2+cx+dy+e=0$ ($ab \neq 0$) :
\begin{enumerate}
    \item Regrouper les termes en $x$ et en $y$ : $a(x^2+\frac{c}{a}x) + b(y^2+\frac{d}{b}y) + e = 0$.
    \item Compléter les carrés :
    \[ a\left[\left(x+\frac{c}{2a}\right)^2 - \frac{c^2}{4a^2}\right] + b\left[\left(y+\frac{d}{2b}\right)^2 - \frac{d^2}{4b^2}\right] + e = 0 \]
    \item Mettre les constantes à droite :
    \[ a\left(x+\frac{c}{2a}\right)^2 + b\left(y+\frac{d}{2b}\right)^2 = \frac{c^2}{4a} + \frac{d^2}{4b} - e = K \]
    \item \textbf{Analyser $K$ et les signes de $a,b$} :
    \begin{itemize}
        \item Si $ab>0$ (ellipse) : $K$ doit avoir le même signe que $a$ (et $b$). Diviser par $K$ pour obtenir $1$ à droite. Centre $\Omega(-\frac{c}{2a}; -\frac{d}{2b})$.
        \item Si $ab<0$ (hyperbole) : $K \neq 0$. Diviser par $K$. Centre $\Omega(-\frac{c}{2a}; -\frac{d}{2b})$.
        \item Si $a=0$ ou $b=0$ (parabole) : Isoler le terme du premier degré pour faire apparaître $y^2=2px$ ou $x^2=2py$. Sommet $\Omega$.
    \end{itemize}
\end{enumerate}
\end{methode}

\begin{exoresolu}[Réduction d'une ellipse]
Réduire $4x^2+9y^2-16x+18y-11=0$.
\textbf{Solution :}
$ab = 36 > 0 \Rightarrow$ Ellipse.
$4(x^2-4x) + 9(y^2+2y) = 11$
$4[(x-2)^2-4] + 9[(y+1)^2-1] = 11$
$4(x-2)^2 - 16 + 9(y+1)^2 - 9 = 11$
$4(x-2)^2 + 9(y+1)^2 = 36$
$\dfrac{(x-2)^2}{9} + \dfrac{(y+1)^2}{4} = 1$
\textbf{Repère centré en $\Omega(2;-1)$ :} $X=x-2$, $Y=y+1$. $\frac{X^2}{3^2}+\frac{Y^2}{2^2}=1$.
$a=3, b=2, c=\sqrt{5}, e=\frac{\sqrt{5}}{3}$.
Foyers dans $(O;\vec{\imath},\vec{\jmath})$ : $F(2-\sqrt{5};-1)$, $F'(2+\sqrt{5};-1)$.
Directrices : $x = 2 \pm \frac{9}{\sqrt{5}}$.
\tcblower
\textbf{Structure de la réponse attendue :} 1. Nature. 2. Calculs de complétion. 3. Équation réduite + repère. 4. $a,b,c,e$. 5. Coordonnées des éléments dans le repère initial.
\end{exoresolu}

\begin{exoresolu}[Réduction d'une hyperbole]
Réduire $x^2-4y^2+2x+16y-19=0$.
\textbf{Solution :}
$ab = -4 < 0 \Rightarrow$ Hyperbole.
$(x^2+2x) - 4(y^2-4y) = 19$
$(x+1)^2 - 1 - 4[(y-2)^2 - 4] = 19$
$(x+1)^2 - 4(y-2)^2 = 4$
$\dfrac{(x+1)^2}{4} - \dfrac{(y-2)^2}{1} = 1$
\textbf{Repère centré en $\Omega(-1;2)$ :} $X=x+1$, $Y=y-2$. $\frac{X^2}{2^2}-\frac{Y^2}{1^2}=1$.
$a=2, b=1, c=\sqrt{5}, e=\frac{\sqrt{5}}{2}>1$.
Asymptotes dans $(O)$ : $Y = \pm \frac{1}{2}X \iff y-2 = \pm \frac{1}{2}(x+1)$.
Foyers : $F(-1-\sqrt{5};2)$, $F'(-1+\sqrt{5};2)$.
Directrices : $x = -1 \pm \frac{4}{\sqrt{5}}$.
\end{exoresolu}

\begin{exoresolu}[Réduction d'une parabole]
Réduire $y^2-4y-4x+8=0$.
\textbf{Solution :}
Terme en $x$ au premier degré, terme en $y^2$ $\Rightarrow$ Parabole (axe horizontal).
$y^2-4y = 4x-8$
$(y-2)^2 - 4 = 4x-8$
$(y-2)^2 = 4x-4 = 4(x-1)$
\textbf{Repère centré au sommet $\Omega(1;2)$ :} $X=x-1$, $Y=y-2$. $Y^2 = 4X$.
$2p=4 \Rightarrow p=2 > 0$ (ouverture vers $x$ croissants).
Foyer : $F(1+\frac{p}{2}; 2) = F(2;2)$.
Directrice : $x = 1-\frac{p}{2} = 0$ (l'axe $Oy$).
\end{exoresolu}

\coursec{Coniques, similitudes et applications concrètes}

\courssub{Image par une similitude directe}
\begin{propriete}[Invariants par similitude directe]
Soit $\mathcal{C}$ une conique de foyer $F$, directrice $\mathcal{D}$, excentricité $e$. Soit $s$ une similitude directe de rapport $k>0$ et d'angle $\theta$.
L'image $\mathcal{C}' = s(\mathcal{C})$ est une conique de \textbf{même nature} et de \textbf{même excentricité $e$}.
\begin{itemize}
    \item Ses foyers sont les images des foyers de $\mathcal{C}$.
    \item Ses directrices sont les images des directrices de $\mathcal{C}$.
    \item Ses axes sont les images des axes de $\mathcal{C}$ (rotés de $\theta$).
    \item Les longueurs caractéristiques ($a, b, c, p$) sont multipliées par $k$.
\end{itemize}
\textbf{Justification :} Une similitude conserve les rapports de distances. Pour $M' = s(M)$, $M'F' = k \cdot MF$ et $d(M',\mathcal{D}') = k \cdot d(M,\mathcal{D})$. Donc $\frac{M'F'}{d(M',\mathcal{D}')} = \frac{MF}{d(M,\mathcal{D})} = e$.
\end{propriete}

\courssub{Situations concrètes et modélisation}
\begin{savaistu}[Coniques dans l'environnement camerounais]
\begin{itemize}
    \item \textbf{Antennes paraboliques (MTN, Orange, CAMTEL) :} Surface parabolique (paraboloïde de révolution). Les ondes arrivant parallèlement à l'axe sont réfléchies vers le foyer où se trouve le récepteur (LNB). \emph{Propriété optique : tout rayon parallèle à l'axe est réfléchi vers le foyer.}
    \item \textbf{Orbites satellites / Planètes :} Ellipses (lois de Kepler). La Terre est à un foyer. Périgée/Apogée = sommets. Excentricité $e$ mesure l'aplatissement.
    \item \textbf{Ponts et arches (Pont sur le Wouri, passages supérieurs Yaoundé) :} Souvent modélisés par des arcs de parabole (charge uniforme sur tablier horizontal) ou d'ellipse (esthétique, poussées).
    \item \textbf{Ombres portées :} L'ombre d'un poteau électrique (cylindre) ou d'un abat-jour conique sur un mur/plat au soleil couchant trace une \textbf{hyperbole} (intersection d'un cône de lumière par un plan vertical).
    \item \textbf{Acoustique / Whispering gallery :} Propriété bifocale de l'ellipse : un son émis en $F$ converge vers $F'$ après réflexion sur la paroi.
\end{itemize}
\end{savaistu}

\begin{exemplebox}[Modélisation : Antenne parabolique]
Une antenne parabolique a un diamètre de \SI{1,2}{m} et une profondeur de \SI{0,3}{m} au centre. Déterminer la position du foyer pour placer le LNB.
\textbf{Modèle :} Parabole $y^2=2px$ (sommet au fond de l'antenne, axe horizontal). Le bord correspond à $x=0,3$ (profondeur) et $y=\pm 0,6$ (rayon).
$0,6^2 = 2p \times 0,3 \Rightarrow 0,36 = 0,6p \Rightarrow p = 0,6$ m.
Foyer à $OF = p/2 = \SI{0,3}{m}$ du sommet (au fond). Le LNB se place à \SI{30}{cm} du fond de l'antenne sur l'axe.
\end{exemplebox}

\coursec{Exercices d'entraînement (Niveau Bac)}

\begin{exercice}[Exercice 1 : Identification et éléments]
On considère l'équation $\mathcal{E} : 9x^2+25y^2-36x-50y-164=0$.
\begin{enumerate}
    \item Montrer que $\mathcal{E}$ représente une ellipse. Déterminer son centre $\Omega$ et son équation réduite dans un repère centré en $\Omega$.
    \item Déterminer $a, b, c, e$. Donner les coordonnées des foyers et les équations des directrices dans le repère initial.
    \item Écrire l'équation de la tangente à $\mathcal{E}$ au point $M_0(4; 2)$.
\end{enumerate}
\end{exercice}

\begin{exercice}[Exercice 2 : Hyperbole et asymptotes]
L'hyperbole $\mathcal{H}$ a pour équation $4x^2-y^2-8x-4y-4=0$.
\begin{enumerate}
    \item Mettre sous forme réduite. Préciser le centre $\Omega$, $a, b, c, e$.
    \item Déterminer les équations des asymptotes dans le repère initial.
    \item Tracer $\mathcal{H}$ et ses asymptotes dans un repère centré en $\Omega$ (échelle 1 cm pour 1 unité).
\end{enumerate}
\end{exercice}

\begin{exercice}[Exercice 3 : Parabole et tangente]
Soit la parabole $\mathcal{P}: y^2-4y-8x+12=0$.
\begin{enumerate}
    \item Déterminer son sommet $S$, son foyer $F$, sa directrice $\mathcal{D}$ et son paramètre $p$.
    \item Déterminer l'équation de la tangente $T$ à $\mathcal{P}$ au point d'abscisse $x=3$ et d'ordonnée positive.
    \item Calculer l'angle (en degrés) que fait cette tangente avec l'axe focal.
\end{enumerate}
\end{exercice}

\begin{exercice}[Exercice 4 : Problème ouvert - Situation concrète]
Un projecteur de stade (modélisé par un point $S$) éclaire un terrain horizontal. Le faisceau lumineux est un cône de révolution d'axe vertical et d'angle au sommet $60^\circ$. Le sol est un plan horizontal.
\begin{enumerate}
    \item Quelle est la nature de la limite de la zone éclairée au sol ? Justifier.
    \item Si on incline le projecteur de $30^\circ$ par rapport à la verticale, quelle devient la nature de la limite ? Justifier par l'angle $\beta$.
\end{enumerate}
\end{exercice}

\coursec{Corrigés détaillés}

\begin{corrige}[Exercice 1]
\textbf{1.} $a=9, b=25 \Rightarrow ab=225>0$ : Ellipse.
$9(x^2-4x) + 25(y^2-2y) = 164$
$9[(x-2)^2-4] + 25[(y-1)^2-1] = 164$
$9(x-2)^2 - 36 + 25(y-1)^2 - 25 = 164$
$9(x-2)^2 + 25(y-1)^2 = 225$
$\dfrac{(x-2)^2}{25} + \dfrac{(y-1)^2}{9} = 1$
Repère $(\Omega; \vec{\imath}, \vec{\jmath})$ avec $\Omega(2;1)$. $X=x-2, Y=y-1$.
$\frac{X^2}{5^2}+\frac{Y^2}{3^2}=1$. $a=5, b=3$. Axe focal horizontal ($Ox$).
$c^2 = a^2-b^2 = 25-9=16 \Rightarrow c=4$.
$e = c/a = 4/5 = 0,8 < 1$ (cohérence).
Foyers dans $(\Omega)$ : $F'(-4;0), F(4;0)$. Dans $(O)$ : $F'(-2;1), F(6;1)$.
Directrices dans $(\Omega)$ : $X = \pm a/e = \pm 5/(4/5) = \pm 25/4 = \pm 6,25$.
Dans $(O)$ : $x = 2 \pm 6,25 \Rightarrow \mathcal{D}': x = -4,25$ ; $\mathcal{D}: x = 8,25$.

\textbf{2.} Vérif $M_0(4;2)$ : $9(16)+25(4)-36(4)-50(2)-164 = 144+100-144-100-164 = -164 \neq 0$.
\textbf{ATTENTION :} $M_0(4;2)$ n'est PAS sur l'ellipse ! (Erreur volontaire dans l'énoncé pour tester la vigilance).
Point sur l'ellipse : $x=4 \Rightarrow 9(4) + 25(y-1)^2 = 225 \Rightarrow 25(y-1)^2 = 189 \Rightarrow (y-1)^2 = 7,56 \Rightarrow y = 1 \pm \sqrt{7,56}$.
Si l'énoncé impose $M_0(4;2)$, on ne peut pas écrire la tangente. On écrit la polaire : $\frac{(4-2)(x-2)}{25} + \frac{(2-1)(y-1)}{9} = 1 \Rightarrow \frac{2(x-2)}{25} + \frac{y-1}{9} = 1$.
\end{corrige}

\begin{corrige}[Exercice 2]
\textbf{1.} $4(x^2-2x) - (y^2+4y) = 4$
$4[(x-1)^2-1] - [(y+2)^2-4] = 4$
$4(x-1)^2 - 4 - (y+2)^2 + 4 = 4$
$4(x-1)^2 - (y+2)^2 = 4$
$\dfrac{(x-1)^2}{1} - \dfrac{(y+2)^2}{4} = 1$
Centre $\Omega(1;-2)$. $X=x-1, Y=y+2$. $\frac{X^2}{1^2} - \frac{Y^2}{2^2} = 1$.
$a=1, b=2$. Axe transverse horizontal.
$c^2 = a^2+b^2 = 5 \Rightarrow c=\sqrt{5}$.
$e = c/a = \sqrt{5} > 1$ (cohérence).

\textbf{2.} Asymptotes dans $(\Omega)$ : $Y = \pm \frac{b}{a}X = \pm 2X$.
Dans $(O)$ : $y+2 = \pm 2(x-1)$.
$y+2 = 2x-2 \Rightarrow y = 2x-4$.
$y+2 = -2x+2 \Rightarrow y = -2x$.

\textbf{3.} Tracé : Centre $\Omega(1;-2)$. Rectangle $-1\le X\le 1, -2\le Y\le 2$. Diagonales = asymptotes. Sommets $A(0;-2), A'(2;-2)$. Branches passant par $A, A'$ à l'intérieur des angles verticaux des asymptotes.
\end{corrige}

\begin{corrige}[Exercice 3]
\textbf{1.} $y^2-4y = 8x-12 \Rightarrow (y-2)^2 - 4 = 8x-12 \Rightarrow (y-2)^2 = 8x-8 = 8(x-1)$.
Sommet $S(1;2)$. $2p=8 \Rightarrow p=4$. $p>0$ : ouverture vers $x$ croissants.
Foyer $F(1+\frac{p}{2}; 2) = F(3;2)$.
Directrice $\mathcal{D}: x = 1-\frac{p}{2} = -1$.

\textbf{2.} Point d'abscisse $x=3$ : $(y-2)^2 = 8(3-1) = 16 \Rightarrow y-2 = \pm 4 \Rightarrow y = 6$ ou $y=-2$.
Ordonnée positive $\Rightarrow M_0(3;6)$. Vérification : $(6-2)^2=16=8(2)$. OK.
Tangente (dédoublement sur $(Y-2)^2=8(X-1)$ avec $X=x-1, Y=y-2$) :
Point $M_0$ dans $(\Omega)$ : $X_0=2, Y_0=4$.
Tangente : $Y_0 Y = 4(X+X_0) \Rightarrow 4Y = 4(X+2) \Rightarrow Y = X+2$.
Retour $(O)$ : $y-2 = (x-1)+2 \Rightarrow y = x+3$.
Vérification par dédoublement direct sur $(y-2)^2 = 8(x-1)$ : $(y-2)(y_0-2) = 4\big((x-1) + (x_0-1)\big) \Rightarrow 4(y-2) = 4(x+1) \Rightarrow \boxed{y = x+3}$.

\textbf{3.} Pente de la tangente $m=1$. Angle avec l'axe horizontal (axe focal) : $\arctan(1) = 45^\circ$.
\end{corrige}

\begin{corrige}[Exercice 4]
\textbf{1.} Projecteur vertical $\Rightarrow$ Axe du cône $\perp$ au sol (plan horizontal). $\beta = 90^\circ$. Angle au sommet $2\alpha = 60^\circ \Rightarrow \alpha = 30^\circ$.
$\beta = 90^\circ > \alpha = 30^\circ$. Le plan ne coupe qu'une nappe et n'est pas parallèle à une génératrice.
\textbf{C'est une ellipse.} (En fait un cercle car axe du cône $\perp$ plan $\Rightarrow$ section circulaire).

\textbf{2.} Projecteur incliné de $30^\circ$ par rapport à la verticale $\Rightarrow$ Axe du cône fait $30^\circ$ avec la normale au sol.
Angle entre plan (sol) et axe du cône : $\beta = 90^\circ - 30^\circ = 60^\circ$.
Toujours $\beta = 60^\circ > \alpha = 30^\circ$.
\textbf{C'est encore une ellipse} (non circulaire cette fois, car $\beta \neq 90^\circ$).
\textit{Note :} Pour obtenir une parabole, il faudrait $\beta = \alpha = 30^\circ$ (projecteur incliné de $60^\circ$ par rapport à la verticale). Pour une hyperbole, $\beta < 30^\circ$ (projecteur presque horizontal).
\end{corrige}

\begin{aretenir}[Checklist finale - Prêt pour le Bac]
\begin{itemize}
    \item[$\square$] Je définis une conique comme section d'un cône (position du plan vs génératrice).
    \item[$\square$] J'écris la définition monofocale $MF = e \cdot d(M,\mathcal{D})$ et je lie $e$ au type (ellipse $<1$, parabole $=1$, hyperbole $>1$).
    \item[$\square$] Je connais par cœur les 3 équations réduites, les relations $c^2=a^2\mp b^2$, les coordonnées de $F, F', \mathcal{D}, \mathcal{D}',$ sommets, asymptotes.
    \item[$\square$] Je calcule $e=c/a$ (ellipse/hyperbole) et je vérifie $e<1$ ou $e>1$.
    \item[$\square$] Je construis une conique point par point (règle du compas / Thalès).
    \item[$\square$] J'écris la tangente par dédoublement \textbf{après vérification $M_0 \in \mathcal{C}$}.
    \item[$\square$] Je reconnais la nature par le signe de $ab$ dans $ax^2+by^2+cx+dy+e=0$.
    \item[$\square$] Je réduis une équation (complétion carrés + translation repère) sans erreur de signe.
    \item[$\square$] J'explique l'invariance de la nature et de $e$ par similitude directe.
    \item[$\square$] Je modélise une situation réelle (antenne, orbite, ombre, pont) par une conique.
    \item[$\square$] Ma rédaction est structurée : \textbf{Nature $\to$ Réduction $\to$ Paramètres $\to$ Éléments $\to$ Tangente/Problème}.
\end{itemize}
\end{aretenir}

\findecours

\end{document}
